\chapter{Galois representations of elliptic curves}\label{ch:galrep}

The paper works throughout with the $\ell$-adic representation $\rho=\rho_{E,\ell}$ of $G_\Q$ on the Tate module $T=T_\ell E$ and
with its reduction $\bar\rho$ on $E[\ell]$. Every statement it makes about these representations, as opposed to the algebra of
representations (\Cref{ch:reps}), the field theory (\Cref{ch:galois,ch:ant,ch:local}) or the geometry of the curve
(\Cref{ch:ellcurves,ch:reduction}), is collected and proved here: the local representation at every prime $p\ne\ell$ by reduction
type, the mod $3$ picture of inertia behind Lemma lem:IV and Remark rem:types, the representation at $\ell$, the Serre conductor
of $\bar\rho$, the possible images of $\bar\rho$, and the three statements on which the hypothesis of Theorem thm:main rests:
Proposition prop:ns, Lemma lem:h0 and Proposition prop:Sf. The chapter supplies the elliptic-curve inputs; the representation
theory it uses is proved in \Cref{ch:reps}, the group theory of $\GL_2(\F_3)$ in \Cref{ch:groups}, the local field theory in
\Cref{ch:local}, and the reduction theory (N\'eron models, Kodaira types, Tate curves, Kraus's theorem) in \Cref{ch:reduction}.

\section*{Core results at a glance}

Throughout, $E/\Q$ has conductor $N$ and a global minimal model, $\ell$ is a prime (odd unless said otherwise), $T=T_\ell E$,
$V=T\otimes\Ql$, $\rho=\rho_{E,\ell}$, $\bar\rho=\bar\rho_{E,\ell}$, and $q\ne\ell$ is a prime; $\Phig_q=\Phi_q(\Fbar_q)$ is the geometric
component group at $q$.

\begin{enumerate}[label=\textbf{\thechapter.\arabic*}]
\item \textbf{Conventions.} $\Frob_p$ is an arithmetic Frobenius and $\Frob_p^{-1}$ a geometric one; $\chi_\ell$ is the
$\ell$-adic cyclotomic character, $\chi_d$ the quadratic character of $\Q(\sqrt d)$, $\chi_{q^*}$ the one of conductor $q$, and
$\bar\chi_3=\chi_{-3}$. \emph{(convention)} \Cref{conv:galrep:notation}.
\item \textbf{The representations.} $\rho\colon G_\Q\to\Aut_{\Zl}(T)\simeq\GL_2(\Zl)$, $\rho\bmod\ell^k$ on $E[\ell^k]=T/\ell^kT$, $\bar\rho$ on
$E[\ell]$, $A_f=\ad V$, $\delta_p=\dim V^{I_p}$. \emph{(definition)} \Cref{def:galrep:rho}.
\item \textbf{Basic properties.} $\rho$ is continuous and $\rho\bmod\ell^k$ cuts out $\Q(E[\ell^k])$; $\det\rho=\chi_\ell$, so $\det\bar\rho=\bar\chi_\ell$
is onto $\Fl^\times$ and $\bar\rho(G_{\Q(\sqrt{\ell^*})})$ is the set of elements of $\bar\rho(G_\Q)$ with square determinant; for $\ell=3$,
$\det\bar\rho_{E,3}=\chi_{-3}$ cuts out $\Q(\sqrt{-3})$. For $\ell$ odd and $c$ a complex conjugation, $\rho(c)$ is conjugate in $\GL_2(\Zl)$ to
$\operatorname{diag}(1,-1)$. $\rho$ is unramified outside $N\ell$. \emph{(proved here, citing \Cref{ch:galois,ch:ellcurves})}
\Cref{prop:galrep:rho-basic}.
\item \textbf{Faltings.} $V_\ell E$ is a semisimple $\Ql[G_\Q]$-module, and $\Hom(E,E')\otimes\Zl\simeq\Hom_{G_\Q}(T_\ell E,T_\ell E')$.
\emph{(stated; reference)} \Cref{thm:galrep:faltings}.
\item \textbf{Twists and isogenies.} $\rho_{E^{(d)},\ell}\simeq\rho_{E,\ell}\otimes\chi_d$, so $\ad\rho$, the projective image of $\bar\rho$ and the
irreducibility of $E[\ell]$ are twist invariant; twisting by $\chi_{-3}$ exchanges the mod $3$ images \texttt{3B.1.1}, \texttt{3B.1.2} and
fixes the others; twisting by $\chi_d$, $d\not\equiv1,-3$ modulo squares, sends \texttt{3B.1.1}, \texttt{3B.1.2} to \texttt{3B} and
\texttt{3Cs.1.1} to \texttt{3Cs}, and fixes \texttt{3Ns}, \texttt{3Nn} and the surjective image. If $E[\ell]$ is irreducible, every curve
isogenous to $E$ has an isomorphic $T_\ell$. \emph{(proved here)} \Cref{prop:galrep:twist-isogeny}.
\item \textbf{Irreducibility.} For $\ell$ odd: $\bar\rho$ is irreducible iff absolutely irreducible; it is reducible iff $E$ has a rational
$\ell$-isogeny (for $\ell=3$ iff $\psi_3$ has a rational root); $\bar\rho^{\mathrm{ss}}$ is determined by $a_p\bmod\ell$ on any set of primes
of density one, and so is $\bar\rho$ when it is irreducible. \emph{(proved here, citing \Cref{ch:reps,ch:galois})}
\Cref{prop:galrep:irreducibility}.
\item \textbf{N\'eron, Ogg, Shafarevich.} For $K/\Q_p$ finite and $\ell\ne p$: good reduction $\iff$ $I_K$ trivial on $T_\ell E$ $\iff$ $I_K$
trivial on $E[m]$ for all (equivalently infinitely many) $m$ prime to $p$ \emph{(proved here)}; potentially good reduction $\iff$
$\rho(I_K)$ finite; potentially multiplicative $\iff$ $\rho(I_K)$ infinite \emph{(proof sketched here)}. So $\rho_{E,\ell}$ is unramified at
$p\ne\ell$ iff $p\nmid N$. \Cref{thm:galrep:nos}.
\item \textbf{Inertia invariants and component groups.} For additive reduction at $q$, $E[\ell^k]^{I_q}\simeq\Phig_q[\ell^k]$; for multiplicative
reduction $0\to\mu_{\ell^k}\otimes\varepsilon\to E[\ell^k]^{I_q}\to\Phig_q[\ell^k]\to0$ with $\Phig_q\simeq\Z/n_q$. \emph{(stated; proved in
\Cref{ch:reduction})} \Cref{lem:galrep:component-invariants}.
\item \textbf{Serre-Tate.} For potentially good reduction at $q$: $\rho_{E,\ell}(I_q)\simeq\Gal(L/\Q_q^{\ur})$, $L$ the smallest extension of $\Q_q^{\ur}$
with good reduction, independently of $\ell\ne q$, with characteristic polynomials in $\Z[X]$ \emph{(stated; reference)}; it lies in
$\SL_2(\Zl)$, injects into $\GL_2(\Fl)$ for $\ell$ odd, and $V^{I_q}=0$ unless the reduction is good \emph{(proved here)}; it is cyclic of
order $12/\gcd(v_q(\Delta_{\min}),12)$ for $q\ge5$ \emph{(proved here from Kraus's theorem)}. \Cref{thm:galrep:serre-tate}.
\item \textbf{Good primes.} For $p\nmid N\ell$, $\rho$ is unramified at $p$, $\Frob_p$ has characteristic polynomial $X^2-a_pX+p$ on $V$, both
Frobenii have eigenvalues $\alpha/\beta,1,\beta/\alpha$ on $\ad V$, $\delta_p=2$, and $\tr\bar\rho(\Frob_p)=a_p\bmod\ell$. \emph{(proved here, citing
\Cref{ch:ellcurves})} \Cref{prop:galrep:good}.
\item \textbf{The Tate shape.} At a multiplicative prime $q$, in a $\Zl$-basis $e_1,e_2$ of $T$,
$\rho(g)=\sm{\chi_\ell\varepsilon(g)}{\varepsilon\kappa(g)}{0}{\varepsilon(g)}$ on $G_{\Q_q}$, with $\varepsilon$ unramified, $\varepsilon(\Frob_q)=a_q(E)$, and $\kappa$
the Kummer cocycle of the Tate parameter $q_E$. On $I_q$, $\rho(\sigma)=\sm{1}{n_qt_\ell(\sigma)}{0}{1}$; $V^{I_q}=\Ql e_1$ ($\delta_q=1$), the unique
stable line; $(\ad V)^{I_q}=\Ql e$, $e=e_1\otimes e_2^\vee$, with $\Frob_q$ acting by $q$. \emph{(proved here, given \Cref{ch:reduction})}
\Cref{thm:galrep:tate-shape}.
\item \textbf{Lemma lem:tamagawa-exp.} At a multiplicative prime $q$, $\rho\bmod\ell^t$ is unramified at $q$ iff $\ell^t\mid n_q$, split or not;
$\#E[\ell^t]^{I_q}=\ell^{t+\min(t,v_\ell(n_q))}$; if $\ell\nmid n_q$, $\bar\rho|_{D_q}$ is a non-split extension of $\varepsilon$ by $\varepsilon\bar\chi_\ell$.
\emph{(proved here, two proofs)} \Cref{lem:galrep:tamagawa-exp}.
\item \textbf{Potentially multiplicative additive primes.} If $v_q(j)<0$ and the reduction is additive ($q=2$ allowed), $\rho|_{G_{\Q_q}}\simeq\rho'\otimes\chi$
with $\chi$ ramified quadratic and $\rho'$ of Tate shape with $v_q(q_E)=-v_q(j)$; $V^{I_q}=0$, $\ad V\simeq\ad V'$, $\rho\bmod\ell^t$ is ramified at $q$
for every $t\ge1$, and $f_q=2$ ($q$ odd), $4$ or $6$ ($q=2$). \emph{(proved here, given \Cref{ch:reduction})} \Cref{prop:galrep:potmult}.
\item \textbf{$\delta_p$.} $\delta_p=2,1,0$ for good, multiplicative, additive reduction at $p\ne\ell$, and $f_p=2-\delta_p+\operatorname{sw}_p(V)$.
\emph{(proved here)} \Cref{cor:galrep:delta}.
\item \textbf{Additive potentially good reduction, $q\ge5$.} $\rho(I_q)=\langle g\rangle$ with $g=-1$ ($\kod{I}_0^*$), $\sigma$ ($\kod{IV}$, $\kod{IV}^*$),
$\tau$ ($\kod{III}$, $\kod{III}^*$), $-\sigma$ ($\kod{II}$, $\kod{II}^*$), where $\sigma^2+\sigma+1=0$ and $\tau^2=-1$; $(\ad V)^{I_q}$ is $\ad V$, $\Ql(2\sigma+1)$,
$\Ql\tau$, $\Ql(2\sigma+1)$, and $\Frob_q$ (arithmetic or geometric) acts on the line by $\chi_{-3}(q)$, $\chi_{-4}(q)$, $\chi_{-3}(q)$; $V^{I_q}=0$ and
$f_q=2$. \emph{(proved here, given Kraus's theorem)} \Cref{thm:galrep:additive-shapes}.
\item \textbf{Lemma lem:IV.} At a prime $q\ne3$ of type $\kod{IV}$ or $\kod{IV}^*$ the reduction is potentially good, $\rho_{E,\ell}(I_q)$ is cyclic
of order $3$ for every $\ell\ne q$, its image on $E[3]$ is a nontrivial unipotent group, $E[3]^{I_q}$ is a line, $\bar\rho_{E,3}$ is tame at $q$
with Serre exponent $1$, $f_q=2$, $T_3E\simeq\Z_3[\zeta_3]$, and $\Frob_q$ acts on $(\ad V)^{I_q}=\Ql(2\sigma+1)$ by $\chi_{-3}(q)$.
\emph{(proved here, two proofs)} \Cref{lem:galrep:IV}.
\item \textbf{Types $\kod{IV}$, $\kod{IV}^*$ at $2$.} $f_2=2$, $v_2(\Delta_{\min})=4$, resp.\ $8$, the Serre exponent of $\bar\rho_{E,3}$ at $2$ is $1$,
$L_2(A_f,s)^{-1}=1+2^{-s}$, $P_2(X)=1+2X$, $e_3(E,2)=1$ and $e_\ell(E,2)=0$ for $\ell\ge5$. \emph{(proved here)} \Cref{cor:galrep:IV-at-2}.
\item \textbf{Inertia mod $3$ at $q\ge5$ (Remark rem:types).} $E[3]^{I_q}=0$ and the Serre exponent of $\bar\rho_{E,3}$ at $q$ is $2=f_q$ for
every additive type except $\kod{IV}$, $\kod{IV}^*$ (exponent $1$); the twist by $\chi_{q^*}$ lowers the Serre exponent exactly for
$\kod{II}$, $\kod{II}^*$, $\kod{I}_n^*$, $\kod{I}_0^*$ and the conductor exponent of $E$ exactly for $\kod{I}_n^*$, $\kod{I}_0^*$. \emph{(proved here)}
\Cref{prop:galrep:mod3-inertia}.
\item \textbf{Good reduction at $\ell$.} For $\ell\nmid N$: $E[\ell^k]$ is the generic fibre of a finite flat group scheme over $\Zl$
\emph{(proof sketched here)}; $V$ is crystalline at $\ell$ with Hodge-Tate weights $0$, $1$, and the Serre weight of $\bar\rho$ is $2$ \emph{(stated;
reference)}. \Cref{thm:galrep:good-at-l}.
\item \textbf{Fundamental characters.} $\omega\colon I_\ell\to\Fl^\times$ (equal to $\bar\chi_\ell|_{I_\ell}$) and $\omega_2\colon I_\ell\to\F_{\ell^2}^\times$ of order
$\ell^2-1$ with $\omega_2^{\ell+1}=\omega$; conjugation by $\Frob_\ell$ sends $\omega_2$ to $\omega_2^\ell$, so $\omega_2$ does not extend to $G_{\Q_\ell}$.
\emph{(definition; proved here, given \Cref{ch:local})} \Cref{def:galrep:fundamental-characters},
\Cref{prop:galrep:fundamental-characters}.
\item \textbf{Deligne and Fontaine.} For $\ell$ odd, $\ell\nmid N$: if $\ell\nmid a_\ell$, $\bar\rho|_{D_\ell}\simeq\sm{\bar\chi_\ell\mu(a_\ell^{-1})}{*}{0}{\mu(a_\ell)}$; if
$\ell\mid a_\ell$, $\bar\rho|_{I_\ell}\otimes\F_{\ell^2}\simeq\omega_2\oplus\omega_2^\ell$ and $\bar\rho|_{D_\ell}$ is irreducible. Hence $a_\ell\bmod\ell$ is determined by
$\bar\rho|_{D_\ell}$. \emph{(proof sketched here)} \Cref{thm:galrep:deligne-fontaine}.
\item \textbf{Consequences at $\ell$.} For $\ell$ odd, $\ell\nmid N$: if $\bar\rho$ is reducible, $E$ is ordinary at $\ell$ and the two characters
restrict to $I_\ell$ as $\omega$ and $1$; if $\bar\rho\simeq\bar\rho\otimes\chi_{\ell^*}$, then $\ell=3$ and $E$ is ordinary at $3$. \emph{(proved here, given
\Cref{thm:galrep:deligne-fontaine})}
\Cref{cor:galrep:ordinary}.
\item \textbf{Serre conductor.} For $\bar\sigma\colon G_\Q\to\GL_n(\Fbar_\ell)$ and $p\ne\ell$, $a_p(\bar\sigma)=n-\dim\bar\sigma^{I_p}+\operatorname{sw}_p(\bar\sigma)$ and
$N(\bar\sigma)=\prod_{p\ne\ell}p^{a_p(\bar\sigma)}$; $a_p(\bar\sigma)=0$ iff $\bar\sigma$ is unramified at $p$; $a_p$ is invariant under extension of scalars and
twists unramified at $p$. \emph{(definition; proved here)} \Cref{def:galrep:serre-conductor},
\Cref{prop:galrep:serre-conductor-basic}.
\item \textbf{Serre conductor of $\bar\rho_{E,\ell}$.} For $p\ne\ell$, $a_p(\bar\rho_{E,\ell})=f_p(E)-\dim_{\Fl}\Phig_p[\ell]$. So $N(\bar\rho_{E,\ell})$ divides
the prime-to-$\ell$ part of $N$; the exponent at a multiplicative $p$ is $0$ or $1$ as $\ell\mid n_p$ or not; at an additive $p$ and $\ell$
odd it is $f_p-1$ for $\ell=3$ and types $\kod{IV}$, $\kod{IV}^*$, and $f_p$ otherwise. \emph{(proved here)}
\Cref{prop:galrep:serre-conductor-E}.
\item \textbf{Open image.} For $E$ without CM, $\rho_{E,\ell}(G_\Q)$ is open in $\GL_2(\Zl)$ and $\bar\rho_{E,\ell}$ is surjective for almost all $\ell$; with
CM by an order in $K$ and $\ell$ odd and unramified in $K$ the image of $\bar\rho_{E,\ell}$ lies in the normaliser of a Cartan subgroup and not in the Cartan
subgroup; a rational $\ell$-isogeny exists only for $\ell\in\{2,3,5,7,11,13,17,19,37,43,67,163\}$. \emph{(stated; reference)}
\Cref{thm:galrep:open-image}.
\item \textbf{Constraints on images.} For $\ell$ odd, $H=\bar\rho(G_\Q)$ has $\det H=\Fl^\times$ and contains a conjugate of $\operatorname{diag}(1,-1)$; so
$H\not\subseteq\SL_2(\Fl)$, $H\not\subseteq C_{ns}(\Fl)$, and $H$ is $\GL_2(\Fl)$, or lies in a Borel subgroup, or in the normaliser of a Cartan
subgroup (not in the Cartan subgroup if $\bar\rho$ is irreducible), or has projective image $A_4$, $S_4$ or $A_5$. \emph{(proved here)}
\Cref{prop:galrep:image-constraints}.
\item \textbf{The mod $3$ images.} For every $E/\Q$, CM or not, $\bar\rho_{E,3}(G_\Q)$ is conjugate to exactly one of $\GL_2(\F_3)$, $N_{ns}(\F_3)$
(\texttt{3Nn}), $N_s(\F_3)$ (\texttt{3Ns}), $\mathrm B(\F_3)$ (\texttt{3B}), $\sm{1}{*}{0}{*}$ (\texttt{3B.1.1}), $\sm{*}{*}{0}{1}$ (\texttt{3B.1.2}), $C_s(\F_3)$
(\texttt{3Cs}), $\{\operatorname{diag}(1,\pm1)\}$ (\texttt{3Cs.1.1}); never to $C_{ns}(\F_3)$. The orbits on the four lines are the factorisation pattern of
$\psi_3$. \emph{(proved here; counts below conductor $10{,}000$)} \Cref{thm:galrep:mod3-images}.
\item \textbf{$3$-isogeny paths.} For every $E/\Q$, CM or not, the rational $3$-isogenies make each isogeny class a disjoint union of paths; the
ends of a path of length at least one have one stable line in $E[3]$, of characters $\eta$ and $\chi_{-3}\eta^{-1}$, the interior vertices two. So
the ends pair the images \texttt{3B.1.1} and \texttt{3B.1.2}, which occur equally often in every isogeny class. \emph{(proved here)}
\Cref{prop:galrep:3-isogeny-paths}.
\item \textbf{Labels.} Sutherland's labels name the subgroup; the dictionary for $\ell=3$ and the labels occurring at $\ell=5,7,11,13$
below conductor $10{,}000$. \emph{(convention; data)} \Cref{conv:galrep:labels}.
\item \textbf{Zywina.} For non-CM $E/\Q$ and $\ell\le13$ the possible images of $\bar\rho_{E,\ell}$ are listed, with the corresponding families of
$j$-invariants. \emph{(stated; reference)} \Cref{thm:galrep:zywina}.
\item \textbf{Dihedral representations.} For $\ell$ odd, $\bar\rho$ irreducible and $K$ quadratic: $\bar\rho|_{G_K}$ not absolutely irreducible $\iff$
$\bar\rho\simeq\bar\rho\otimes\chi_K$ $\iff$ $\bar\rho\otimes\Fbar_\ell\simeq\Ind_{G_K}^{G_\Q}\psi$ $\iff$ $\ell\mid a_p$ for almost all $p$ inert in $K$; then every
isomorphism $\bar\rho\to\bar\rho\otimes\chi_K$ has trace $0$ and $H^0(\Q,\ad\bar\rho\otimes\chi_K)$ is a line. \emph{(proved here)}
\Cref{prop:galrep:dihedral}.
\item \textbf{$\Q(\sqrt{\ell^*})$ is automatic for $\ell\ge5$ (Proposition prop:Sf).} For $\ell$ odd, $\ell\nmid N$, $\bar\rho$ irreducible: if $\ell\ge5$,
$\bar\rho$ is absolutely irreducible on $G_{\Q(\sqrt{\ell^*})}$; if $\ell=3$ and it is not, $E$ is ordinary at $3$. \emph{(proved here, given
\Cref{thm:galrep:deligne-fontaine})}
\Cref{thm:galrep:lstar-automatic}.
\item \textbf{Proposition prop:ns.} If $E[3]$ is irreducible, the image is $\GL_2(\F_3)$, $N_{ns}(\F_3)$ or $N_s(\F_3)$; $\bar\rho_{E,3}$ is absolutely
irreducible on $G_{\Q(\sqrt{-3})}$ unless the image is $N_s(\F_3)$, and then $\bar\rho_{E,3}\otimes\F_9\simeq\Ind\psi$ with $\psi$ of order $4$,
$\bar\rho_{E,3}\simeq\bar\rho_{E,3}\otimes\chi_{-3}$, and $h_3(E)=1$; $\psi_3$ is a product of two irreducible quadratics in that case and irreducible
otherwise. CM allowed. \emph{(proved here)} \Cref{prop:galrep:ns}.
\item \textbf{Lemma lem:h0.} For $\ell$ odd and $\bar\rho$ irreducible, $H^0(\Q,\ad T(1)\otimes\Qlmod)\simeq\Z/\ell^{h_\ell(E)}$ with $h_\ell(E)$ the largest $k$
such that $\rho\bmod\ell^k\simeq(\rho\otimes\chi_\ell)\bmod\ell^k$; $h_\ell(E)=0$ for every $\ell\ge5$ and $h_3(E)\le1$, with $h_3(E)=1$ iff
$\bar\rho_{E,3}\simeq\bar\rho_{E,3}\otimes\chi_{-3}$ iff the image is $N_s(\F_3)$; for every $E$, CM or not. \emph{(proved here)}
\Cref{lem:galrep:h0}.
\item \textbf{Proposition prop:ns-structure.} If the image of $\bar\rho_{E,3}$ is $N_s(\F_3)$: $\bar\rho(I_p)\subseteq\langle J\rangle$ for $p\ne3$; every
potentially multiplicative $p\ne3$ has $3\mid v_p(j)$; no $p\ne3$ has type $\kod{IV}$, $\kod{IV}^*$ and no $p\ge5$ type $\kod{II}$, $\kod{II}^*$; the
set $R$ of primes $p\ne3$ with $\rho_{E,3}(I_p)$ cyclic of order $4$ is nonempty and is the set of primes $p\ne3$ ramified in either quadratic
field of $\psi_3$; for $p\in R$, $\Frob_p$ acts on $(\ad V)^{I_p}$ by $\chi_{-3}(p)$ and $e_3(E,p)=v_3(p-\chi_{-3}(p))\ge1$; if $p\in R$, $p\ge5$, the type is
$\kod{III}$ or $\kod{III}^*$ and $p\equiv\pm1\pmod{12}$. \emph{(proved here)} \Cref{prop:galrep:ns-structure}.
\item \textbf{$S_f$.} For $\ell$ odd, $\ell\nmid N$: $\ell\notin S_f$ $\iff$ $E[\ell]$ irreducible and $h_\ell(E)=0$ $\iff$ $E[\ell]$ irreducible and not
($\ell=3$ and image $N_s(\F_3)$). \emph{(proved here, given \Cref{thm:galrep:deligne-fontaine} for $\ell\ge5$)} \Cref{cor:galrep:Sf}.
\item \textbf{Reducible $\bar\rho$.} For $\ell$ odd, $\ell\nmid N$, $E[\ell]$ reducible: $\bar\rho^{\mathrm{ss}}=\eta_1\oplus\eta_2$, $\eta_1\eta_2=\bar\chi_\ell$, $E$ ordinary
at $\ell$, $\{\eta_i|_{I_\ell}\}=\{\omega,1\}$, both unramified at multiplicative primes; if $E$ is semistable, $\{\eta_1,\eta_2\}=\{1,\bar\chi_\ell\}$,
$a_p\equiv1+p\pmod\ell$, and $E$ or an $\ell$-isogenous curve has a rational point of order $\ell$. \emph{(proved here, given
\Cref{thm:galrep:deligne-fontaine})}
\Cref{prop:galrep:reducible}.
\item \textbf{Eisenstein labels at $3$.} $3\mid\#E(\Q)_{\mathrm{tors}}$ iff the image is \texttt{3B.1.1} or \texttt{3Cs.1.1}; \texttt{3Cs.1.1} means
$E[3]\simeq\Z/3\oplus\mu_3$, not the split Cartan subgroup; without CM such a curve has no rational point of order $9$. \emph{(proved here,
the last statement citing \Cref{ch:ellcurves})} \Cref{prop:galrep:eis-labels}.
\item \textbf{Optimal level.} An odd irreducible $\bar\sigma\colon G_\Q\to\GL_2(\Fbar_\ell)$ is modular of level $N(\bar\sigma)$ and weight $k(\bar\sigma)$
(Khare-Wintenberger); for $\ell\nmid N(E')$, $E'[\ell]$ irreducible, and image not $N_s(\F_3)$ if $\ell=3$, $\bar\rho_{E',\ell}$ arises from a newform of
weight $2$ on $\Gamma_0(N(\bar\rho_{E',\ell}))$. \emph{(stated; reference; the elliptic-curve inputs proved here)}
\Cref{thm:galrep:optimal-level}.
\end{enumerate}

\section*{Where the paper uses this chapter}

\begin{itemize}
\item \textbf{Section 1.1} ($V$, $T$, $\rho$, $\bar\rho$, $A_f=\ad V$, $\delta_p=\dim V^{I_p}$, the arithmetic and geometric Frobenius):
\Cref{conv:galrep:notation}, \Cref{def:galrep:rho}, \Cref{prop:galrep:good}, \Cref{cor:galrep:delta}.
\item \textbf{The hypothesis of Theorem thm:main} ($\ell\notin S_f$) and \textbf{Proposition prop:Sf}: \Cref{prop:galrep:irreducibility},
\Cref{prop:galrep:dihedral}, \Cref{thm:galrep:lstar-automatic}, \Cref{cor:galrep:Sf}.
\item \textbf{Definition def:terms} ($t_\ell(E,q)$, potentially multiplicative primes, $\kod{I}_n^*$ at $2$) and \textbf{Section 1.5, Remark rem:ko}
(the Tamagawa exponent): \Cref{lem:galrep:tamagawa-exp}, \Cref{prop:galrep:potmult}.
\item \textbf{Lemma lem:tamagawa-exp}: \Cref{thm:galrep:tate-shape}, \Cref{lem:galrep:tamagawa-exp}.
\item \textbf{Lemma lem:IV} and \textbf{Remark rem:types}: \Cref{lem:galrep:component-invariants}, \Cref{thm:galrep:serre-tate},
\Cref{lem:galrep:IV}, \Cref{cor:galrep:IV-at-2}, \Cref{prop:galrep:mod3-inertia}.
\item \textbf{Lemma lem:euler, Proposition prop:shapes, Lemma lem:tamadd} (the images of inertia and their invariants on $V$ and $\ad V$):
\Cref{thm:galrep:serre-tate}, \Cref{thm:galrep:tate-shape}, \Cref{prop:galrep:potmult}, \Cref{thm:galrep:additive-shapes}, \Cref{lem:galrep:IV};
the $L$-factors themselves are computed in \Cref{ch:lfunctions}.
\item \textbf{Proposition prop:forms, Remark rem:wild, Section sec:wild} (the types $\kod{IV}$, $\kod{IV}^*$ at $2$): \Cref{lem:galrep:IV},
\Cref{cor:galrep:IV-at-2}, \Cref{rem:galrep:wild-inertia}.
\item \textbf{Lemma lem:old} (the local analysis at $p\nmid\ell N$, at $p=\ell$, at $p\,\|\,N$ ramified or not): \Cref{prop:galrep:irreducibility}(iii),
\Cref{thm:galrep:tate-shape}, \Cref{lem:galrep:tamagawa-exp}, \Cref{thm:galrep:deligne-fontaine}.
\item \textbf{The optimal-level paragraph of Section 3 and Theorem thm:qual} (cases (a), (b), (c), the level $M=N(\bar\rho)$):
\Cref{thm:galrep:good-at-l}, \Cref{def:galrep:serre-conductor}, \Cref{prop:galrep:serre-conductor-E}, \Cref{lem:galrep:IV},
\Cref{prop:galrep:mod3-inertia}, \Cref{thm:galrep:optimal-level}, \Cref{ex:galrep:level-lowering}.
\item \textbf{The proof of Theorem thm:main} (the vanishing of $H^0(\Q,W^D)$; the potentially multiplicative primes, $2$ included):
\Cref{lem:galrep:h0}, \Cref{prop:galrep:potmult}.
\item \textbf{Theorem thm:quant} ($\ad T$ is unchanged by the twist): \Cref{prop:galrep:twist-isogeny}.
\item \textbf{Proposition prop:ns}, \textbf{Proposition prop:ns-structure}, \textbf{Remark rem:ns}, \textbf{Section sec:anomaly}, \textbf{Conjecture
conj:euler} ($h_\ell(E)$): \Cref{thm:galrep:mod3-images}, \Cref{prop:galrep:ns}, \Cref{lem:galrep:h0}, \Cref{prop:galrep:ns-structure}.
\item \textbf{Section sec:eis, Proposition prop:no9, Conjecture conj:eis}: \Cref{prop:galrep:reducible}, \Cref{prop:galrep:eis-labels}.
\item \textbf{Section sec:data and the appendix} (the mod $3$ images checked against $\psi_3$; the function \texttt{ns3}):
\Cref{thm:galrep:mod3-images}, \Cref{rem:galrep:computing-images}.
\end{itemize}

\section{The representations \texorpdfstring{$\rho_{E,\ell}$ and $\bar\rho_{E,\ell}$}{rho and rhobar}}\label{sec:galrep:reps}

This section fixes the objects of the paper's Section 1.1 and collects their basic properties: the determinant, complex
conjugation, the effect of twists and isogenies, and the meaning of irreducibility. Most proofs were given in
\Cref{ch:galois} and \Cref{ch:ellcurves}; here they are assembled in the form the paper uses them. The paper uses these facts in the
hypothesis of Theorem thm:main, in Proposition prop:ns (``$\det\bar\rho_{E,3}=\chi_{-3}$ cuts out $\Q(\sqrt{-3})$'', complex conjugation),
in Lemma lem:h0 and Lemma lem:old (``irreducible, hence absolutely irreducible'') and in Theorem thm:quant (twists).

\begin{convention}[notation of this chapter]\label{conv:galrep:notation}
\begin{enumerate}[label=(\alph*)]
\item \emph{Objects.} $E/\Q$ is given by a global minimal model and has conductor $N$; $p$, $q$, $\ell$ are primes, $\ell$ is the coefficient
prime (odd unless said otherwise) and $q$ usually a bad prime; $\ell^*=(-1)^{(\ell-1)/2}\ell$, $q^*=(-1)^{(q-1)/2}q$ for odd $\ell,q$. The letter $q$ is
never the variable of a $q$-expansion in this chapter; the Tate parameter of $E$ at $q$ is $q_E$ (never $u$), and $u=q_Eq^{-v_q(q_E)}$ is its unit
part.
\item \emph{Local data.} $D_q\simeq G_{\Q_q}$ and $I_q$ are a decomposition and an inertia group at $q$ for a fixed embedding $\Qbar\to\Qbar_q$
(\Cref{ch:ant}); $f_q=f_q(E)$ is the conductor exponent and $\operatorname{sw}_q$ the Swan conductor; at a multiplicative prime $n_q=v_q(\Delta_{\min})$;
$\Phi_q(\cdot)$ is the component group scheme of the N\'eron model at $q$, always written with an argument, $\Phig_q=\Phi_q(\Fbar_q)$, and
$c_q=\#\Phi_q(\F_q)$ divides $\#\Phig_q$ (\Cref{ch:reduction}).
\item \emph{Frobenius.} $\Frob_p$ is an \emph{arithmetic} Frobenius at $p$ (inducing $x\mapsto x^p$ on residue fields), well defined modulo $I_p$
and up to conjugation; the geometric Frobenius is written $\Frob_p^{-1}$, never $\Phi_p$, because $\Phi_q(\cdot)$ is the component group here.
Group-theoretic statements (conjugation, the Tate matrices, the tame relation $\Frob\,t\,\Frob^{-1}=t^q$) use $\Frob$; $L$-factors use
$\Frob^{-1}$ and are the business of \Cref{ch:lfunctions}, with $L_q(A_f,s)=\det(1-\Frob_q^{-1}q^{-s}\mid(\ad V)^{I_q})^{-1}$,
$L_q(\Sym f,s)=L_q(A_f,s-1)=P_q(q^{-s})^{-1}$ and $e_\ell(E,q)=v_\ell\bigl(L_q(A_f,1)^{-1}/L^{\nv}_q(A_f,1)^{-1}\bigr)$.
\item \emph{Characters.} $\chi_\ell\colon G_\Q\to\Zl^\times$ is the $\ell$-adic cyclotomic character and $\bar\chi_\ell$ its reduction; a letter
subscript ($\chi_\ell$, $\chi_p$) always means a cyclotomic character, and at a numerical prime we write $\chi_3^{\mathrm{cyc}}$. For a non-square
$d\in\Q^\times$, $\chi_d$ is the quadratic character of $\Q(\sqrt d)$, written with the discriminant ($\chi_{-3}$, $\chi_{-4}$, $\chi_{-7}$,
$\chi_{-11}$); the quadratic character of conductor $q$ is $\chi_{q^*}$, never $\chi_q$. Since $\Q(\sqrt{\ell^*})\subset\Q(\zeta_\ell)$,
$\chi_{\ell^*}=\bar\chi_\ell^{(\ell-1)/2}$ (\Cref{ch:galois}); in particular $\bar\chi_3=\chi_{-3}$, while $\chi_{\ell^*}\ne\bar\chi_\ell$ for $\ell\ge5$. The letter
$\varepsilon$ denotes an unramified character of $G_{\Q_q}$ of order at most $2$, and $\omega$, $\omega_2$ the fundamental characters of
\Cref{def:galrep:fundamental-characters}. $E^{(d)}=E\otimes\chi_d$ is the quadratic twist, and $E^{(q)}=E^{(q^*)}$ for odd $q$.
\item \emph{Matrices.} As in \Cref{ch:groups}: $J=\sm{0}{1}{-1}{0}$ (the paper's $J$), $w=\sm0110$, $u=\sm1101$, $h=\operatorname{diag}(1,-1)$; the
subgroups $C_s$, $N_s$, $C_{ns}$, $N_{ns}$, $\mathrm B$ (upper triangular), $\mathrm U=\langle u\rangle$, $Q_8$ and $C_4=\langle J\rangle$ of $\GL_2(\F_3)$; with
$\F_9=\F_3[i]$, $J$ acts as multiplication by $-i$. The letter $\sigma$ denotes an element of order $3$ of an image of inertia (as in Lemmas
lem:IV and lem:euler) and $\tau$ one with $\tau^2=-1$.
\item \emph{Labels.} Curves carry Cremona's labels (\curve{11a1}); mod $\ell$ images carry Sutherland's labels (\texttt{3Ns},
\texttt{3B.1.1}), always together with the subgroup they name (\Cref{conv:galrep:labels}). The basis of $\ad V$ attached to a basis
$e_1,e_2$ of $V$ is $e=e_1\otimes e_2^\vee$, $h=e_1\otimes e_1^\vee-e_2\otimes e_2^\vee$, $f=e_2\otimes e_1^\vee$; $G$ acts on $\ad V$ by conjugation.
\end{enumerate}
\end{convention}

\begin{definition}\label{def:galrep:rho}
Let $E/\Q$ be an elliptic curve and $\ell$ a prime. The \emph{$\ell$-adic Tate module} is $T=T_\ell E=\varprojlim_kE[\ell^k]$ (transition maps
multiplication by $\ell$), free of rank $2$ over $\Zl$, and $V=V_\ell E=T\otimes_{\Zl}\Ql$ (\Cref{ch:ellcurves}). The group $G_\Q$ acts on
$E[\ell^k]\subset E(\Qbar)$ through the action on coordinates, compatibly with the transition maps; this defines
\[
\rho=\rho_{E,\ell}\colon G_\Q\to\Aut_{\Zl}(T)\simeq\GL_2(\Zl),
\]
its reductions $\rho\bmod\ell^k\colon G_\Q\to\Aut(E[\ell^k])\simeq\GL_2(\Z/\ell^k)$, $E[\ell^k]=T/\ell^kT$, and $\bar\rho=\bar\rho_{E,\ell}=\rho\bmod\ell$ on $E[\ell]$.
The \emph{adjoint} is $A=\ad T$, the trace-zero endomorphisms of $T$ with $G_\Q$ acting by conjugation, and $A_f=\ad V=A\otimes\Ql$ (the
paper's notation). For a prime $p\ne\ell$, $\delta_p=\dim_{\Ql}V^{I_p}$.
\end{definition}

Why these objects. A single $E[\ell^k]$ is a finite Galois module, and it forgets how the torsion of different levels fits together; the Tate
module packages all levels at once into a lattice $T$ over the complete discrete valuation ring $\Zl$, on which $G_\Q$ acts by $2\times2$
matrices. Passing to $V=T\otimes\Ql$ gives a vector space, where semisimplicity and eigenvalues behave as in linear algebra over a field of
characteristic $0$; the lattice $T$ itself keeps the integral information that the paper's divisibility statements are about. The adjoint
$\ad V$ appears because the paper's $L$-function is $L(A_f,s)=L(\ad V,s)$, and the integer $\delta_p$ appears because the local factor of an
$L$-function at $p$ only sees the inertia invariants (\Cref{conv:galrep:notation}(c)); \Cref{cor:galrep:delta} computes $\delta_p$ for every
reduction type.

A small example shows what the action looks like at the bottom level. Write $E$ as $y^2=g(x)$ with $g\in\Q[x]$ a separable cubic (possible for every
$E/\Q$ after a change of coordinates). Then $E[2]=\{O,(e_1,0),(e_2,0),(e_3,0)\}$ with $e_1,e_2,e_3$ the roots of $g$, and $G_\Q$ acts on the three
nonzero points by permuting the roots. Since $\GL_2(\F_2)$ is the group of all permutations of the three nonzero vectors of $\F_2^2$, the image of
$\bar\rho_{E,2}$ is, as a permutation group, the Galois group of $g$: trivial if $g$ splits, of order $2$ if it has exactly one rational root, and
$C_3$ or $S_3$ if it is irreducible. For $\ell$ odd the analogue of $g$ is the $\ell$-division polynomial, and for $\ell=3$ it is used throughout
\Cref{sec:galrep:images}.

\begin{proposition}[basic properties]\label{prop:galrep:rho-basic}
Let $E/\Q$ and $\ell$ be as above.
\begin{enumerate}[label=(\roman*)]
\item $\rho$ is continuous, and the kernel of $\rho\bmod\ell^k$ is $G_{\Q(E[\ell^k])}$, where $\Q(E[\ell^k])$ is the finite Galois extension generated
by the coordinates of the points of $E[\ell^k]$.
\item $\det\rho=\chi_\ell$. Hence $\Q(\zeta_\ell)\subseteq\Q(E[\ell])$, $\det\bar\rho=\bar\chi_\ell$ is surjective onto $\Fl^\times$, and
\[
\bar\rho\bigl(G_{\Q(\sqrt{\ell^*})}\bigr)=\{g\in\bar\rho(G_\Q):\det g\in(\Fl^\times)^2\}\qquad(\ell\text{ odd}).
\]
For $\ell=3$: $\det\bar\rho_{E,3}=\chi_{-3}$, the field it cuts out is $\Q(\sqrt{-3})$, and $\bar\rho_{E,3}(G_{\Q(\sqrt{-3})})=\bar\rho_{E,3}(G_\Q)\cap\SL_2(\F_3)$.
\item Let $\ell$ be odd and $c\in G_\Q$ a complex conjugation. Then $T=T^{c=1}\oplus T^{c=-1}$ with both summands free of rank one, so $\rho(c)$ is
conjugate in $\GL_2(\Zl)$ to $\operatorname{diag}(1,-1)$; in particular $\bar\rho(c)$ has determinant $-1$, trace $0$ and the two distinct
eigenvalues $1$, $-1$.
\item $\rho$ is unramified at every prime $p\nmid N\ell$.
\end{enumerate}
\end{proposition}

\begin{proof}
(i) Each $E[\ell^k]$ is a finite set of points with algebraic coordinates, so the action of $G_\Q$ on it factors through the finite quotient
$\Gal(\Q(E[\ell^k])/\Q)$, and the stabiliser of every point is open. The preimage under $\rho$ of the basic open subgroup $1+\ell^kM_2(\Zl)$ of $\GL_2(\Zl)$ is the kernel of $\rho\bmod\ell^k$, which is
therefore open; this is continuity (\Cref{ch:galois}). The kernel of the action on
$E[\ell^k]$ consists of the $g$ fixing all coordinates of all $\ell^k$-torsion points, which is $G_{\Q(E[\ell^k])}$; the extension is Galois because
$E[\ell^k]$ is $G_\Q$-stable.

(ii) The Weil pairing $e_{\ell^k}\colon E[\ell^k]\times E[\ell^k]\to\mu_{\ell^k}$ is bilinear, alternating, nondegenerate and Galois equivariant, and the
pairings are compatible as $k$ varies (\Cref{ch:ellcurves}). For a basis $P,Q$ of $E[\ell^k]$, $\zeta=e_{\ell^k}(P,Q)$ is a primitive $\ell^k$-th root
of unity, and for $g\in G_\Q$ with $\rho(g)=\sm abcd$ in that basis,
\[
g(\zeta)=e_{\ell^k}(gP,gQ)=e_{\ell^k}(aP+cQ,bP+dQ)=\zeta^{ad-bc},
\]
so $\det(\rho(g)\bmod\ell^k)=\chi_\ell(g)\bmod\ell^k$ for every $k$, whence $\det\rho=\chi_\ell$. An element fixing $E[\ell]$ pointwise fixes $\zeta_\ell=e_\ell(P,Q)$,
so $\Q(\zeta_\ell)\subseteq\Q(E[\ell])$. The reduction $\bar\chi_\ell$ is the character $G_\Q\to\Gal(\Q(\zeta_\ell)/\Q)\simeq\Fl^\times$, which is surjective. For $\ell$
odd, $\Q(\sqrt{\ell^*})$ is the fixed field in $\Q(\zeta_\ell)$ of the squares of $\Fl^\times$ (\Cref{ch:galois}); so $g\in G_{\Q(\sqrt{\ell^*})}$ iff
$\bar\chi_\ell(g)=\det\bar\rho(g)$ is a square. Both inclusions of the displayed equality follow. If $g\in G_{\Q(\sqrt{\ell^*})}$, then $\bar\rho(g)$ is an
element of $\bar\rho(G_\Q)$ with square determinant. Conversely, an element $h$ of $\bar\rho(G_\Q)$ is $\bar\rho(g)$ for some $g\in G_\Q$, and if $\det h$
is a square, then $\bar\chi_\ell(g)=\det h$ is a square, so this very $g$ lies in $G_{\Q(\sqrt{\ell^*})}$ and $h\in\bar\rho(G_{\Q(\sqrt{\ell^*})})$. (The point is that
membership of $g$ in $G_{\Q(\sqrt{\ell^*})}$ depends only on $\det\bar\rho(g)$, so it does not matter which preimage of $h$ is chosen.) For $\ell=3$, $\F_3^\times=\{\pm1\}$, its squares are $\{1\}$, and $\bar\chi_3(g)=\pm1$ according as $g(\zeta_3)=\zeta_3^{\pm1}$, that is,
according as $g$ fixes $\sqrt{-3}=2\zeta_3+1$ or not: $\bar\chi_3=\chi_{-3}$.

(iii) Since $c^2=1$, $\rho(c)^2=1$, and $\det\rho(c)=\chi_\ell(c)=-1$ because $c(\zeta)=\zeta^{-1}$ on all roots of unity. As $2\in\Zl^\times$, the
endomorphisms $\operatorname{pr}_\pm=\frac12(1\pm\rho(c))$ of $T$ are complementary idempotents, so $T=\operatorname{pr}_+T\oplus\operatorname{pr}_-T$ with $\operatorname{pr}_\pm T=T^{c=\pm1}$. Both are direct
summands of a free $\Zl$-module of rank $2$, hence free; their ranks add to $2$; and they cannot be $2$ and $0$, since then $\rho(c)=\pm1$ would
have determinant $1$. So both have rank one, and a basis $e_+,e_-$ adapted to the decomposition gives $\rho(c)=\operatorname{diag}(1,-1)$. Reducing
modulo $\ell$ gives the statement on $\bar\rho(c)$; $1\ne-1$ in $\Fl$ as $\ell$ is odd.

(iv) The conductor exponent $f_p$ is defined in \Cref{ch:reduction} as the Artin conductor of $V$, so $p\nmid N$ means
$\dim V-\dim V^{I_p}+\operatorname{sw}_p(V)=0$; both terms are nonnegative, so $V^{I_p}=V$, and $I_p$ acts trivially on $V$, hence on $T\subset V$
(see also \Cref{thm:galrep:nos} below).
\end{proof}

\begin{theorem}[Faltings]\label{thm:galrep:faltings}
Let $E$, $E'$ be elliptic curves over $\Q$ and $\ell$ a prime. Then $V_\ell E$ is a semisimple $\Ql[G_\Q]$-module, and the natural map
$\Hom_\Q(E,E')\otimes\Zl\to\Hom_{G_\Q}(T_\ell E,T_\ell E')$ is an isomorphism. In particular $E$ and $E'$ are isogenous over $\Q$ iff $V_\ell E\simeq V_\ell E'$ as
$G_\Q$-modules, iff $a_p(E)=a_p(E')$ for almost all $p$.
\end{theorem}

We do not prove this; see Faltings \cite{Faltings83} (the Tate conjecture for abelian varieties over number fields). The last equivalence
follows from the first two with Chebotarev's theorem and the Brauer-Nesbitt theorem (\Cref{ch:galois}), since semisimple representations
in characteristic $0$ are determined by their traces. The book uses \Cref{thm:galrep:faltings} only as background (for instance in
\Cref{ch:ant}); every argument about the paper's curves uses $\bar\rho$, for which the needed statements are proved below.

\begin{proposition}[twists and isogenies]\label{prop:galrep:twist-isogeny}
Let $E/\Q$ and let $\ell$ be odd.
\begin{enumerate}[label=(\roman*)]
\item For $d\in\Q^\times$ not a square, $\rho_{E^{(d)},\ell}\simeq\rho_{E,\ell}\otimes\chi_d$. Hence $\ad\rho_{E^{(d)},\ell}\simeq\ad\rho_{E,\ell}$, the images of
$\bar\rho_{E^{(d)},\ell}$ and $\bar\rho_{E,\ell}$ in $\PGL_2(\Fl)$ coincide, and $E^{(d)}[\ell]$ is irreducible iff $E[\ell]$ is.
\item For $\ell=3$ (the labels are those of \Cref{thm:galrep:mod3-images}): twisting by $\chi_{-3}$ exchanges \texttt{3B.1.1} and \texttt{3B.1.2}
and fixes every other image; twisting by $\chi_d$ with $d\not\equiv1,-3$ modulo $\Q^{\times2}$ sends \texttt{3B.1.1} and \texttt{3B.1.2} to
\texttt{3B}, sends \texttt{3Cs.1.1} to \texttt{3Cs}, and fixes \texttt{3Ns}, \texttt{3Nn} and the surjective image.
\item If $E[\ell]$ is irreducible, every curve $E'$ isogenous to $E$ over $\Q$ has $T_\ell E'\simeq T_\ell E$ as $\Zl[G_\Q]$-modules; so $E$ and $E'$ have
the same $\rho$, $\bar\rho$, image label at $\ell$, adjoint $\ad T$, and $h_\ell$ (\Cref{lem:galrep:h0}).
\end{enumerate}
\end{proposition}

\begin{proof}
(i) The first statement is proved in \Cref{ch:ellcurves}: an isomorphism $E\to E^{(d)}$ defined over $\Q(\sqrt d)$ transforms the action of
$g\in G_\Q$ into $\chi_d(g)$ times it. Concretely, take $E\colon y^2=x^3+ax+b$ and $E^{(d)}\colon dy^2=x^3+ax+b$, and $\phi(x,y)=(x,y/\sqrt d)$, an
isomorphism $E\to E^{(d)}$ over $\Q(\sqrt d)$. For $P=(x,y)\in E(\Qbar)$ and $g\in G_\Q$,
\[
g\bigl(\phi(P)\bigr)=\Bigl(g(x),\frac{g(y)}{g(\sqrt d)}\Bigr)=\Bigl(g(x),\frac{g(y)}{\chi_d(g)\sqrt d}\Bigr)=[\chi_d(g)]\,\phi\bigl(g(P)\bigr),
\]
because $[-1](x,y)=(x,-y)$ on $E^{(d)}$. So, transported to $E$ by $\phi$, the action of $g$ on $E^{(d)}[\ell^k]$ is $\chi_d(g)\rho(g)$, compatibly in $k$. Conjugation by $\chi_d(g)\rho(g)$ equals conjugation by $\rho(g)$, so the adjoint representations agree; and
$\chi_d(g)\bar\rho(g)$ and $\bar\rho(g)$ have the same image in $\PGL_2(\Fl)$. A line stable under $\bar\rho$ is stable under $\bar\rho\otimes\chi_d$, and conversely.

(ii) Write $H$ for the image of $\bar\rho=\bar\rho_{E,3}$ and $H'$ for that of $\bar\rho\otimes\chi_d$; then $H'\subseteq H\cdot\{\pm1\}$ and $H\subseteq H'\cdot\{\pm1\}$, because
$(\bar\rho\otimes\chi_d)(s)=\chi_d(s)\bar\rho(s)$ with $\chi_d(s)=\pm1$ for every $s\in G_\Q$.
If $H$ contains an element $g$ with $g^2=-1$, choose $s\in G_\Q$ with $\bar\rho(s)=g$: then $((\bar\rho\otimes\chi_d)(s))^2=g^2=-1$, so $-1\in H'$, hence
$H'\subseteq H\{\pm1\}=H$ and $H\subseteq H'\{\pm1\}=H'$, and $H'=H$. The groups $\GL_2(\F_3)$, $N_{ns}(\F_3)$ and $N_s(\F_3)$ all contain $J$, with $J^2=-1$; so
the surjective image, \texttt{3Nn} and \texttt{3Ns} are fixed by every twist. In the reducible cases write
$\bar\rho\sim\sm{\eta_1}{*}{0}{\eta_2}$ with $\eta_1\eta_2=\chi_{-3}$ (\Cref{prop:galrep:rho-basic}(ii)); then $\bar\rho\otimes\chi_d\sim\sm{\eta_1\chi_d}{*}{0}{\eta_2\chi_d}$,
and the twist is split iff $\bar\rho$ is (by (i), the two representations have the same stable lines). By \Cref{thm:galrep:mod3-images}, \texttt{3B.1.1} is the non-split case with $\eta_1=1$, \texttt{3B.1.2} the
non-split case with $\eta_1=\chi_{-3}$, \texttt{3B} the non-split case with $\eta_1\notin\{1,\chi_{-3}\}$, and \texttt{3Cs.1.1}, \texttt{3Cs} the split cases
with $\{\eta_1,\eta_2\}=\{1,\chi_{-3}\}$, respectively $\{\eta_1,\eta_2\}\ne\{1,\chi_{-3}\}$. Twisting by $\chi_{-3}$ sends $\eta_1=1$ to $\chi_{-3}$ and $\chi_{-3}$ to $1$,
and sends a character outside $\{1,\chi_{-3}\}$ to another such; twisting by $\chi_d$ with $\chi_d\notin\{1,\chi_{-3}\}$ sends both $1$ and $\chi_{-3}$ outside
$\{1,\chi_{-3}\}$. This gives every case.

(iii) Over $\Q$ every isogeny between curves of the class is a composite of isogenies of prime degree with $G_\Q$-stable cyclic kernels
(\Cref{ch:ellcurves}). An isogeny of degree $\ell$ would have a $G_\Q$-stable kernel of order $\ell$ inside $E[\ell]$ (or inside $E''[\ell]$ for an
intermediate curve $E''$, whose $\ell$-torsion is isomorphic to $E[\ell]$ by induction), contradicting irreducibility. So $E$ and $E'$ are linked
by isogenies $\varphi$ of degree prime to $\ell$, and such a $\varphi$ induces a $G_\Q$-equivariant injection $T_\ell E\to T_\ell E'$ whose cokernel has order the
$\ell$-part of $\deg\varphi$, which is $1$; so it is an isomorphism. In detail: $\varphi$ is defined over $\Q$, so the map $\varphi_*\colon T_\ell E\to T_\ell E'$ it induces on
compatible systems of torsion points commutes with $G_\Q$; the dual isogeny satisfies $\hat\varphi\circ\varphi=[\deg\varphi]$ on $E$ and $\varphi\circ\hat\varphi=[\deg\varphi]$ on $E'$
(\Cref{ch:ellcurves}); and multiplication by the integer $\deg\varphi$, a unit of $\Zl$, is an automorphism of both Tate modules. So $\hat\varphi_*\circ\varphi_*$ and
$\varphi_*\circ\hat\varphi_*$ are invertible, and $\varphi_*$ is a $G_\Q$-isomorphism. The other objects are built from $T$ alone.
\end{proof}

\begin{proposition}[irreducibility]\label{prop:galrep:irreducibility}
Let $\ell$ be odd.
\begin{enumerate}[label=(\roman*)]
\item $\bar\rho_{E,\ell}$ is irreducible iff it is absolutely irreducible.
\item $\bar\rho_{E,\ell}$ is reducible iff $E$ has a rational isogeny of degree $\ell$; for $\ell=3$, iff the $3$-division polynomial $\psi_3$ has a root
in $\Q$.
\item The semisimplification $\bar\rho^{\mathrm{ss}}$ is determined up to isomorphism by the classes $a_p(E)\bmod\ell$ for $p$ in any set of primes of
density one not dividing $N\ell$; if $\bar\rho$ is irreducible, $\bar\rho$ itself is. If $g$ is a newform of weight $2$ with coefficients in the ring of
integers $\mathcal O$ of a number field, $\lambda\mid\ell$ a prime of $\mathcal O$ with residue field $k$, and $a_p(g)\equiv a_p(E)\pmod\lambda$ for a set of
primes of density one, then $\bar\rho_{g,\lambda}\simeq\bar\rho_{E,\ell}\otimes_{\Fl}k$ over a common finite extension of $k$ whenever $\bar\rho_{E,\ell}$ is irreducible.
\end{enumerate}
\end{proposition}

\begin{proof}
(i) Suppose $\bar\rho$ is irreducible and $L\subset E[\ell]\otimes\Fbar_\ell$ is a line stable under $G_\Q$. Choose the basis of
\Cref{prop:galrep:rho-basic}(iii), in which $\bar\rho(c)=\operatorname{diag}(1,-1)$. The line $L$ is stable under $\bar\rho(c)$, hence an eigenline of it;
the eigenvalues $1\ne-1$ are distinct, so the eigenlines are $\Fbar_\ell e_+$ and $\Fbar_\ell e_-$, and $L=\Fbar_\ell e_\pm$ for one sign. Then for every $g$,
$\bar\rho(g)e_\pm\in L\cap E[\ell]=\Fl e_\pm$, so $\Fl e_\pm$ is a stable line of $E[\ell]$: a contradiction. The converse is trivial. (This is the
criterion ``distinct eigenvalues in the base field'' of \Cref{ch:reps}.)

(ii) A $G_\Q$-stable line of $E[\ell]$ is a $G_\Q$-stable subgroup $C$ of order $\ell$, and these are exactly the kernels of isogenies $E\to E/C$
defined over $\Q$ of degree $\ell$ (\Cref{ch:ellcurves}). For $\ell=3$, a line is $\{O,P,-P\}$ and is stable iff $g(P)=\pm P$ for all $g$, iff the
$x$-coordinate $x(P)$, which is the same for $P$ and $-P$ and different for different lines, is rational. The four roots of $\psi_3$ are the
$x$-coordinates of the four lines (\Cref{ch:ellcurves}).

(iii) For $p\nmid N\ell$ the characteristic polynomial of $\bar\rho(\Frob_p)$ is $X^2-a_pX+p\bmod\ell$ (\Cref{prop:galrep:good} below). By
Chebotarev's theorem every element of the finite group $\bar\rho(G_\Q)$ is $\bar\rho(\Frob_p)$ for a set of primes of positive density, which meets every
set of density one (the complement of a set of density one has density $0$, so it cannot contain a set of positive density); so the characteristic polynomials of all elements of the image are determined by the $a_p\bmod\ell$ on the given set,
and by the Brauer-Nesbitt theorem so is the semisimplification (\Cref{ch:galois}, \Cref{ch:reps}). If $\bar\rho$ is irreducible it is
semisimple and equal to $\bar\rho^{\mathrm{ss}}$. For the newform, $\bar\rho_{g,\lambda}$ (\Cref{ch:modgalrep}) has characteristic polynomial $X^2-a_p(g)X+p$ at
$\Frob_p$ for $p$ prime to the level and $\ell$, so the same argument over $k$ gives $\bar\rho_{g,\lambda}^{\mathrm{ss}}\simeq(\bar\rho_{E,\ell}\otimes k)^{\mathrm{ss}}$
after extension of scalars to a common finite field; when $\bar\rho_{E,\ell}$ is irreducible it is absolutely irreducible by (i), so
$\bar\rho_{E,\ell}\otimes k$ is irreducible and semisimple, and $\bar\rho_{g,\lambda}$, having it as its semisimplification and the same dimension, is isomorphic
to it (a representation whose semisimplification is irreducible is irreducible). The statement over $\Fl$ alone needs the absolute
irreducibility: \Cref{ch:galois} gives the details.
\end{proof}

\begin{paperlink}
Lemma lem:old, the case $p\nmid\ell N$: ``Consequently $\bar\rho_h\simeq\bar\rho$ over a common finite extension of $\F_\ell$, by Chebotarev and
Brauer-Nesbitt, $\bar\rho$ being irreducible, hence absolutely irreducible.'' This is \Cref{prop:galrep:irreducibility}(iii) with (i), the paper's
newform $h$ being the book's $g$. The same
reading of ``irreducible'' is used in Lemma lem:h0 (Schur's lemma for $\bar\rho$) and in Proposition prop:Sf; for $\bar\rho_{E,\ell}$, $\ell$ odd, the two
notions agree, but for the restriction of $\bar\rho$ to $\Gal(\Qbar/\Q(\sqrt{\ell^*}))$ they do not (\Cref{prop:galrep:ns}).
\end{paperlink}

\begin{example}\label{ex:galrep:eleven}
For \curve{11a1}, $y^2+y=x^3-x^2-10x-20$, the $3$-division polynomial is $\psi_3=3x^4-4x^3-60x^2-237x-21$, irreducible over $\Q$, with discriminant
$-3^3\cdot11^{10}$ and Galois group $S_4$; so $E[3]$ is irreducible and, by \Cref{thm:galrep:mod3-images}, $\bar\rho_{\curve{11a1},3}$ is surjective. At $\ell=5$
the label is \texttt{5Cs.1.1}: $E[5]\simeq\Z/5\oplus\mu_5$ (the rational point $(5,5)$ of order $5$ spans the first summand), so $E[5]$ is reducible
and \curve{11a1} has two rational $5$-isogenies, to \curve{11a2} and \curve{11a3}. All of this was recomputed with PARI \cite{PARI}.

Here is the work. The coefficients are $a_1=0$, $a_2=-1$, $a_3=1$, $a_4=-10$, $a_6=-20$, so $b_2=a_1^2+4a_2=-4$, $b_4=2a_4+a_1a_3=-20$, $b_6=a_3^2+4a_6=-79$ and
$b_8=a_1^2a_6+4a_2a_6-a_1a_3a_4+a_2a_3^2-a_4^2=80-1-100=-21$, and $\psi_3=3x^4+b_2x^3+3b_4x^2+3b_6x+b_8$ is the quartic above; the discriminant of
the curve is $-11^5$. To see the Galois group, factor $\psi_3$ modulo small primes of good reduction not dividing $3\disc(\psi_3)$: modulo $5$ it is
irreducible, so Frobenius at $5$ acts on the roots as a $4$-cycle; modulo $7$ it is a linear factor times an irreducible cubic, so Frobenius at $7$
acts as a $3$-cycle (\Cref{ch:galois}). A subgroup of $S_4$ containing a $4$-cycle and a $3$-cycle has order divisible by $12$ and contains an odd
permutation, so it is $S_4$; in particular $\psi_3$ is irreducible, and no line of $E[3]$ is rational. For the point $P=(5,5)$: it lies on the
curve ($5^2+5=30=5^3-5^2-10\cdot5-20$), and PARI gives $5P=O$ with $P\ne O$, so $P$ has order $5$.
\end{example}

\section[N\'eron-Ogg-Shafarevich and inertia invariants]{N\'eron-Ogg-Shafarevich, component groups and inertia invariants}\label{sec:galrep:nos}

This section relates the action of inertia on the torsion of $E$ to the reduction of $E$. The criterion of N\'eron, Ogg and
Shafarevich says when inertia acts trivially; the component-group lemma computes the invariants of inertia on $E[\ell^k]$ at every bad
prime; the theorem of Serre and Tate describes the image of inertia at a prime of potentially good reduction. The paper uses all
three: $\delta_p$ in Section 1.1, ``potentially good'' and the finiteness of $\rho_{E,3}(I_q)$ in Lemma lem:IV, the unramifiedness of $\bar\rho_h$ in Lemma lem:old, the
sentence ``$N(\bar\sigma)$ is prime to every prime at which $\bar\sigma$ is unramified'' in the optimal-level paragraph, and the remark before
Lemma lem:IV (``At an additive prime $q\ne\ell$ the inertia invariants of the torsion are those of the component group'').

\begin{theorem}[N\'eron, Ogg, Shafarevich]\label{thm:galrep:nos}
Let $K/\Q_p$ be a finite extension with residue field $k$, $E/K$ an elliptic curve and $\ell\ne p$ a prime. The following are equivalent:
\begin{enumerate}[label=(\alph*)]
\item $E$ has good reduction over $K$;
\item $I_K$ acts trivially on $T_\ell E$;
\item $I_K$ acts trivially on $E[m]$ for every integer $m\ge1$ prime to $p$;
\item $I_K$ acts trivially on $E[m]$ for infinitely many integers $m$ prime to $p$.
\end{enumerate}
Moreover:
\begin{enumerate}[label=(\alph*),start=5]
\item $E$ has potentially good reduction iff $\rho_{E,\ell}(I_K)$ is finite, for one, equivalently for every, $\ell\ne p$;
\item $E$ has potentially multiplicative reduction ($v_K(j)<0$) iff $\rho_{E,\ell}(I_K)$ is infinite, and then $\rho_{E,\ell}(I_K)$ contains a unipotent
element $\ne1$.
\end{enumerate}
For $E/\Q$ and $p\ne\ell$: $\rho_{E,\ell}$ is unramified at $p$ iff $p\nmid N$.
\end{theorem}

\begin{proof}
(a)$\Rightarrow$(c). Let $m$ be prime to $p$, $P\in E[m]$, and let $L/K$ be a finite Galois extension containing the coordinates of all points of
$E[m]$. A minimal model of $E$ over $K$ with good reduction still has good reduction over $L$. Let $\sigma\in I_K$. Since $\sigma$ acts trivially on
the residue field of $L$, $\sigma(P)$ and $P$ have the same reduction, so $\sigma(P)-P$ lies in $E_1(L)$, the kernel of reduction, and in $E[m]$. But
$E_1(L)$ has no torsion of order prime to $p$ (\Cref{ch:reduction}: it is isomorphic to the formal group, on which multiplication by $m$ is
bijective). So $\sigma(P)=P$.

(c)$\Rightarrow$(b) and (b)$\Rightarrow$(d). $T_\ell E=\varprojlim E[\ell^k]$, so $I_K$ acts trivially on $T_\ell E$ iff it acts trivially on every $E[\ell^k]$; and
$m=\ell^k$, $k\ge1$, are infinitely many integers prime to $p$.

(d)$\Rightarrow$(a). Suppose $E$ has bad reduction, and let $m$ be prime to $p$ with $I_K$ acting trivially on $E[m]$. The points of $E[m]$ then have
coordinates in the maximal unramified extension $K^{\ur}$, hence in a finite unramified extension $K'$ of $K$, with finite residue field $k'$.
The reduction type, the identity component and the component group do not change under unramified extension (\Cref{ch:reduction}), so
$E$ has bad reduction over $K'$ and there are exact sequences
\[
0\to E_1(K')\to E_0(K')\to\tilde E_{\mathrm{ns}}(k')\to0,\qquad 0\to E_0(K')\to E(K')\to\Phi_E(k')\to0,
\]
where $\Phi_E(k')\subseteq\Phi_E(\bar k)$ is a subgroup of the geometric component group, a finite group whose order does not depend on $K'$
(\Cref{ch:reduction}). The group $E_1(K')$ has no $m$-torsion, so $E_0(K')[m]$ injects into $\tilde E_{\mathrm{ns}}(k')[m]$. Here $\tilde E_{\mathrm{ns}}(k')$ is
$k'^+$ (additive reduction), which has no $m$-torsion as $p\nmid m$, or $k'^\times$ or the kernel of the norm from the quadratic extension of $k'$
(multiplicative reduction), subgroups of $\bar k^\times$, whose $m$-torsion is cyclic of order at most $m$ (it consists of roots of $X^m-1$, which has at
most $m$ roots in a field, and a finite subgroup of the multiplicative group of a field is cyclic). Hence $\#E_0(K')[m]\le m$ and
\[
m^2=\#E[m]=\#E(K')[m]\le\#E_0(K')[m]\cdot\#\Phi_E(k')\le m\cdot\#\Phi_E(\bar k) .
\]
So $m\le\#\Phi_E(\bar k)$: only finitely many $m$ are possible, contradicting (d).

(e), (f): see the sketch below. For $E/\Q$, the conductor exponent $f_p$ is the Artin conductor $\dim V-\dim V^{I_p}+\operatorname{sw}_p(V)$ of $V=V_\ell E$
(its definition in \Cref{ch:reduction}), a sum of two nonnegative terms; so $f_p=0$ iff $V^{I_p}=V$ (the Swan term then vanishes), iff $\rho_{E,\ell}$ is
unramified at $p$, iff $E$ has good reduction at $p$ by (a)$\Leftrightarrow$(b) over $K=\Q_p$.
\end{proof}

\begin{proof}[Sketch of proof of (e) and (f)]
By the semistable reduction theorem (\Cref{ch:reduction}) there is a finite extension $L/K$ over which $E$ has good or multiplicative reduction,
and which is good iff $v_K(j)\ge0$. If it is good, (a)$\Rightarrow$(b) over $L$ shows that $I_L=I_K\cap G_L$, a subgroup of finite index of $I_K$, acts
trivially on $T_\ell E$, so $\rho(I_K)$ is finite. If it is multiplicative, \Cref{thm:galrep:tate-shape} over $L$ (after a further unramified quadratic
extension, the curve is a Tate curve) shows that $\rho(I_L)$ consists of the matrices $\sm{1}{n\,t_\ell(\sigma)}{0}{1}$ with $n=-v_L(j)>0$ and $t_\ell$ onto
$\Zl$, an infinite group of unipotent elements. The two cases are exclusive, which gives (e) and (f); the condition on $j$ does not involve $\ell$.
The references are Serre and Tate \cite[Section 1 and Section 2]{SerreTate68} and \cite[VII.7]{SilvermanAEC}.
\end{proof}

The following lemma is proved in \Cref{ch:reduction}; we restate it because every inertia computation of this chapter rests on it.

\begin{lemma}[inertia invariants and component groups]\label{lem:galrep:component-invariants}
Let $q\ne\ell$ be primes, $E/\Q_q$ with bad reduction, $\Phig_q=\Phi_q(\Fbar_q)$ and $k\ge1$.
\begin{enumerate}[label=(\roman*)]
\item If the reduction is additive, there is a $\Gal(\Fbar_q/\F_q)$-equivariant isomorphism $E[\ell^k]^{I_q}\simeq\Phig_q[\ell^k]$; in particular
$E(\Q_q)[\ell^\infty]\simeq\Phi_q(\F_q)[\ell^\infty]$.
\item If the reduction is multiplicative, there is an exact sequence of $\Gal(\Fbar_q/\F_q)$-modules
$0\to\mu_{\ell^k}\otimes\varepsilon\to E[\ell^k]^{I_q}\to\Phig_q[\ell^k]\to0$, where $\Phig_q\simeq\Z/n_q$ and $\varepsilon$ is the unramified character of
\Cref{thm:galrep:tate-shape}.
\item Consequently, for $\ell$ odd, $\dim_{\Fl}E[\ell]^{I_q}$ is $1+[\ell\mid n_q]$ for type $\kod{I}_{n}$ ($n\ge1$), $1$ for types $\kod{IV}$, $\kod{IV}^*$ when
$\ell=3$, and $0$ for every other additive type.
\end{enumerate}
\end{lemma}

\begin{proof}[Recollection of the proof]
This is proved in \Cref{ch:reduction}. In one paragraph: $E[\ell^k]^{I_q}=E(\Q_q^{\ur})[\ell^k]$, and $\Q_q^{\ur}$ is the union of the finite unramified
extensions $K'$. For each, $E(K')/E_0(K')\simeq\Phi_q(k')$, $E_1(K')$ is uniquely $\ell$-divisible, and $E_0(K')/E_1(K')\simeq\tilde E_{\mathrm{ns}}(k')$ is $k'^+$
(additive: uniquely $\ell$-divisible) or a torus (multiplicative: $\ell$-divisible after enlarging $k'$, with $\ell^k$-torsion $\mu_{\ell^k}$ twisted by the
action of Frobenius on the tangent directions at the node). The snake lemma for multiplication by $\ell^k$ and the union over $K'$ give (i)
and (ii). Part (iii) uses the orders of $\Phig_q$: $n$ for $\kod{I}_n$, $1,2,3,4,4,3,2,1$ for $\kod{II},\kod{III},\kod{IV},\kod{I}_0^*,\kod{I}_n^*,\kod{IV}^*,
\kod{III}^*,\kod{II}^*$ (\Cref{ch:reduction}); for $\ell$ odd the only additive types with nonzero $\Phig_q[\ell]$ are $\kod{IV}$ and $\kod{IV}^*$ at $\ell=3$, with
$\Phig_q\simeq\Z/3$.
\end{proof}

\begin{warning}\label{warn:galrep:component-invariants}
The lemma distinguishes additive from multiplicative reduction, and it concerns $\Phig_q$, not the Tamagawa number $c_q=\#\Phi_q(\F_q)$. For
good reduction $E[\ell^k]^{I_q}=E[\ell^k]$ (\Cref{thm:galrep:nos}). The Frobenius-invariant form $E(\Q_q)[\ell^\infty]\simeq\Phi_q(\F_q)[\ell^\infty]$ holds only for
additive reduction: at a split multiplicative prime with $q\equiv1\pmod\ell$, $E(\Q_q)$ contains $\mu_\ell$ although $\ell$ may not divide $c_q$. The lemma
holds at every $q\ne\ell$, including $q=2,3$; this is what makes it useful at the wild primes, where the Kodaira type does not determine
the image of inertia.
\end{warning}

\begin{theorem}[Serre and Tate; Kraus]\label{thm:galrep:serre-tate}
Let $q$ be a prime, $E/\Q_q$ with potentially good reduction, and $\ell\ne q$.
\begin{enumerate}[label=(\roman*)]
\item The kernel of $\rho_{E,\ell}|_{I_q}$ is $\Gal(\Qbar_q/L)$, where $L=\Q_q^{\ur}(E[m])$ for any $m\ge3$ prime to $q$ is the smallest extension of $\Q_q^{\ur}$
over which $E$ has good reduction. So $\rho_{E,\ell}(I_q)\simeq\Gal(L/\Q_q^{\ur})$ does not depend on $\ell$; moreover the characteristic polynomial of
$\rho_{E,\ell}(\sigma)$, $\sigma\in I_q$, has coefficients in $\Z$ and does not depend on $\ell$.
\item $\rho_{E,\ell}(I_q)$ is a finite subgroup of $\SL_2(\Zl)$. For $\ell$ odd, reduction modulo $\ell$ is injective on it. A finite subgroup of $\SL_2(\Ql)$
with a nonzero fixed vector is trivial; so $V^{I_q}=0$ unless the reduction is good.
\item $\rho_{E,\ell}(I_q)$ is cyclic of order $1$, $2$, $3$, $4$ or $6$ if $q\ge5$; one of these or $\Z/3\rtimes\Z/4$ (order $12$) if $q=3$; one of these,
$Q_8$ or $\SL_2(\F_3)$ if $q=2$.
\item If $q\ge5$, $\rho_{E,\ell}(I_q)$ has order prime to $q$, factors through the tame quotient of $I_q$, is cyclic, and has order
$12/\gcd(v_q(\Delta_{\min}),12)$, that is, $1,6,4,3,2,3,4,6$ for $\kod{I}_0,\kod{II},\kod{III},\kod{IV},\kod{I}_0^*,\kod{IV}^*,\kod{III}^*,\kod{II}^*$.
\end{enumerate}
\end{theorem}

\begin{proof}
(i) is the theorem of Serre and Tate \cite[Section 2]{SerreTate68}, sketched in \Cref{ch:reduction}; we do not prove it here. (iii) is Kraus's
classification \cite{Kraus90}, stated in \Cref{ch:reduction}.

(ii) Finiteness is (i) (or \Cref{thm:galrep:nos}(e)). The determinant is $\chi_\ell$, which is trivial on $I_q$ because $\mu_{\ell^\infty}$ lies in the
unramified extension $\Q_q^{\ur}$ ($q\ne\ell$); so the image lies in $\SL_2(\Zl)$. The kernel of $\GL_2(\Zl)\to\GL_2(\Fl)$ is $1+\ell M_2(\Zl)$, which for $\ell$ odd is
torsion-free (\Cref{ch:groups}); a finite subgroup meets it trivially, so reduction is injective on it. (The reason, in brief: write $g\ne1$ in the kernel as
$g=1+\ell^rA$ with $r\ge1$ and $A\not\equiv0\pmod\ell$. For $m$ prime to $\ell$, $g^m\equiv1+m\ell^rA\not\equiv1\pmod{\ell^{r+1}}$; and the binomial expansion gives
$g^\ell\equiv1+\ell^{r+1}A\pmod{\ell^{r+2}}$, because $\binom\ell k\ell^{kr}$ is divisible by $\ell^{2r+1}$ for $2\le k<\ell$ and $\ell^{r\ell}$ by $\ell^{r+2}$ when $\ell\ge3$. Iterating, $g^{\ell^j}=1+\ell^{r+j}A_j$ with
$A_j\equiv A\not\equiv0\pmod\ell$, and then $g^{\ell^jm}\ne1$ for $m$ prime to $\ell$: no power of $g$ is $1$.
For $\ell=2$ this fails: $-1=1+2\cdot(-1)$ has order $2$, which is why the injectivity is stated for $\ell$ odd.) Let $G\subseteq\SL_2(\Ql)$ be finite and $v\ne0$
with $gv=v$ for all $g\in G$. Each $g\in G$ has finite order, hence is diagonalisable over $\Qbar_\ell$ (its minimal polynomial divides $X^r-1$,
which has distinct roots in characteristic $0$); one eigenvalue is $1$ and the determinant is $1$, so both eigenvalues are $1$ and $g=1$. Hence
$V^{I_q}\ne0$ forces $\rho(I_q)=1$, which is good reduction by \Cref{thm:galrep:nos}.

(iv) The order of $\rho(I_q)$ divides $24$ by (iii), hence is prime to $q\ge5$. The wild inertia group $P_q$ is a pro-$q$ group, so its image is a
$q$-group of order prime to $q$, hence trivial, and $\rho|_{I_q}$ factors through the tame quotient $I_q/P_q$. A finite quotient of the tame
inertia group is cyclic (\Cref{ch:local}). The order is Kraus's semistability defect $12/\gcd(v_q(\Delta_{\min}),12)$ (\Cref{ch:reduction}, where
it is deduced from (i)); the values by type follow from $v_q(\Delta_{\min})=0,2,3,4,6,8,9,10$ for the potentially good types at $q\ge5$
(\Cref{ch:reduction}).
\end{proof}

\begin{remark}[independence of $\ell$]\label{rem:galrep:wd-independence}
\Cref{thm:galrep:serre-tate}(i), with \Cref{prop:galrep:good}, \Cref{thm:galrep:tate-shape} and \Cref{prop:galrep:potmult} below, shows that the
Frobenius-semisimplified Weil-Deligne representation of $V_\ell E|_{G_{\Q_q}}$ ($q\ne\ell$) is defined over $\Q$ and does not depend on $\ell$: at a good
prime it is unramified with Frobenius polynomial $X^2-a_qX+q$; at a potentially multiplicative prime it is $\mathrm{sp}(2)$ twisted by an
$\ell$-independent quadratic character; at a potentially good prime, inertia acts through the $\ell$-independent group $\Gal(L/\Q_q^{\ur})$ with
integral characteristic polynomials, and Frobenius acts through its action on the good reduction of $E$ over $L$ (\cite[Section 2]{SerreTate68}).
This is the statement quoted in \Cref{ch:local}; we only sketch it, as the paper needs nothing beyond the explicit shapes proved below.
\end{remark}

\section[Good and multiplicative reduction]{Good and multiplicative reduction; Lemma lem:tamagawa-exp}\label{sec:galrep:mult}

This section describes $\rho|_{G_{\Q_p}}$ at the primes $p\ne\ell$ of good and of potentially multiplicative reduction, and proves the paper's
Lemma lem:tamagawa-exp. The matrices of \Cref{thm:galrep:tate-shape} are the ones the paper writes in the proof of Lemma lem:tamfactor
(\Cref{ch:blochkato}) and quotes from \cite[Theorem 3.1(e)]{DDT} in Lemma lem:old; the potentially multiplicative additive case, including
$q=2$, is what the proof of Theorem thm:main needs at the primes of type $\kod{I}_n^*$.

\begin{proposition}[good reduction]\label{prop:galrep:good}
Let $p\nmid N\ell$ and $a_p=p+1-\#\tilde E(\F_p)$. Then $\rho$ is unramified at $p$; the arithmetic Frobenius $\Frob_p$ has characteristic polynomial
$X^2-a_pX+p$ on $V$, with roots $\alpha$, $\beta$; the geometric Frobenius $\Frob_p^{-1}$ has roots $\alpha^{-1}$, $\beta^{-1}$, characteristic polynomial
$X^2-\frac{a_p}{p}X+\frac1p$; on $\ad V$ both act with eigenvalues $\alpha/\beta,1,\beta/\alpha$; $\delta_p=2$. The characteristic polynomial of $\bar\rho(\Frob_p)$
is $X^2-a_pX+p\bmod\ell$, so $\tr\bar\rho(\Frob_p)\equiv a_p$ and $\det\bar\rho(\Frob_p)\equiv p\pmod\ell$. For $\ell=3$ and $p\nmid3\disc(\psi_3)$, the orbits of $\bar\rho(\Frob_p)$
on the four lines of $E[3]$ have the lengths of the degrees of the irreducible factors of $\psi_3\bmod p$.
\end{proposition}

\begin{proof}
Unramified: \Cref{thm:galrep:nos}. Reduction gives a $G_{\Q_p}$-equivariant isomorphism $T_\ell E\simeq T_\ell\tilde E$ (\Cref{thm:galrep:nos}, proof of
(a)$\Rightarrow$(c), and \Cref{ch:reduction}): reduction $E[\ell^k]\to\tilde E[\ell^k]$ is injective because its kernel lies in the kernel of reduction $E_1$, which has no
$\ell$-torsion, and it is then bijective because both groups have $\ell^{2k}$ elements; it commutes with $G_{\Q_p}$, which acts on $\tilde E(\Fbar_p)$ through
$\Gal(\Fbar_p/\F_p)$. On $T_\ell\tilde E$ the element $\Frob_p$ acts as the Frobenius endomorphism $\pi_p$ of $\tilde E/\F_p$ (both raise coordinates to the $p$-th power), and the
characteristic polynomial of $\pi_p$ on $T_\ell\tilde E$ is $X^2-a_pX+p$ (\Cref{ch:ellcurves}: $\deg(1-\pi_p)=\#\tilde E(\F_p)$ and $\deg\pi_p=p$). The inverse has
the inverse eigenvalues. On $\ad V$, if $\rho(\Frob_p)$ has eigenvectors $v_\alpha$, $v_\beta$ (over $\Qbar_\ell$; when $\alpha=\beta$ the eigenvalues on
$\ad V$ are $1,1,1$ in any case), conjugation by it multiplies $v_\alpha\otimes v_\beta^\vee$, $v_\alpha\otimes v_\alpha^\vee-v_\beta\otimes v_\beta^\vee$, $v_\beta\otimes v_\alpha^\vee$
by $\alpha/\beta$, $1$, $\beta/\alpha$; the inverse gives the same set. $\delta_p=2$ as $I_p$ acts trivially. Reduction modulo $\ell$ of the characteristic
polynomial gives that of $\bar\rho(\Frob_p)$. The last statement is \Cref{ch:galois} (Frobenius acting on the roots of a polynomial whose
discriminant is prime to $p$), applied to the four roots of $\psi_3$, which correspond to the four lines (\Cref{prop:galrep:irreducibility}(ii)).
\end{proof}

\begin{theorem}[the Tate shape]\label{thm:galrep:tate-shape}
Let $q$ be a prime of multiplicative reduction, $\ell\ne q$, $n_q=v_q(\Delta_{\min})=-v_q(j)$, $\varepsilon\colon G_{\Q_q}\to\{\pm1\}$ the unramified character with
$\varepsilon(\Frob_q)=a_q(E)$ ($+1$ for split, $-1$ for nonsplit reduction), $q_E\in\Q_q^\times$ the Tate parameter, and $\kappa\colon G_{\Q_q}\to\Zl$ the Kummer cocycle
of $q_E$ (after fixing a generator of $\Zl(1)$), $g(q_E^{1/\ell^k})=\zeta_{\ell^k}^{\kappa(g)}q_E^{1/\ell^k}$. There is a $\Zl$-basis $e_1,e_2$ of $T$ in which
\[
\rho(g)=\begin{pmatrix}\chi_\ell(g)\varepsilon(g)&\varepsilon(g)\kappa(g)\\0&\varepsilon(g)\end{pmatrix}\qquad(g\in G_{\Q_q}).
\]
Consequently:
\begin{enumerate}[label=(\roman*)]
\item on $I_q$, $\rho(\sigma)=\sm{1}{n_q\,t_\ell(\sigma)}{0}{1}$, where $t_\ell\colon I_q\to\Zl$ is the $\ell$-adic tame character, which is onto; $\rho$ factors on $I_q$
through its maximal pro-$\ell$ quotient, and $\rho(I_q)$ is infinite;
\item $V^{I_q}=\Ql e_1$, so $\delta_q=1$, and $\Frob_q$ acts on it by $a_q(E)\,q$;
\item $\Ql e_1$ is the unique $G_{\Q_q}$-stable line of $V$, and $V$ is a non-split extension of $\varepsilon$ by $\varepsilon\chi_\ell$;
\item the Weil-Deligne representation of $V|_{G_{\Q_q}}$ is $\mathrm{sp}(2)\otimes\mu_\varepsilon$, of conductor exponent $1$;
\item $(\ad V)^{I_q}=\Ql e$, $e=e_1\otimes e_2^\vee$, on which $\Frob_q$ acts by $q$ and the geometric Frobenius by $q^{-1}$.
\end{enumerate}
\end{theorem}

\begin{proof}
By \Cref{ch:reduction}, $E$ over $\Q_q$ is the twist by $\varepsilon$ of the Tate curve $E_{q_E}$ (trivial twist when the reduction is split), and on
$T_\ell E_{q_E}$, with $e_1=(\zeta_{\ell^k})_k$ and $e_2=(q_E^{1/\ell^k})_k$ (compatible systems in $\Qbar_q^\times/q_E^{\Z}$), $g$ acts by $\sm{\chi_\ell(g)}{\kappa(g)}{0}{1}$: indeed
$g\zeta_{\ell^k}=\zeta_{\ell^k}^{\chi_\ell(g)}$ and $g(q_E^{1/\ell^k})=\zeta_{\ell^k}^{\kappa(g)}q_E^{1/\ell^k}$. Twisting multiplies every matrix by the scalar $\varepsilon(g)$, which gives the
displayed shape; the corner is $\varepsilon\kappa$, as the cocycle relation $c(gh)=\chi_\ell\varepsilon(g)c(h)+c(g)\varepsilon(h)$ requires. (The convention is that the
columns of $\rho(g)$ are the images of the basis vectors: $ge_1=\chi_\ell\varepsilon(g)e_1$ and $ge_2=\varepsilon\kappa(g)e_1+\varepsilon(g)e_2$. The Kummer cocycle satisfies
$\kappa(gh)=\kappa(g)+\chi_\ell(g)\kappa(h)$, since $gh(q_E^{1/\ell^k})=g(\zeta_{\ell^k}^{\kappa(h)}q_E^{1/\ell^k})=\zeta_{\ell^k}^{\chi_\ell(g)\kappa(h)+\kappa(g)}q_E^{1/\ell^k}$; multiplying by
$\varepsilon(gh)=\varepsilon(g)\varepsilon(h)$ gives exactly the upper right entry $\chi_\ell\varepsilon(g)\cdot\varepsilon\kappa(h)+\varepsilon\kappa(g)\cdot\varepsilon(h)$ of the product
$\rho(g)\rho(h)$, so the displayed formula is a homomorphism.) The equalities $\varepsilon(\Frob_q)=a_q(E)$ and
$n_q=v_q(q_E)=-v_q(j)$ are in \Cref{ch:reduction}.

(i) On $I_q$, $\varepsilon=1$ and $\chi_\ell=1$ ($q\ne\ell$). Write $q_E=q^{n_q}u$ with $u\in\Z_q^\times$; the Kummer cocycle is additive in its argument,
$\kappa=n_q\kappa_q+\kappa_u$. The $\ell^k$-th roots of the unit $u$ lie in $\Q_q^{\ur}$ (Hensel's lemma, as $q\ne\ell$), so $\kappa_u$ vanishes on $I_q$; and $\kappa_q$
restricted to $I_q$ is the tame character $t_\ell$, which is onto $\Zl$ (\Cref{ch:local}). Since $n_q\ne0$, the image is infinite, and since $t_\ell$
factors through the maximal pro-$\ell$ quotient of $I_q$, so does $\rho|_{I_q}$.

(ii) A vector $xe_1+ye_2$ is fixed by $\sm{1}{s}{0}{1}$ iff $sy=0$; some $s=n_qt_\ell(\sigma)$ is nonzero, so $V^{I_q}=\Ql e_1$; $\Frob_q$ acts on $e_1$ by
$\chi_\ell\varepsilon(\Frob_q)=q\,a_q(E)$.

(iii) A $G_{\Q_q}$-stable line is stable under $\rho(\sigma)=\sm{1}{s}{0}{1}$ with $s\ne0$, whose only eigenline is $\Ql e_1$. So $\Ql e_1$ is the only stable
line, and $V$ is not a direct sum of two characters: the extension $0\to\varepsilon\chi_\ell\to V\to\varepsilon\to0$ is non-split.

(iv) The Weil-Deligne representation attached to this shape is computed in \Cref{ch:local}: the monodromy operator is $N=\sm0100$ (from the
logarithm of the unipotent action of $t_\ell$), Frobenius acts through $\varepsilon\chi_\ell$ and $\varepsilon$, and the conductor exponent is
$\dim V-\dim V^{I_q}=1$.

(v) In the basis $e,h,f$ of $\ad V$ (\Cref{conv:galrep:notation}), conjugation by $\sm1s01$ sends $e\mapsto e$, $h\mapsto h-2se$, $f\mapsto f+sh-s^2e$. So
$ae+bh+cf$ is fixed iff $(-2bs-cs^2)e+cs\,h=0$, and for $s\ne0$ this means $c=0$ and then $b=0$ (in $\Ql$, $2$ is invertible):
$(\ad V)^{I_q}=\Ql e$. For an upper triangular $g=\sm{a}{b}{0}{d}$ one has $geg^{-1}=\frac ad\,e$; for a Frobenius lift this is $\chi_\ell\varepsilon/\varepsilon=\chi_\ell(\Frob_q)=q$.
The geometric Frobenius acts by $q^{-1}$.
\end{proof}

\begin{paperlink}
Lemma lem:old, the cases $p\,\|\,N$: ``$\rho_{E,\ell}|_{D_p}$ is an extension of an unramified character $\varepsilon$ by $\varepsilon\chi_\ell$ with
$\varepsilon(\Frob_p)=a_p(f)=\pm1$ \cite[Theorem 3.1(e)]{DDT}''; and ``A non-split extension of two characters has a unique stable line''. These are
\Cref{thm:galrep:tate-shape} and its part (iii) (and, modulo $\ell$, \Cref{lem:galrep:tamagawa-exp}(iii)). In the proof of Lemma lem:tamfactor
the paper uses the matrix of the arithmetic Frobenius, with corner $\varepsilon\kappa$; the factor $1-q^{-1-s}$ of the table of Section 1 comes from
the geometric Frobenius acting by $q^{-1}$ on $(\ad V)^{I_q}$, (v). The two conventions are each right in their place; the paper's
Section 1.1 says so (``The matrices in the proof of Lemma~lem:tamfactor are those of the arithmetic Frobenius'').
\end{paperlink}

\begin{lemma}[the paper's Lemma lem:tamagawa-exp]\label{lem:galrep:tamagawa-exp}
Let $q$ be a prime of multiplicative reduction, $n_q=v_q(\Delta_E)$, and $\ell\ne q$ a prime. Then $\rho_{E,\ell}\bmod\ell^t$ is unramified at $q$ if and only
if $\ell^t\mid n_q$. This holds for split and for nonsplit reduction. Moreover:
\begin{enumerate}[label=(\roman*)]
\item the Tamagawa exponent $t_q=\max\{t\ge0:\rho\bmod\ell^t\text{ is unramified at }q\}$ equals $v_\ell(n_q)=v_\ell(-v_q(j))=t_\ell(E,q)$;
\item $\#E[\ell^t]^{I_q}=\ell^{\,t+\min(t,v_\ell(n_q))}$;
\item if $\ell\nmid n_q$, then $\bar\rho|_{D_q}$ is ramified, a non-split extension of $\varepsilon$ by $\varepsilon\bar\chi_\ell$ with a unique stable line, and
$\bar\rho(I_q)$ has order $\ell$;
\item if $\ell\mid n_q$, then $\bar\rho|_{D_q}$ is unramified and $\bar\rho(\Frob_q)$ has eigenvalues $a_q(E)q$ and $a_q(E)$ modulo $\ell$.
\end{enumerate}
\end{lemma}

\begin{proof}[First proof (the paper's, expanded)]
Let $K=\Q_q$ if the reduction is split and the unramified quadratic extension of $\Q_q$ if not. Then $I_K=I_q$ (an unramified extension does not change
the inertia group), and over $K$ the curve is isomorphic to the Tate curve $E_{q_E}$ (\Cref{ch:reduction}: over the unramified quadratic extension
the reduction becomes split, and $\varepsilon$ becomes trivial). By Tate's theorem $E_{q_E}(\Qbar_q)\simeq\Qbar_q^\times/q_E^{\Z}$ Galois-equivariantly, and
$E[\ell^t]$ is the subgroup of classes $\zeta^a\alpha^b$, where $\zeta$ is a primitive $\ell^t$-th root of unity, $\alpha$ an $\ell^t$-th root of $q_E$, and
$a,b\in\Z/\ell^t$. Since $q\ne\ell$, $\mu_{\ell^t}\subset\Q_q^{\ur}$, so $I_q$ fixes $\zeta$. For $\sigma\in I_q$, $\sigma(\alpha)/\alpha$ is an $\ell^t$-th root of unity; the class
of $\alpha$ is fixed iff $\sigma(\alpha)/\alpha\in q_E^{\Z}$, and as $q_E^{\Z}\cap\mu_{\ell^t}=\{1\}$ ($v_q(q_E)>0$), iff $\sigma(\alpha)=\alpha$. So $I_q$ acts trivially on $E[\ell^t]$ iff
$\alpha\in\Q_q^{\ur}$ for some, equivalently every, choice of $\alpha$, that is (Kummer theory over $\Q_q^{\ur}$, which contains $\mu_{\ell^t}$), iff $q_E$ is an
$\ell^t$-th power in $\Q_q^{\ur}$. Write $q_E=q^{n_q}u$ with $u$ a unit. The unit $u$ is an $\ell^t$-th power in $\Q_q^{\ur}$: the polynomial $X^{\ell^t}-u$ is
separable modulo $q$ and has a root modulo $q$ in $\Fbar_q$, which lifts by Hensel's lemma to a finite unramified extension (\Cref{ch:local}).
So $q_E$ is an $\ell^t$-th power in $\Q_q^{\ur}$ iff $q^{n_q}$ is; and $q$ is a uniformiser of $\Q_q^{\ur}$, whose valuation group is $\Z$, so $q^{n_q}=y^{\ell^t}$ forces
$\ell^t\mid n_q$, while conversely $\ell^t\mid n_q$ gives $q^{n_q}=(q^{n_q/\ell^t})^{\ell^t}$. Hence $\rho\bmod\ell^t$ is unramified iff $\ell^t\mid n_q$. Nothing depended on
the reduction being split.
\end{proof}

\begin{proof}[Second proof (component groups)]
By \Cref{lem:galrep:component-invariants}(ii), $\#E[\ell^t]^{I_q}=\#\mu_{\ell^t}\cdot\#\Phig_q[\ell^t]=\ell^t\cdot\#(\Z/n_q)[\ell^t]=\ell^{t+\min(t,v_\ell(n_q))}$, which is (ii). So
$E[\ell^t]^{I_q}=E[\ell^t]$, of order $\ell^{2t}$, iff $\min(t,v_\ell(n_q))=t$, iff $\ell^t\mid n_q$. (A third, one-line proof: by \Cref{thm:galrep:tate-shape}(i),
$\rho\bmod\ell^t$ on $I_q$ is $\sm{1}{n_qt_\ell}{0}{1}\bmod\ell^t$, trivial iff $\ell^t\mid n_q$ since $t_\ell$ is onto $\Zl$.)
\end{proof}

\begin{proof}[Proof of the supplements]
(i) is the main statement, with $n_q=-v_q(j)$ (\Cref{thm:galrep:tate-shape}) and the definition of $t_\ell(E,q)$. (ii) is the second proof (or the
matrices: $(x,y)\in(\Z/\ell^t)^2$ is fixed by all $\sm{1}{n_qs}{0}{1}$, $s\in\Zl$, iff $n_qy\equiv0\pmod{\ell^t}$, which leaves $\ell^t$ choices of $x$ and
$\ell^{\min(t,v_\ell(n_q))}$ of $y$). (iii) If $\ell\nmid n_q$, $\bar\rho(\sigma)=\sm{1}{n_qt_\ell(\sigma)}{0}{1}\bmod\ell$ runs through the whole unipotent group of order
$\ell$; the extension of the unramified $\varepsilon$ by $\varepsilon\bar\chi_\ell$ (unramified too) is ramified, hence non-split, and as in
\Cref{thm:galrep:tate-shape}(iii) the line $\Fl e_1$ is the only stable line. (iv) If $\ell\mid n_q$, $\bar\rho$ is trivial on $I_q$ and
$\bar\rho(\Frob_q)=\sm{q\varepsilon(\Frob_q)}{*}{0}{\varepsilon(\Frob_q)}\bmod\ell$ with $\varepsilon(\Frob_q)=a_q(E)$.
\end{proof}

\begin{paperlink}
The paper's proof (quoted from Lemma lem:tamagawa-exp): ``By Kummer theory over $\Q_q^{\mathrm{ur}}$, which contains $\mu_{\ell^t}$ as $q\ne\ell$,
inertia acts trivially on $E[\ell^t]$ if and only if $q_E$ is an $\ell^t$-th power in $\Q_q^{\mathrm{ur}}$; the units of its valuation ring are $\ell^t$-th
powers by Hensel's lemma, so this holds if and only if $\ell^t\mid n_q$.'' The first proof above separates the two steps this sentence
compresses: Kummer theory reduces the question to $q_E$ being an $\ell^t$-th power, and Hensel's lemma reduces it to the valuation. Section 1.5
uses part (i): ``for multiplicative $q$ one has $t_q=v_\ell(n_q)=t_\ell(E,q)$ whether or not the reduction is split''. Lemma lem:old uses (iii)
and (iv) in its two cases $p\,\|\,N$, and the optimal-level paragraph uses the consequence that $N(\bar\rho)$ is prime to $q$ iff $\ell\mid n_q$
(\Cref{prop:galrep:serre-conductor-E}).
\end{paperlink}

A worked case. The curve \curve{14a1}, $y^2+xy+y=x^3+4x-6$, has $\Delta_{\min}=-2^6\cdot7^3$, so it has type $\kod{I}_6$ at $2$ and $\kod{I}_3$ at $7$, with $a_2=-1$ (nonsplit) and
$a_7=1$ (split). Take $\ell=3$. Since $3\mid6$ and $3\mid3$, \Cref{lem:galrep:tamagawa-exp} says that $\bar\rho_{E,3}$ is unramified at $2$ and at $7$, while $\rho\bmod9$ is
ramified at both ($v_3(6)=v_3(3)=1$, so $t_2=t_7=1$), and $\#E[3^t]^{I_q}=3^{t+1}$ for $q=2,7$ and every $t\ge1$. By (iv), $\bar\rho(\Frob_2)$ has eigenvalues
$a_2\cdot2=-2\equiv1$ and $a_2=-1\equiv2\pmod3$, and $\bar\rho(\Frob_7)$ has eigenvalues $7\equiv1$ and $1$. All of this agrees with the global structure
$E[3]\simeq\Z/3\oplus\mu_3$ of \curve{14a1} (image \texttt{3Cs.1.1}, \Cref{thm:galrep:mod3-images}): this module is unramified at every prime $q\ne3$, and $\Frob_q$
acts on it with eigenvalues $1$ and $q\bmod3$. (Recomputed with PARI.) The reader can do the same for \curve{11a1} and \curve{11a3} at $\ell=5$ in
\Cref{ex:galrep:eleven-five}.

The next proposition treats the primes at which $E$ has additive reduction but $v_q(j)<0$. They look additive, but after a ramified quadratic twist they
become multiplicative, so everything just proved applies to them up to a sign.

\begin{proposition}[potentially multiplicative additive reduction]\label{prop:galrep:potmult}
Let $q\ne\ell$ ($q=2$ allowed), $\ell$ odd, and let $E$ have additive reduction at $q$ with $v_q(j)<0$; put $\nu=-v_q(j)$. Then there is a ramified
quadratic character $\chi$ of $G_{\Q_q}$ such that the twist of $E$ by $\chi$ has multiplicative reduction, and $\rho|_{G_{\Q_q}}\simeq\rho'\otimes\chi$ with
$\rho'$ of the shape of \Cref{thm:galrep:tate-shape} for a Tate parameter of valuation $\nu$. For $q$ odd one may take $\chi=\chi_{q^*}|_{G_{\Q_q}}$; then
$E^{(q)}$ has type $\kod{I}_\nu$ at $q$ and $E$ has type $\kod{I}_\nu^*$. Moreover:
\begin{enumerate}[label=(\roman*)]
\item $\rho(I_q)=\bigl\{\chi(\sigma)\sm{1}{\nu t_\ell(\sigma)}{0}{1}\bigr\}$ is infinite and contains $-1$; $V^{I_q}=0$, so $\delta_q=0$;
\item $\ad V\simeq\ad V'$ as $G_{\Q_q}$-modules, so $(\ad V)^{I_q}$ is a line on which $\Frob_q$ acts by $q$;
\item $\rho\bmod\ell^t$ is ramified at $q$ for every $t\ge1$, while $(\rho\otimes\chi)\bmod\ell^t$ is unramified at $q$ iff $\ell^t\mid\nu$;
\item the conductor exponent is $f_q=2a(\chi)$, where $a(\chi)$ is the conductor exponent of $\chi$: $f_q=2$ for $q$ odd, and $f_2\in\{4,6\}$.
\end{enumerate}
At $q=2$ the index $m$ of the Kodaira symbol $\kod{I}_m^*$ is not $\nu$ (\Cref{ch:reduction} gives $m=\nu+2f_2-4$), and a symbol $\kod{I}_m^*$ at $2$ can also
occur with potentially good reduction.
\end{proposition}

\begin{proof}
By \Cref{ch:reduction}, $E$ over $\Q_q$ is the quadratic twist of the Tate curve $E_{q_E}$ by $\gamma=-c_4/c_6$, and the reduction is additive iff
$\Q_q(\sqrt\gamma)/\Q_q$ is ramified; take $\chi$ the character of this extension, and $\rho'=\rho_{E_{q_E},\ell}$, of Tate shape with $\varepsilon=1$ and
$v_q(q_E)=\nu$. The twist of $E$ by $\chi$ is $E_{q_E}$, which is split multiplicative. For $q$ odd every ramified quadratic character of $G_{\Q_q}$ is
$\chi_{q^*}$ times an unramified one (\Cref{ch:local}), and an unramified quadratic twist keeps the Tate shape (with $\varepsilon$ the unramified
factor); the global twist $E^{(q)}$ has the local twist as its completion at $q$, of type $\kod{I}_\nu$, and $E$ has type $\kod{I}_\nu^*$ by the
exchange of types (\Cref{ch:reduction}).

(i) The formula follows from \Cref{thm:galrep:tate-shape}(i). The pair $(t_\ell,\chi)\colon I_q\to\Zl\times\{\pm1\}$ is onto: its image is a closed subgroup
projecting onto both factors, and since $\Zl$ is pro-$\ell$ with $\ell$ odd while $\{\pm1\}$ has order $2$, such a subgroup is the whole product. (Explicitly: let
$X$ be the image, and pick $(x,-1)\in X$. Its $\ell^k$-th power is $(\ell^kx,-1)$, as $\ell^k$ is odd; as $k\to\infty$ this tends to $(0,-1)$, which lies in $X$ because $X$
is closed. For any $y\in\Zl$ some $(y,\pm1)$ lies in $X$, and multiplying by $(0,-1)$ if necessary gives both $(y,1)$ and $(y,-1)$.) So
some $\sigma$ has $t_\ell(\sigma)=0$, $\chi(\sigma)=-1$, and $\rho(\sigma)=-1$, which fixes no nonzero vector: $V^{I_q}=0$. (ii) Conjugation by $\chi(g)\rho'(g)$ equals
conjugation by $\rho'(g)$; then \Cref{thm:galrep:tate-shape}(v). (iii) $-1\in\rho(I_q)$ and $-1\not\equiv1\pmod{\ell^t}$ for $\ell$ odd; the second statement is
\Cref{lem:galrep:tamagawa-exp} for $E_{q_E}$ (or its unramified twist), whose $\nu$ plays the role of $n_q$. (iv) The wild inertia group acts
through $\chi$ only ($\rho'$ is tame on $I_q$, factoring through the pro-$\ell$ tame quotient), so on the wild inertia group $V$ is the sum of two copies of
the character $\chi$. The Swan conductor is additive and depends only on the wild inertia group, and $a(\chi)=1+\operatorname{sw}(\chi)$ for a ramified character
(\Cref{ch:local}); so $\operatorname{sw}(V)=2\operatorname{sw}(\chi)=2(a(\chi)-1)$ and
$f_q=\dim V-\dim V^{I_q}+\operatorname{sw}(V)=2+2(a(\chi)-1)=2a(\chi)$. For $q$ odd $\chi$ is tame, $a(\chi)=1$; for $q=2$ the ramified quadratic characters of
$\Q_2$ have conductor exponent $2$ or $3$ (\Cref{ch:local}). The last sentence is \Cref{ch:reduction}.
\end{proof}

\begin{paperlink}
The proof of Theorem thm:main bounds $v_\ell(\Tam_p(\ad T))$ below by $t_\ell(E,p)$ ``at every potentially multiplicative prime $p$,
multiplicative or additive, $p=2$ included, by Lemma~lem:tamfactor''. At an additive such prime, Lemma lem:tamfactor is applied to the
Tate curve of which $E$ is a twist; this is legitimate because $\ad T$ is unchanged by the twist, \Cref{prop:galrep:potmult}(ii) (integrally:
conjugation by $\chi\rho'$ equals conjugation by $\rho'$ on $\ad T$). The argument is local and needs no global twist, which is why it covers
$q=2$, where the global twist of Lemma lem:twist is not available. Lemma lem:euler's row $\kod{I}_n^*$ (``$V=V'\otimes\chi_{q^*}$ with $V'$ Steinberg
at $q$'') is \Cref{prop:galrep:potmult} for $q$ odd.
\end{paperlink}

\begin{corollary}\label{cor:galrep:delta}
For $p\ne\ell$, $\delta_p=\dim V^{I_p}$ is $2$, $1$, $0$ according as $E$ has good, multiplicative or additive reduction at $p$, and the conductor
exponent is $f_p=2-\delta_p+\operatorname{sw}_p(V)$.
\end{corollary}

\begin{proof}
Good: \Cref{prop:galrep:good}. Multiplicative: \Cref{thm:galrep:tate-shape}(ii). Additive: \Cref{thm:galrep:serre-tate}(ii) if the reduction is
potentially good (the image of inertia is then nontrivial, by \Cref{thm:galrep:nos}), \Cref{prop:galrep:potmult}(i) if not. The conductor
exponent of \Cref{ch:reduction} is defined as the Artin conductor $a(V_\ell E)=\dim V-\dim V^{I_p}+\operatorname{sw}_p(V)$; that it agrees with Ogg's formula
is the theorem of Ogg and Saito, stated there.
\end{proof}

\section[Additive potentially good reduction]{Additive potentially good reduction: the image of inertia by Kodaira type}\label{sec:galrep:additive}

This section supplies the Kodaira-type column of the shapes table of \Cref{ch:local}: for each additive potentially good type it gives
the image of inertia on $T$, its invariants on $\ad V$ and the action of Frobenius on them, and the reduction modulo $3$ with its invariants
and Serre exponent. It proves the paper's Lemma lem:IV (at every prime $q\ne3$) and Remark rem:types, and gives the inertia half of Lemma
lem:euler; the $L$-factors and the terms $e_\ell(E,q)$ are then computed in \Cref{ch:lfunctions}. The type $\kod{IV}$ at $2$ is covered as well
(\Cref{cor:galrep:IV-at-2}).

\begin{theorem}[inertia at additive potentially good primes, $q\ge5$]\label{thm:galrep:additive-shapes}
Let $q\ge5$, $\ell\ne q$, and let $E$ have additive potentially good reduction at $q$. Then $\rho(I_q)=\langle g\rangle$ is cyclic, $g\in\SL_2(\Zl)$, and, with
$\sigma$ denoting an element with $\sigma^2+\sigma+1=0$ and $\tau$ one with $\tau^2=-1$:
\begin{center}\small\setlength{\tabcolsep}{4pt}
\begin{tabular}{llllll}
\toprule
type & $v_q(\Delta_{\min})$ & $\#\rho(I_q)$ & generator & $(\ad V)^{I_q}$ & $\Frob_q$ acts on it by\\
\midrule
$\kod{I}_0^*$ & $6$ & $2$ & $-1$ & $\ad V$ & eigenvalues $\alpha'/\beta',1,\beta'/\alpha'$\\
$\kod{IV}$, $\kod{IV}^*$ & $4$, $8$ & $3$ & $\sigma$ & $\Ql(2\sigma+1)$ & $\chi_{-3}(q)$\\
$\kod{III}$, $\kod{III}^*$ & $3$, $9$ & $4$ & $\tau$ & $\Ql\tau$ & $\chi_{-4}(q)$\\
$\kod{II}$, $\kod{II}^*$ & $2$, $10$ & $6$ & $-\sigma$ & $\Ql(2\sigma+1)$ & $\chi_{-3}(q)$\\
\bottomrule
\end{tabular}
\end{center}
Here $\alpha',\beta'$ are the roots of $X^2-a_q(E^{(q)})X+q$. In every case $V^{I_q}=0$ and $f_q=2$; for every Frobenius lift,
$\rho(\Frob_q)\,g\,\rho(\Frob_q)^{-1}=g^q$; and the action on the line is the same for the arithmetic and the geometric Frobenius. Moreover $(2\sigma+1)^2=-3$; $\End_{I_q}(T)=\Zl[g]\simeq\Zl\otimes\Z[\zeta_3]$
(resp.\ $\Zl\otimes\Z[i]$) for orders $3$, $6$ (resp.\ $4$), which is not a domain when $\ell\equiv1\pmod3$ (resp.\ $\pmod4$); for $\kod{I}_0^*$,
$\rho\otimes\chi_{q^*}$ is unramified at $q$ and equals $\rho_{E^{(q)},\ell}$ on $G_{\Q_q}$, $E^{(q)}$ having good reduction at $q$. For $\kod{III}$, $\kod{III}^*$ the line is
spanned by $\tau$ itself, which is diagonalisable over $\Ql$ only when $\ell\equiv1\pmod4$.
\end{theorem}

\begin{proof}
By \Cref{thm:galrep:serre-tate}(iv) the image is cyclic of the stated order $m\in\{2,3,4,6\}$ ($m\ne1$ as the reduction is not good) and lies in
$\SL_2(\Zl)$. Let $g$ be a generator. It is diagonalisable over $\Qbar_\ell$ with eigenvalues $\lambda,\lambda^{-1}$, roots of unity of the same order as
each other, and the order of $g$ is the order of $\lambda$. So: for $m=2$, $\lambda=-1$ and $g=-1$; for $m=3$, $\lambda$ is a primitive cube root of
unity and $g$ has characteristic polynomial $X^2+X+1$, so $g=\sigma$ satisfies $\sigma^2+\sigma+1=0$ (Cayley-Hamilton); for $m=4$, $\lambda=\pm i$ and $g=\tau$
has characteristic polynomial $X^2+1$; for $m=6$, $\lambda$ is a primitive sixth root of unity, $-\lambda$ a primitive cube root, and $-g=\sigma$
satisfies $\sigma^2+\sigma+1=0$. This proves the generator column; (\Cref{ch:reps} lists the same shapes.)

Invariants. For $m=2$ every endomorphism commutes with $-1$, so $(\ad V)^{I_q}=\ad V$. For $m\in\{3,4,6\}$, $g$ has distinct eigenvalues, so its
commutant in $\End(V)$ is $\Ql[g]=\Ql\oplus\Ql g$, of dimension $2$ (its dimension over $\Qbar_\ell$ is $2$, and $\Ql[g]$ is contained in it and
two-dimensional). The trace-zero part is the line spanned by $g-g^{-1}$, whose trace is $(\lambda+\lambda^{-1})-(\lambda^{-1}+\lambda)=0$, and which is nonzero. For
$g=\sigma$: $\sigma^{-1}=\sigma^2=-\sigma-1$, so $g-g^{-1}=2\sigma+1$, and $(2\sigma+1)^2=4\sigma^2+4\sigma+1=-3$; for $g=-\sigma$, $g-g^{-1}=-(2\sigma+1)$, the same line;
for $g=\tau$, $g-g^{-1}=2\tau$, the line $\Ql\tau$, and $\tau$ is diagonalisable over $\Ql$ iff its eigenvalues $\pm i$ lie in $\Ql$, iff $\ell\equiv1\pmod4$.

Frobenius. The image of inertia has order prime to $q$, so it is a quotient of the tame inertia group, and $g=\rho(t)$ for a topological
generator $t$ of tame inertia. The tame relation $\Frob\,t\,\Frob^{-1}=t^q$ (\Cref{ch:local}) gives $\rho(\Frob_q)g\rho(\Frob_q)^{-1}=g^q$ for every Frobenius lift.
Then conjugation by $\rho(\Frob_q)$ sends $g-g^{-1}$ to $g^q-g^{-q}$. For $m=3$: $g^q=g$ if $q\equiv1\pmod3$ and $g^q=g^2=g^{-1}$ if $q\equiv2\pmod3$, so the line is
multiplied by $\chi_{-3}(q)$; for $m=6$: $q\equiv\pm1\pmod6$, and $q\equiv1\pmod6$ iff $q\equiv1\pmod3$, again $\chi_{-3}(q)$; for $m=4$: $g^q=g^{\pm1}$ according as
$q\equiv\pm1\pmod4$, that is, $\chi_{-4}(q)$. A sign $c=\pm1$ is its own inverse, so the geometric Frobenius acts by the same sign. (Here the choice of
Frobenius could have mattered: for the geometric Frobenius the tame relation reads $\Frob^{-1}t\,\Frob=t^{q'}$ with $qq'=1$ in the relevant quotient, so it raises $g$
to the power $q^{-1}\bmod m$. For $m\in\{3,4,6\}$ and $q$ prime to $m$ one has $q^2\equiv1\pmod m$, so $q^{-1}\equiv q\pmod m$ and the two Frobenii conjugate $g$ to the same
power. At a multiplicative prime there is no such coincidence, \Cref{thm:galrep:tate-shape}(v).) For $m=2$: the
restriction of $\rho$ to $I_q$ is the quadratic character $\sigma\mapsto\pm1$ of $I_q$; since $q$ is odd, $I_q$ has a unique continuous character of order
$2$ (it factors through the tame quotient $\prod_{r\ne q}\Z_r(1)$, and only the factor $\Z_2(1)$ has a quotient of order $2$), namely
$\chi_{q^*}|_{I_q}$. So $\rho\otimes\chi_{q^*}$ is unramified at $q$; it is $\rho_{E^{(q)},\ell}$ (\Cref{prop:galrep:twist-isogeny}), and $E^{(q)}$ has good
reduction at $q$ by \Cref{thm:galrep:nos}. Then $\ad V=\ad V_\ell E^{(q)}$, unramified with the eigenvalues of \Cref{prop:galrep:good}.

The rest. $V^{I_q}=0$ by \Cref{thm:galrep:serre-tate}(ii); the wild inertia acts trivially, so $\operatorname{sw}=0$ and $f_q=2$. The lattice $T$ is free of rank
one over $\Zl[g]\simeq\Zl[X]/(\Phi_m(X))$ for $m\in\{3,4,6\}$ (\Cref{ch:ant}, where this is proved for every odd $\ell$; for $\ell=2$ it follows in the
same way when $\Phi_m$ stays irreducible modulo $2$, and we need only $\ell$ odd), so $\End_{I_q}(T)=\End_{\Zl[g]}(T)=\Zl[g]$, and
$\Zl[X]/(\Phi_m)\simeq\Zl\otimes\Z[\zeta_m]$. If $\ell\equiv1\pmod m$, $\Phi_m$ has a root in $\Zl$ by Hensel's lemma and $\Zl[X]/(\Phi_m)\simeq\Zl\times\Zl$ is not a domain.
\end{proof}

\begin{paperlink}
Lemma lem:euler, types $\kod{II},\kod{II}^*,\kod{IV},\kod{IV}^*$: ``On $\ad V$ the sign disappears and the invariants form the line spanned by
$2\sigma+1$, whose square is $-3$, inside $\operatorname{End}_{I_q}(T)=\Z_\ell[\sigma]\simeq\Z_\ell\otimes\Z[\zeta_3]$ (not a field when
$\ell\equiv1\pmod3$)''; types $\kod{III}$, $\kod{III}^*$: ``the invariant line of $\ad V$ is spanned by $\tau$ itself, which has trace $0$ (it is
diagonalisable only over $\Q_\ell(i)$)''. These are \Cref{thm:galrep:additive-shapes}. The paper's phrase ``for either choice of Frobenius'' is the
statement that a sign is its own inverse; at a multiplicative prime the two Frobenii act by $q$ and $q^{-1}$ (\Cref{thm:galrep:tate-shape}(v)), and
only the geometric one gives the factor $1-q^{-1-s}$ of the table of Section 1. The $\kod{I}_0^*$ row holds at $q=3$ as well, where the image of
inertia is $\{\pm1\}$ (\Cref{ch:reduction}).
\end{paperlink}

\begin{remark}[Kraus's table from component groups]\label{rem:galrep:kraus-from-components}
At $q\ge5$ the dictionary between types and orders in \Cref{thm:galrep:additive-shapes} can also be derived from the component groups alone,
without Kraus's formula, given that the order $m$ lies in $\{2,3,4,6\}$ and does not depend on $\ell$ (\Cref{thm:galrep:serre-tate}(i), (iv)). For
a generator $g$ of order $m$, the fixed spaces of its reductions on $E[2]$ and $E[3]$ have dimensions $(2,0)$, $(0,1)$, $(1,0)$, $(0,0)$ for
$m=2,3,4,6$: modulo $2$, $-1$ is trivial, $\sigma$ has the irreducible characteristic polynomial $X^2+X+1$, and $\tau$ is a nontrivial unipotent (if
$\tau=1+2A$ with $A\in M_2(\Z_2)$, then $\tau^2=-1$ gives $A+A^2=-\frac12$, impossible); modulo $3$, $-1$ and $-\sigma$ have only the eigenvalue $-1$,
$\sigma$ is a nontrivial unipotent, and $\tau$ has the irreducible characteristic polynomial $X^2+1$. By \Cref{lem:galrep:component-invariants} at
$\ell=2$ and $\ell=3$ these dimensions are those of $\Phig_q[2]$ and $\Phig_q[3]$, which are $(2,0)$ for $\kod{I}_0^*$, $(0,1)$ for $\kod{IV}$, $\kod{IV}^*$,
$(1,0)$ for $\kod{III}$, $\kod{III}^*$ and $(0,0)$ for $\kod{II}$, $\kod{II}^*$. Since the four pairs are distinct, the type determines $m$:
$\kod{I}_0^*\mapsto2$, $\kod{IV},\kod{IV}^*\mapsto3$, $\kod{III},\kod{III}^*\mapsto4$, $\kod{II},\kod{II}^*\mapsto6$. The details are \Cref{ex:galrep:kraus}.
\end{remark}

\begin{lemma}[the paper's Lemma lem:IV]\label{lem:galrep:IV}
Let $q\ne3$ be a prime at which $E$ has reduction of type $\kod{IV}$ or $\kod{IV}^*$. Then:
\begin{enumerate}[label=(\roman*)]
\item the reduction at $q$ is potentially good, and the image of $I_q$ in $\Aut(T_\ell E)$ is cyclic of order $3$ for every prime $\ell\ne q$;
\item its image in $\Aut(E[3])$ is a nontrivial unipotent subgroup; $E[3]^{I_q}$ is a line; $\bar\rho_{E,3}$ is tamely ramified at $q$; the Serre
conductor of $\bar\rho_{E,3}$ has exponent $1$ at $q$; and the conductor of $E$ has exponent $f_q=2$;
\item if $\sigma$ generates $\rho_{E,3}(I_q)$, then $T_3E$ is free of rank one over $\Z_3[X]/(X^2+X+1)\simeq\Z_3[\zeta_3]$ with $X$ acting as $\sigma$, and
$E[3]\simeq\Z_3[\zeta_3]/(\pi^2)$ with $\pi=1-\zeta_3$, $\sigma-1$ acting as multiplication by $-\pi$;
\item for $\ell\ne q$ and $\sigma$ a generator of $\rho_{E,\ell}(I_q)$, $(\ad V_\ell E)^{I_q}=\Ql(2\sigma+1)$ is a line on which $\Frob_q$ and $\Frob_q^{-1}$ act by $\chi_{-3}(q)$; so, with the conventions of
\Cref{conv:galrep:notation}(c), $L_q(A_f,s)^{-1}=1-\chi_{-3}(q)q^{-s}$, $P_q(X)=1-\chi_{-3}(q)qX$ and $e_\ell(E,q)=v_\ell(q-\chi_{-3}(q))$.
\end{enumerate}
\end{lemma}

\begin{proof}[Proof (the paper's, every step expanded)]
(i) and (ii). By \Cref{lem:galrep:component-invariants}(i) (the reduction is additive, $q\ne3$), $E[3]^{I_q}\simeq\Phig_q[3]$, and $\Phig_q\simeq\Z/3$ for the types
$\kod{IV}$, $\kod{IV}^*$ (\Cref{ch:reduction}); so $E[3]^{I_q}$ is a line. Inertia acts nontrivially on $E[3]$: otherwise $E[3]^{I_q}=E[3]$ would have order $9$, not $3$.

The reduction is potentially good. Otherwise, by \Cref{prop:galrep:potmult}, $\rho_{E,3}(I_q)$ consists of the matrices
$\chi(\sigma)\sm{1}{\nu t_3(\sigma)}{0}{1}$ with $\chi$ ramified quadratic, and some $\sigma$ acts by $-1$ (proof of \Cref{prop:galrep:potmult}(i)), which fixes no line of $E[3]$: a contradiction. (The paper phrases the same
argument as: the subgroup of $I_q$ on which the corner vanishes modulo $3$ has index dividing $3$, so it is not contained in the kernel of $\chi$,
of index $2$. Alternatively, types $\kod{IV}$, $\kod{IV}^*$ have $v_q(j)\ge0$, \Cref{ch:reduction}.)

So $G=\rho_{E,3}(I_q)$ is finite (\Cref{thm:galrep:nos}(e)), lies in $\SL_2(\Z_3)$, and reduction modulo $3$ is injective on it
(\Cref{thm:galrep:serre-tate}(ii)). Its image $\bar G\subseteq\SL_2(\F_3)$ is nontrivial and fixes a nonzero vector $v$. In a basis $(v,v')$ of $E[3]$ the elements
of $\bar G$ are $\sm{1}{b}{0}{d}$ with $d=\det=1$, so $\bar G$ lies in the unipotent group $\mathrm U=\{\sm1b01\}$ of order $3$, and equals it (\Cref{ch:groups}:
the elements of $\SL_2(\F_3)$ fixing a nonzero vector form a group of order $3$). Hence $G\simeq\bar G$ is cyclic of order $3$, its image in $\Aut(E[3])$ is a
nontrivial unipotent group, and by \Cref{thm:galrep:serre-tate}(i) the image of $I_q$ in $\Aut(T_\ell E)$ is cyclic of order $3$ for every $\ell\ne q$. The order
$3$ is prime to $q$ (as $q\ne3$), so the image is a quotient of the tame inertia group: the ramification is tame, $\operatorname{sw}=0$. The Serre exponent of
$\bar\rho_{E,3}$ at $q$ is $\dim E[3]-\dim E[3]^{I_q}+0=2-1=1$, and the conductor exponent of $E$ is $\dim V-\dim V^{I_q}+0=2-0=2$
(\Cref{thm:galrep:serre-tate}(ii)).

(iii) $\sigma$ has determinant $\chi_3^{\mathrm{cyc}}(\sigma)=1$ and order $3$, so its characteristic polynomial is $X^2+X+1$ (proof of
\Cref{thm:galrep:additive-shapes}) and $\sigma^2+\sigma+1=0$ on $T=T_3E$. So $T$ is a module over $\Z_3[X]/(X^2+X+1)$ with $X$ acting as $\sigma$, and this ring is
$\Z_3[\zeta_3]$, a discrete valuation ring with uniformiser $\pi=1-\zeta_3$ and residue field $\F_3$, $\pi^2=-3\zeta_3$ (\Cref{ch:ant}). The module $T$ is
torsion-free of $\Z_3$-rank $2$, hence a torsion-free finitely generated module of rank one over the DVR $\Z_3[\zeta_3]$, hence free of rank one
(\Cref{ch:ant}). Then $E[3]=T/3T\simeq\Z_3[\zeta_3]/(3)=\Z_3[\zeta_3]/(\pi^2)$, on which $\sigma-1$ acts as $\zeta_3-1=-\pi$: nonzero, with square zero, a nontrivial
unipotent whose fixed line is $\pi\Z_3[\zeta_3]/(\pi^2)$.

(iv) This is the computation of \Cref{thm:galrep:additive-shapes} for $m=3$, which used only that $\rho(I_q)$ is cyclic of order $3$, prime to $q$: the
commutant of $\sigma$ meets $\ad V$ in $\Ql(2\sigma+1)$, and $\rho(\Frob_q)\sigma\rho(\Frob_q)^{-1}=\sigma^q$ by the tame relation, which is $\sigma^{\chi_{-3}(q)}$. So $\Frob_q^{-1}$
acts on $(\ad V)^{I_q}$ by $\chi_{-3}(q)$, $L_q(A_f,s)^{-1}=1-\chi_{-3}(q)q^{-s}$, $P_q(X)=L_q(A_f,s-1)^{-1}|_{X=q^{-s}}=1-\chi_{-3}(q)qX$, and, the naive factor being $1$ at an
additive prime, $e_\ell(E,q)=v_\ell(1-\chi_{-3}(q)q^{-1})=v_\ell(q-\chi_{-3}(q))$ since $v_\ell(q)=0$.
\end{proof}

\begin{proof}[Second proof of (i) and (ii) for $q\ge5$]
By Kraus's theorem (\Cref{thm:galrep:serre-tate}(iv)) the image of $I_q$ is cyclic of order $12/\gcd(v_q(\Delta_{\min}),12)=3$, as $v_q(\Delta_{\min})\in\{4,8\}$; it is
tame because $3\ne q$. A generator $\sigma$ of $\rho_{E,3}(I_q)$ gives, as in (iii), $E[3]\simeq\Z_3[\zeta_3]/(\pi^2)$ with $\sigma-1$ acting as $-\pi$: a nontrivial unipotent
with a fixed line. The Serre exponent is $2-1=1$.
\end{proof}

\begin{paperlink}
The paper states Lemma lem:IV for every $q\ne3$, $q=2$ included. What makes $q=2$ accessible is the component-group argument,
which \Cref{ch:reduction} also gives for the inertia statement at $2$; Kraus's formula, used in the second proof above, is available only for $q\ge5$. At $q=3$ the types $\kod{IV}$, $\kod{IV}^*$ are
wildly ramified (\Cref{ch:reduction}), and the lemma does not apply: an image of order $3$ at $q=3$ is not tame. Theorem thm:qual (a) uses
part (ii): ``By Lemma~lem:IV the Serre conductor $M=N(\bar\rho)$ has exponent exactly $1$ at $q$, so $M\mid N/q$.''
\end{paperlink}

A worked case of \Cref{lem:galrep:IV}. The curve \curve{121c1}, $y^2+xy=x^3+x^2-2x-7$, has $\Delta_{\min}=-11^4$ and type $\kod{IV}$ at $q=11$. Kraus's formula gives an
image of inertia of order $12/\gcd(4,12)=3$, as the lemma predicts. Since $11\equiv2\pmod3$, $\chi_{-3}(11)=-1$: Frobenius inverts $\sigma$ and acts on $(\ad V)^{I_{11}}$ by
$-1$, so $L_{11}(A_f,s)^{-1}=1+11^{-s}$, $P_{11}(X)=1+11X$ (PARI's \texttt{lfuneuler(lfunsympow(E,2),11)} returns $1/(1+11X)$), and $e_\ell(E,11)=v_\ell(12)$, which is
$1$ for $\ell=3$ and $0$ for $\ell\ge5$. The Serre exponent of $\bar\rho_{E,3}$ at $11$ is $2-1=1$, one less than $f_{11}=2$; this is why
$N(\bar\rho_{\curve{121c1},3})=11$ in \Cref{ex:galrep:serre-conductor-examples}.

\begin{corollary}[types $\kod{IV}$, $\kod{IV}^*$ at $2$]\label{cor:galrep:IV-at-2}
Let $E/\Q$ have type $\kod{IV}$ or $\kod{IV}^*$ at $2$. Then $f_2(E)=2$ (so $4\,\|\,N$), $v_2(\Delta_{\min})=4$ for type $\kod{IV}$ and $8$ for type $\kod{IV}^*$, $\rho_{E,\ell}(I_2)$ is
cyclic of order $3$ for every odd $\ell$, $E[3]^{I_2}$ is a line, and the Serre conductor of $\bar\rho_{E,3}$ has exponent $1$ at $2$. Moreover $\Frob_2$ inverts
the generator $\sigma$ of $\rho(I_2)$ and acts on $(\ad V)^{I_2}=\Ql(2\sigma+1)$ by $\chi_{-3}(2)=-1$; so $L_2(A_f,s)^{-1}=1+2^{-s}$, $P_2(X)=1+2X$, $e_3(E,2)=v_3(3/2)=1$
and $e_\ell(E,2)=0$ for $\ell\ge5$.
\end{corollary}

\begin{proof}
Everything except the discriminant is \Cref{lem:galrep:IV} at $q=2$, with $\chi_{-3}(2)=-1$ ($2\equiv-1\pmod3$, so $\sigma^2=\sigma^{-1}$). Ogg's formula
$f_2=v_2(\Delta_{\min})+1-m$ (\Cref{ch:reduction}), with $m$ the number of components of the special fibre of the minimal regular model, $m=3$ for
$\kod{IV}$ and $m=7$ for $\kod{IV}^*$, gives $v_2(\Delta_{\min})=f_2-1+m=4$, resp.\ $8$.
\end{proof}

All $4{,}750$ curves of conductor below $10{,}000$ with type $\kod{IV}$ ($2{,}094$ curves) or $\kod{IV}^*$ ($2{,}656$) at $2$ were checked with PARI against
\Cref{cor:galrep:IV-at-2}: every one has $f_2=2$, $v_2(\Delta_{\min})=4$, resp.\ $8$, $v_2(j)\ge0$, and, computed from the fields of definition of the
$3$-torsion points, an image of $I_2$ on $E[3]$ of order $3$ with an inertia-fixed point. On a random sample of $60$ of them, and on \curve{20a1},
\curve{52a1}, \curve{140b1} and \curve{980i1}, \texttt{lfuneuler(lfunsympow(E,2),2)} returns $1/(1+2X)$, that is, $P_2(X)=1+2X$.

\begin{paperlink}
The proof of Proposition prop:forms treats the types $\kod{IV}$, $\kod{IV}^*$ at $2$ by Lemma lem:IV (``of type $\mathrm{IV},\mathrm{IV}^*$ and
$\ell=3$, $w_3(E,2)=1=e_3(E,2)$ by Lemma~lem:IV''), and the paragraph after the proof reads: ``On the $92{,}812$ working-set curves with type
$\mathrm{IV}$ or $\mathrm{IV}^*$ at $2$, PARI's factor at $2$ is $1+2X$ for every one of them, as Lemma~lem:IV requires.'' This is \Cref{cor:galrep:IV-at-2}. The same corollary is what lets case
(a) of Theorem thm:qual hold at $q=2$: the Serre exponent of $\bar\rho_{E,3}$ at $2$ is $1$ while $4\,\|\,N$, so the level drops at $2$ exactly as at a prime
$q\ge5$; the paper's Example, case (a) at $q=2$, is \curve{52a1} (\Cref{ex:galrep:level-lowering}).
\end{paperlink}

\begin{proposition}[inertia modulo $3$ at $q\ge5$; the paper's Remark rem:types]\label{prop:galrep:mod3-inertia}
Let $q\ge5$ and $\ell=3$. The image $\bar\rho_{E,3}(I_q)$, the fixed space, the Serre exponent $a_q(\bar\rho_{E,3})$, the conductor exponent $f_q(E)$, and the
Serre exponent after the twist by $\chi_{q^*}$ are as follows (types exchanged by the twist as in \Cref{ch:reduction}).
\begin{center}\small
\begin{tabular}{lllllll}
\toprule
type at $q$ & $\bar\rho(I_q)$ generated by & order & $\dim E[3]^{I_q}$ & $a_q(\bar\rho_{E,3})$ & $f_q(E)$ & $a_q(\bar\rho_{E^{(q)},3})$\\
\midrule
$\kod{I}_0$ & $1$ & $1$ & $2$ & $0$ & $0$ & $2$\\
$\kod{I}_n$, $3\nmid n$ & unipotent $\ne1$ & $3$ & $1$ & $1$ & $1$ & $2$\\
$\kod{I}_n$, $3\mid n$ & $1$ & $1$ & $2$ & $0$ & $1$ & $2$\\
$\kod{I}_n^*$, $3\nmid n$ & $-1$ and a unipotent $\ne1$ & $6$ & $0$ & $2$ & $2$ & $1$\\
$\kod{I}_n^*$, $3\mid n$ & $-1$ & $2$ & $0$ & $2$ & $2$ & $0$\\
$\kod{I}_0^*$ & $-1$ & $2$ & $0$ & $2$ & $2$ & $0$\\
$\kod{II}$, $\kod{II}^*$ & $-u$, $u$ unipotent $\ne1$ & $6$ & $0$ & $2$ & $2$ & $1$\\
$\kod{III}$, $\kod{III}^*$ & $\bar\tau$, $\bar\tau^2=-1$ & $4$ & $0$ & $2$ & $2$ & $2$\\
$\kod{IV}$, $\kod{IV}^*$ & unipotent $\ne1$ & $3$ & $1$ & $1$ & $2$ & $2$\\
\bottomrule
\end{tabular}
\end{center}
In particular: (a) $E[3]^{I_q}=0$ and $a_q(\bar\rho_{E,3})=2=f_q$ for every additive type other than $\kod{IV}$, $\kod{IV}^*$ (including $\kod{I}_0^*$, $\kod{III}$,
$\kod{III}^*$), and $E[3]^{I_q}=0$ for these types also at $q=2$; (b) the twist by $\chi_{q^*}$ lowers the Serre exponent of $\bar\rho_{E,3}$ at $q$ exactly for
$\kod{II}$, $\kod{II}^*$, $\kod{I}_n^*$ and $\kod{I}_0^*$, and lowers the conductor exponent of $E$ exactly for $\kod{I}_n^*$ and $\kod{I}_0^*$ (for $\kod{II}\leftrightarrow\kod{IV}^*$
both curves have $f_q=2$). For $\ell\ge5$ and $q\ge5$, $q\ne\ell$: every additive type has $E[\ell]^{I_q}=0$ and Serre exponent $2$, and $\kod{I}_n$ has exponent $0$ or
$1$ according as $\ell\mid n$ or not.
\end{proposition}

\begin{proof}
Reduce the generators of \Cref{thm:galrep:tate-shape}, \Cref{prop:galrep:potmult} and \Cref{thm:galrep:additive-shapes} modulo $3$. A unipotent
$\sm{1}{nt}{0}{1}$ with $t$ running through $\Z_3$ is trivial modulo $3$ iff $3\mid n$, and otherwise generates the group $\mathrm U$ of order $3$; for $\kod{I}_n^*$ the
image is $\{\pm1\}$ times this (the pair $(t_3,\chi)$ is onto, proof of \Cref{prop:galrep:potmult}), of order $6$ or $2$. For the potentially good types
reduction is injective (\Cref{thm:galrep:serre-tate}(ii)): $-1$ has order $2$; $\sigma\bmod3$ is a nontrivial unipotent (\Cref{lem:galrep:IV}(iii)); $-\sigma\bmod3$ is
minus a nontrivial unipotent, of order $6$; $\tau\bmod3$ has square $-1$ and characteristic polynomial $X^2+1$, irreducible over $\F_3$. The fixed spaces
follow (an element with only the eigenvalue $-1$, or with irreducible characteristic polynomial, fixes no nonzero vector), and agree with
\Cref{lem:galrep:component-invariants}(iii). At $q\ge5$ every image of inertia has order prime to $q$ (unipotents modulo $3$ have order $3$), so the Swan
term vanishes and $a_q=2-\dim E[3]^{I_q}$. The column $f_q(E)$ is \Cref{cor:galrep:delta} ($f_q=1$ multiplicative, $2$ additive at $q\ge5$). The last
column is the same computation for $E^{(q)}$, whose type is obtained by the exchanges $\kod{I}_0\leftrightarrow\kod{I}_0^*$, $\kod{I}_n\leftrightarrow\kod{I}_n^*$,
$\kod{II}\leftrightarrow\kod{IV}^*$, $\kod{III}\leftrightarrow\kod{III}^*$, $\kod{IV}\leftrightarrow\kod{II}^*$ (\Cref{ch:reduction}). Statements (a) and (b) are read off the table;
the case $q=2$ in (a) is \Cref{lem:galrep:component-invariants}, since $\Phig_2[3]=0$ for all additive types other than $\kod{IV}$, $\kod{IV}^*$. For $\ell\ge5$:
$\Phig_q[\ell]=0$ for every additive type, so $E[\ell]^{I_q}=0$; and $\kod{I}_n$ is \Cref{lem:galrep:tamagawa-exp}.
\end{proof}

All rows of the table were checked with PARI on curves of conductor below $10{,}000$ with surjective mod $3$ image: for each Kodaira type at a
prime $q\ge5$ and each residue of $n$ modulo $3$, on up to $12$ curves (all of them when fewer exist), the least common multiple of the
ramification indices of $q$ in the fields $\Q(P)$, $P\in E[3]$, equals the order in the table, and an inertia-fixed point exists exactly in the
rows with $\dim E[3]^{I_q}\ge1$.

\begin{paperlink}
The paper's Remark rem:types: ``For all these types $E[3]^{I_q}=0$ and the Serre conductor of $\bar\rho_{E,3}$ has exponent $2$ at $q$. The
twist by $\chi_{q^*}$ lowers that exponent for $\mathrm{I}_0^*$, $\mathrm{I}_n^*$, $\mathrm{II}$, $\mathrm{II}^*$; it lowers the conductor of $E$
only for $\mathrm{I}_0^*$ and $\mathrm{I}_n^*$.'' This is (a) and (b) of \Cref{prop:galrep:mod3-inertia}. The remark distinguishes the two exponents: in the case
$\kod{II}\leftrightarrow\kod{IV}^*$ the twist lowers the Serre exponent of $\bar\rho_{E,3}$ but not the conductor exponent of $E$. Among the additive rows, Theorem thm:qual treats exactly those in
which, without the twist or with the twist by $\chi_{q^*}$, the Serre exponent at $q$ is strictly smaller than the conductor exponent at $q$ of the
corresponding curve ($E$, respectively $E^{(q)}$): $\kod{IV}$, $\kod{IV}^*$ without the twist ($1<2$, case (a)); $\kod{II}$, $\kod{II}^*$ with it ($1<2$, case
(b)); and $\kod{I}_n^*$ with $3\mid n$ with it ($0<1$, case (c) at $\ell=3$). In the rows $\kod{I}_0^*$ ($0=0$ after the twist), $\kod{I}_n^*$ with $3\nmid n$ ($1=1$)
and $\kod{III}$, $\kod{III}^*$ ($2=2$ either way) the Serre exponent never drops below the conductor exponent, and no level is gained.
\end{paperlink}

\begin{remark}[the wild primes]\label{rem:galrep:wild-inertia}
At $q=2$ and $q=3$ the Kodaira type does not determine $\rho(I_q)$ (\Cref{thm:galrep:serre-tate}(iii)). At $q=2$ the image can be $Q_8$ or $\SL_2(\F_3)$
(non-abelian, so absolutely irreducible on $V$, $(\ad V)^{I_2}=0$ and $P_2=1$), cyclic of order $2$, $4$ or $6$ (with a nontrivial wild part, since the tame
quotient $\prod_{r\ne2}\Z_r(1)$ of $I_2$ has no quotient of even order), or cyclic of order $3$. The last happens exactly for the types $\kod{IV}$, $\kod{IV}^*$.
These types have it by \Cref{cor:galrep:IV-at-2}. Conversely, if $\rho(I_2)$ has order $3$, the reduction is additive and potentially good
(\Cref{thm:galrep:nos}), reduction modulo $3$ is injective on $\rho_{E,3}(I_2)$ (\Cref{thm:galrep:serre-tate}(ii)), and an element of order $3$ of
$\SL_2(\F_3)$ has characteristic polynomial $X^2+X+1=(X-1)^2$, so it is a nontrivial unipotent and fixes a line; hence $E[3]^{I_2}\ne0$, so
$\Phig_2[3]\ne0$ by \Cref{lem:galrep:component-invariants}(i), and among the additive types only $\kod{IV}$ and $\kod{IV}^*$ have a component group of order
divisible by $3$ (\Cref{ch:reduction}). At $q=3$ an image of order divisible by $3$ is wild.
What survives at every $q\ne\ell$ is \Cref{lem:galrep:component-invariants} and the conductor formula \Cref{prop:galrep:serre-conductor-E}. On samples
of $40$ curves per type at $q=2$ (conductor below $10{,}000$, surjective mod $3$ image), for the potentially good types $\kod{II}$, $\kod{II}^*$, $\kod{III}$,
$\kod{III}^*$, $\kod{I}_0^*$ and $\kod{I}_m^*$ the least common multiple of the ramification indices of $2$ in the fields $\Q(P)$, $P\in E[3]$, was $2$, $4$, $6$ or $8$,
for the potentially multiplicative symbols $\kod{I}_m^*$ it was $2$ or $6$, and in all these cases $E[3]^{I_2}=0$. The Euler factors at $2$ and $3$ are
computed in \Cref{ch:lfunctions}; Remark rem:wild itself is discussed in \Cref{ch:reduction}.
\end{remark}

\section[The representation at \texorpdfstring{$\ell$}{l}]{The representation at \texorpdfstring{$\ell$}{l} when \texorpdfstring{$E$}{E} has good reduction at \texorpdfstring{$\ell$}{l}}\label{sec:galrep:at-l}

At $p=\ell$ the restriction $\rho|_{G_{\Q_\ell}}$ is not unramified even when $E$ has good reduction, and its description needs $p$-adic Hodge
theory. The paper uses only three facts about it, all for $\ell\nmid N$: that $\bar\rho|_{D_\ell}$ comes from a finite flat group scheme (``good'' in the
optimal-level theorem), that $V$ is crystalline (through \cite{DFG04}, \Cref{ch:blochkato}), and that $a_\ell\bmod\ell$ is determined by $\bar\rho|_{D_\ell}$
(Lemma lem:old, case $p=\ell$). This section states them, sketches the elliptic-curve proofs, and draws the consequence that drives
Proposition prop:Sf: a twist by $\chi_{\ell^*}$ can preserve $\bar\rho$ only when $\ell=3$.

\begin{theorem}[good reduction at $\ell$]\label{thm:galrep:good-at-l}
Let $\ell\nmid N$ and let $\mathcal E$ be the N\'eron model of $E$ over $\Zl$, an elliptic curve over $\Zl$.
\begin{enumerate}[label=(\roman*)]
\item For every $k\ge1$, $\mathcal E[\ell^k]$ is a finite flat group scheme over $\Zl$ of order $\ell^{2k}$ whose generic fibre is $E[\ell^k]$; so $\bar\rho|_{D_\ell}$ is
\emph{finite} (flat), which is what ``good'' means in the optimal-level theorem \cite[Theorem 3.15]{DDT}.
\item $V|_{G_{\Q_\ell}}$ is crystalline with Hodge-Tate weights $0$ and $1$, and the crystalline Frobenius on $\Dcris(V^\vee)$ has characteristic polynomial
$X^2-a_\ell X+\ell$.
\item If $\bar\rho$ is irreducible, its Serre weight is $k(\bar\rho)=2$.
\end{enumerate}
\end{theorem}

\begin{proof}[Sketch of proof]
(i) $\mathcal E$ is an abelian scheme of relative dimension one over $\Zl$ (\Cref{ch:reduction}: good reduction means the minimal model is smooth,
and it is then the N\'eron model). Multiplication by $\ell^k$ on an abelian scheme is quasi-finite and proper, hence finite, and flat by the
fibrewise criterion, of degree $\ell^{2k}$ on both fibres; its kernel is the finite flat group scheme $\mathcal E[\ell^k]$, with generic fibre $E[\ell^k]$. See
\cite[Section 2]{DDT} for finite flat group schemes and the meaning of ``finite'' for $\bar\rho|_{D_\ell}$. (ii) is a theorem of Fontaine (and of
Faltings for general abelian varieties); an introduction is Berger \cite{Berger}; \Cref{ch:blochkato} uses it. We do not prove it. (iii) is
Serre's recipe for the weight \cite[Section 2]{Serre87}, in the form of Edixhoven \cite[Section 2]{Edixhoven92}; for $\bar\rho_{E,\ell}$ with $\ell\nmid N$ the
recipe gives $2$ in both the ordinary and the supersingular case of \Cref{thm:galrep:deligne-fontaine}. (For orientation only: at a
multiplicative prime $\ell$ the weight is $2$ or $\ell+1$ according as $\ell\mid v_\ell(\Delta)$ or not; the paper never needs $\ell\mid N$ in a proof.)
\end{proof}

\begin{definition}[fundamental characters]\label{def:galrep:fundamental-characters}
Let $\ell$ be a prime and $I_\ell\subset G_{\Q_\ell}$ the inertia group. For $n\ge1$ choose a root $\varpi_n$ of $X^{\ell^n-1}=\ell$ in $\Qbar_\ell$, and define
\[
\omega_n\colon I_\ell\to\F_{\ell^n}^\times,\qquad\omega_n(\sigma)=\frac{\sigma(\varpi_n)}{\varpi_n}\bmod\mathfrak m,
\]
where $\sigma(\varpi_n)/\varpi_n\in\mu_{\ell^n-1}$ is reduced modulo the maximal ideal of $\bar\Z_\ell$ (which identifies $\mu_{\ell^n-1}$ with $\F_{\ell^n}^\times$). We
write $\omega=\omega_1$ (the fundamental character of level $1$) and $\omega_2$ (level $2$).
\end{definition}

Why these characters. Any continuous character of $I_\ell$ with values in a finite field of characteristic $\ell$ has order prime to $\ell$, so it kills the
pro-$\ell$ wild inertia group and factors through the tame quotient, which is procyclic. The characters $\omega_n$ are the standard generators of the
characters of tame inertia with values in $\F_{\ell^n}^\times$: every such character is a power of $\omega_n$ (by (i) of the next proposition, $\omega_n$ maps the
tame quotient onto the cyclic group $\F_{\ell^n}^\times$). They are the language in which \Cref{thm:galrep:deligne-fontaine} describes $\bar\rho|_{I_\ell}$: level $1$ in the
ordinary case, level $2$ in the supersingular case. For example, at $\ell=3$, $\omega$ has order $2$ and is the restriction of $\chi_{-3}$ to $I_3$, while
$\omega_2\colon I_3\to\F_9^\times$ has order $8$ and $\omega_2^4=\omega$.

\begin{proposition}\label{prop:galrep:fundamental-characters}
\begin{enumerate}[label=(\roman*)]
\item $\omega_n$ is a well-defined continuous surjective homomorphism, independent of the choice of $\varpi_n$, and factors through the tame
inertia group.
\item $\omega=\bar\chi_\ell|_{I_\ell}$, of order $\ell-1$.
\item $\omega_2$ has order $\ell^2-1$ and $\omega_2^{\ell+1}=\omega$.
\item For a Frobenius lift $F\in G_{\Q_\ell}$ and $\sigma\in I_\ell$, $\omega_2(F\sigma F^{-1})=\omega_2(\sigma)^\ell$. Hence $\omega_2$ does not extend to a character of $G_{\Q_\ell}$,
and for any extension $\psi$ of $\omega_2$ to $G_{\Q_{\ell^2}}$ ($\Q_{\ell^2}$ the unramified quadratic extension) the induced representation
$\Ind_{G_{\Q_{\ell^2}}}^{G_{\Q_\ell}}\psi$ is irreducible, with restriction $\omega_2\oplus\omega_2^\ell$ to $I_\ell$.
\end{enumerate}
\end{proposition}

\begin{proof}
(i) $\mu_{\ell^n-1}$ lies in $\Q_\ell^{\ur}$, which is the fixed field of $I_\ell$; Kummer theory over $\Q_\ell^{\ur}$ (\Cref{ch:local}) shows that
$\sigma\mapsto\sigma(\varpi_n)/\varpi_n$ is a homomorphism $I_\ell\to\mu_{\ell^n-1}$, independent of the root (two roots differ by a factor in $\mu_{\ell^n-1}$, fixed by
$I_\ell$). The extension $\Q_\ell^{\ur}(\varpi_n)$ is totally and tamely ramified of degree $\ell^n-1$ ($v(\varpi_n)=1/(\ell^n-1)$), so the map is onto, and it is
trivial on wild inertia, a pro-$\ell$ group. (ii) Let $\varpi'=\zeta_\ell-1$, a uniformiser of $\Q_\ell(\zeta_\ell)$. For $\sigma\in I_\ell$ with $\sigma(\zeta_\ell)=\zeta_\ell^a$, $a=\chi_\ell(\sigma)\bmod\ell$,
$\sigma(\varpi')/\varpi'=1+\zeta_\ell+\cdots+\zeta_\ell^{a-1}\equiv a\pmod{\varpi'}$. Now $\ell=\prod_{a=1}^{\ell-1}(1-\zeta_\ell^a)$ and $(1-\zeta_\ell^a)/(1-\zeta_\ell)\equiv a\pmod{\varpi'}$, with $(\ell-1)!\equiv-1\pmod\ell$ (Wilson); as $\ell-1$ is even,
this gives $\varpi'^{\,\ell-1}=-\ell\cdot w$ with $w\equiv1\pmod{\varpi'}$, so $\varpi'=\zeta'\varpi_1w'$ where
$(\zeta')^{\ell-1}=-1$ and $(w')^{\ell-1}=w$; both $\zeta'$ and $w'$ lie in $\Q_\ell^{\ur}$ (Hensel's lemma, $\ell\nmid\ell-1$) and are fixed by $I_\ell$. Hence
$\omega(\sigma)=\sigma(\varpi')/\varpi'\bmod\mathfrak m=a$. (iii) $\varpi_2^{\ell+1}$ is a root of $X^{\ell-1}=\ell$, so $\omega_2^{\ell+1}=\omega$; the order is $\ell^2-1$ by (i). (iv) Put
$F^{-1}(\varpi_2)=\zeta\varpi_2$ with $\zeta\in\mu_{\ell^2-1}$. Then $\sigma F^{-1}(\varpi_2)=\zeta\omega_2(\sigma)\varpi_2$ (lifting $\omega_2(\sigma)$ to $\mu_{\ell^2-1}$, fixed by $\sigma$), and
applying $F$, which acts on $\mu_{\ell^2-1}\subset\Q_\ell^{\ur}$ by $x\mapsto x^\ell$, gives $F\sigma F^{-1}(\varpi_2)=F(\zeta)\omega_2(\sigma)^\ell F(\varpi_2)=\omega_2(\sigma)^\ell\varpi_2$, since
$F(\zeta)F(\varpi_2)=F(\zeta\varpi_2)=\varpi_2$. A character of $G_{\Q_\ell}$ is invariant under conjugation, and $\omega_2^\ell\ne\omega_2$ because $\ell^2-1\nmid\ell-1$; so $\omega_2$
does not extend. The restriction of $\Ind\psi$ to $I_\ell\subset G_{\Q_{\ell^2}}$ is $\psi|_{I_\ell}\oplus\psi^F|_{I_\ell}=\omega_2\oplus\omega_2^\ell$, two distinct characters, and a
$G_{\Q_\ell}$-stable line would be stable under $I_\ell$, hence one of the two eigenlines, which $F$ exchanges. So $\Ind\psi$ is irreducible.
\end{proof}

\begin{theorem}[Deligne, Fontaine; the case of elliptic curves]\label{thm:galrep:deligne-fontaine}
Let $\ell$ be odd and $\ell\nmid N$. For $x\in\Fl^\times$ let $\mu(x)$ be the unramified character of $G_{\Q_\ell}$ with $\mu(x)(\Frob_\ell)=x$.
\begin{enumerate}[label=(\roman*)]
\item (Ordinary case.) If $\ell\nmid a_\ell$, then
\[
\bar\rho|_{D_\ell}\simeq\begin{pmatrix}\bar\chi_\ell\,\mu(a_\ell^{-1})&*\\0&\mu(a_\ell)\end{pmatrix},\qquad\bar\rho|_{I_\ell}\simeq\begin{pmatrix}\omega&*\\0&1\end{pmatrix}.
\]
\item (Supersingular case.) If $\ell\mid a_\ell$ (then $a_\ell=0$ for $\ell\ge5$, and $a_3\in\{0,\pm3\}$), then $\bar\rho|_{I_\ell}\otimes\F_{\ell^2}\simeq\omega_2\oplus\omega_2^\ell$ and
$\bar\rho|_{D_\ell}$ is irreducible.
\item Hence $a_\ell\bmod\ell$ is determined by $\bar\rho|_{D_\ell}$: it is $0$ if $\bar\rho|_{D_\ell}$ is irreducible, and otherwise it is the eigenvalue of $\Frob_\ell$ on
the unique unramified one-dimensional quotient of $\bar\rho|_{D_\ell}$.
\end{enumerate}
The same holds for a newform $g$ of weight $2$ and level prime to $\ell$ with coefficients in $\mathcal O$, with $a_\ell(g)\bmod\lambda$ and
$\bar\rho_{g,\lambda}$ (Deligne in the ordinary case, Fontaine in the supersingular case).
\end{theorem}

\begin{proof}[Sketch of proof]
(i) The finite flat group scheme $\mathcal E[\ell]$ over $\Zl$ has a connected-\'etale sequence $0\to\mathcal E[\ell]^0\to\mathcal E[\ell]\to\mathcal E[\ell]^{\mathrm{et}}\to0$. The
\'etale quotient is determined by its special fibre, $\tilde E[\ell](\Fbar_\ell)$, which has order $\ell$ exactly when $\tilde E$ is ordinary, that is, when
$\ell\nmid a_\ell$ (\Cref{ch:ellcurves}). On it $\Frob_\ell$ acts as the Frobenius endomorphism, by the unit root $\alpha_u$ of $X^2-a_\ell X+\ell$ (the root which is
an $\ell$-adic unit; the other root is $\ell/\alpha_u$), and $\alpha_u\equiv a_\ell\pmod\ell$ because $\alpha_u+\ell/\alpha_u=a_\ell$. So the quotient is $\mu(a_\ell)$, unramified, and the
sub is $\det\bar\rho/\mu(a_\ell)=\bar\chi_\ell\mu(a_\ell^{-1})$, which restricts to $I_\ell$ as $\omega$ (\Cref{prop:galrep:fundamental-characters}(ii)). (ii) If $\tilde E$ is
supersingular, $\mathcal E[\ell]$ is connected and the formal group of $\mathcal E$ has height $2$; Serre \cite[Section 1]{Serre72} shows that $I_\ell$ acts on
$E[\ell]\otimes\F_{\ell^2}$ through $\omega_2\oplus\omega_2^\ell$. By \Cref{prop:galrep:fundamental-characters}(iv) no line is stable under $D_\ell$ even over $\F_{\ell^2}$ (a
stable line would be stable under $I_\ell$ and $F$), so $\bar\rho|_{D_\ell}$ is irreducible. The Hasse bound $|a_\ell|\le2\sqrt\ell<\ell$ for $\ell\ge5$ gives $a_\ell=0$; for
$\ell=3$, $|a_3|\le3$ and $3\mid a_3$. (iii) If $\bar\rho|_{D_\ell}$ is irreducible we are in case (ii) and $a_\ell\equiv0$. Otherwise we are in case (i) (in case (ii)
the restriction is irreducible), and the one-dimensional quotients of $\bar\rho|_{D_\ell}$ are among the two characters of its semisimplification,
$\bar\chi_\ell\mu(a_\ell^{-1})$ and $\mu(a_\ell)$; the first is ramified because $\omega\ne1$ for $\ell$ odd, so $\mu(a_\ell)$ is the unique unramified one-dimensional
quotient, and its value at $\Frob_\ell$ is $a_\ell\bmod\ell$. For newforms the description of $\bar\rho_{g,\lambda}|_{D_\ell}$ is due to Deligne (ordinary) and Fontaine
(supersingular); see \cite[Section 2]{Edixhoven92}, as cited by the paper, and \Cref{ch:modgalrep}. We do not prove these. (The paper writes $h$ for this
newform; here it is $g$, and $h$ is the matrix $\operatorname{diag}(1,-1)$.)
\end{proof}

\begin{paperlink}
Lemma lem:old, the case $p=\ell$: ``By the theorems of Deligne and Fontaine describing $\bar\rho_f|_{D_\ell}$ and $\bar\rho_h|_{D_\ell}$
\cite[\S2]{Edixhoven92}, the residue class of $a_\ell$ is determined by the restriction to $D_\ell$: it is $0$ when that restriction is irreducible,
and otherwise the eigenvalue of Frobenius on the unramified quotient.'' This is \Cref{thm:galrep:deligne-fontaine}(iii). The paper then says why that quotient is
well defined: ``the two characters of the semisimplification restrict to $I_\ell$ as $\bar\chi_\ell$ and $1$, and $\bar\chi_\ell|_{I_\ell}\ne1$ since $\ell$ is odd, so the
unramified quotient is the module of $I_\ell$-coinvariants''. In the notation of this chapter $\bar\chi_\ell|_{I_\ell}=\omega$, the fundamental character of level $1$, which is
nontrivial because it has order $\ell-1\ge2$. Example:
$\bar\rho_{\curve{121c1},3}\simeq\bar\rho_{\curve{11a1},3}$ (\Cref{ex:galrep:serre-conductor-examples}), both curves are ordinary at $3$, and indeed
$a_3(\curve{121c1})=2\equiv-1=a_3(\curve{11a1})\pmod3$.
\end{paperlink}

\begin{corollary}\label{cor:galrep:ordinary}
Let $\ell$ be odd and $\ell\nmid N$.
\begin{enumerate}[label=(\roman*)]
\item If $\bar\rho_{E,\ell}$ is reducible, then $E$ is ordinary at $\ell$, and the two characters of $\bar\rho^{\mathrm{ss}}$ restrict to $I_\ell$ as $\omega$ and $1$.
\item If $\bar\rho_{E,\ell}\simeq\bar\rho_{E,\ell}\otimes\chi_{\ell^*}$, then $\ell=3$ and $E$ is ordinary at $3$.
\end{enumerate}
\end{corollary}

\begin{proof}
(i) In the supersingular case $\bar\rho|_{D_\ell}$, hence $\bar\rho$, is irreducible. In the ordinary case the characters of $\bar\rho^{\mathrm{ss}}$, restricted to $I_\ell$, are
the composition factors of $\bar\rho|_{I_\ell}$, which are $\omega$ and $1$ by \Cref{thm:galrep:deligne-fontaine}(i) (Jordan-H\"older, \Cref{ch:modules}).

(ii) Restrict to $I_\ell$ and compare semisimplifications over $\F_{\ell^2}$. As $\chi_{\ell^*}=\bar\chi_\ell^{(\ell-1)/2}$ (\Cref{conv:galrep:notation}(d)),
$\chi_{\ell^*}|_{I_\ell}=\omega^{(\ell-1)/2}$, which is the unique character of order $2$ of the cyclic group $\omega(I_\ell)=\Fl^\times$: so $\omega^{(\ell-1)/2}\ne1$. In the ordinary
case the multisets $\{\omega,1\}$ and $\{\omega^{(\ell+1)/2},\omega^{(\ell-1)/2}\}$ must agree; as $\omega^{(\ell-1)/2}\ne1$, the character $1$ must equal $\omega^{(\ell+1)/2}$, so
$\ell-1\mid(\ell+1)/2$, which for $\ell\ge3$ forces $\ell-1\le(\ell+1)/2$, that is, $\ell=3$. In the supersingular case $\omega^{(\ell-1)/2}=\omega_2^{(\ell+1)(\ell-1)/2}=-1$ (the
element of order $2$ of $\F_{\ell^2}^\times$), and $\{\omega_2,\omega_2^\ell\}=\{-\omega_2,-\omega_2^\ell\}$ with $-\omega_2\ne\omega_2$ would need $\omega_2^\ell=-\omega_2$, that is,
$\omega_2^{\ell-1}=-1=\omega_2^{(\ell^2-1)/2}$; as $\omega_2$ has order $\ell^2-1$ this means $\ell^2-1\mid(\ell^2-1)/2-(\ell-1)=(\ell-1)^2/2$, that is, $\ell+1\mid(\ell-1)/2$,
impossible since $0<(\ell-1)/2<\ell+1$. So only the ordinary case with $\ell=3$ remains.
\end{proof}

Both parts were checked on the data: every curve of conductor below $10{,}000$ with a Borel or split Cartan label at an odd prime
$\ell\nmid N$ ($4{,}571$ pairs of a curve and a prime: $3{,}628$ at $\ell=3$, $637$ at $5$, $238$ at $7$, $60$ at $13$, $8$ at $37$) has $\ell\nmid a_\ell$; and the
$62$ curves with label \texttt{3Ns} and $3\nmid N$ all have $3\nmid a_3$.

\section[Conductors of mod \texorpdfstring{$\ell$}{l} representations]{Conductors of mod \texorpdfstring{$\ell$}{l} representations}\label{sec:galrep:conductors}

The optimal-level theorem lowers the level of a newform to the Serre conductor $N(\bar\rho)$, and every case of Theorem thm:qual is a
computation of $N(\bar\rho)$, or of $N(\bar\rho\otimes\chi_{q^*})$, at one prime. This section defines the Serre conductor and proves the formula
$a_p(\bar\rho_{E,\ell})=f_p(E)-\dim\Phig_p[\ell]$, which contains the divisibility $N(\bar\rho_{E,\ell})\mid N$ that the paper quotes from \cite{Carayol89}, and every
local exponent the paper uses.

\begin{definition}[Serre conductor]\label{def:galrep:serre-conductor}
Let $\bar\sigma\colon G_\Q\to\GL_n(\Fbar_\ell)$ be continuous (so its image is finite), on a space $W$, and let $p\ne\ell$. The \emph{Artin exponent} of $\bar\sigma$ at $p$ is
\[
a_p(\bar\sigma)=\dim W-\dim W^{I_p}+\operatorname{sw}_p(\bar\sigma),
\]
where the Swan conductor $\operatorname{sw}_p$ is defined from the ramification groups of the finite image of $D_p$ exactly as for complex
representations (\Cref{ch:local}). The \emph{Serre conductor} is $N(\bar\sigma)=\prod_{p\ne\ell}p^{a_p(\bar\sigma)}$; it is prime to $\ell$ by definition. The
\emph{Serre weight} $k(\bar\sigma)$ of an odd irreducible two-dimensional $\bar\sigma$ is defined from $\bar\sigma|_{I_\ell}$ by Serre's recipe \cite[Section 2]{Serre87};
we use it only through \Cref{thm:galrep:good-at-l}(iii).
\end{definition}

The definition copies the conductor of an $\ell$-adic or complex representation, so that for $\bar\rho_{E,\ell}$ it can be compared term by term with the conductor
exponent $f_p(E)$ of $V_\ell E$; the comparison is \Cref{prop:galrep:serre-conductor-E}. The prime $\ell$ is left out because ramification at $\ell$ is measured by
the weight instead. Small examples, all one-dimensional: a quadratic character $\chi_d$ at an odd prime $p\mid d$ has $W^{I_p}=0$ and tame ramification, so
$a_p(\chi_d)=1$; at $p=2$, $\chi_{-4}$ and $\chi_{\pm8}$ have conductors $4$ and $8$, so $a_2(\chi_{-4})=2$ and $a_2(\chi_{\pm8})=3$, the excess over $1$ being the Swan conductor.
For instance, $\chi_{-3}$ viewed modulo $\ell\ne3$ has Serre conductor $3$.

\begin{proposition}\label{prop:galrep:serre-conductor-basic}
\begin{enumerate}[label=(\roman*)]
\item $a_p(\bar\sigma)\ge0$, and $a_p(\bar\sigma)=0$ iff $\bar\sigma$ is unramified at $p$. So $N(\bar\sigma)$ is prime to every prime at which $\bar\sigma$ is unramified.
\item $a_p$ is unchanged by extension of scalars $\Fbar_\ell\supseteq k'\supseteq k$ and by twisting by a character unramified at $p$.
\item If the image of the wild inertia group $P_p$ acts trivially, $a_p(\bar\sigma)=\dim W-\dim W^{I_p}$.
\end{enumerate}
\end{proposition}

\begin{proof}
(i) Both $\dim W-\dim W^{I_p}$ and $\operatorname{sw}_p$ are nonnegative (\Cref{ch:local}). If $\bar\sigma$ is unramified, $W^{I_p}=W$ and every ramification group acts
trivially, so $\operatorname{sw}_p=0$ and $a_p=0$. Conversely $a_p=0$ forces $W^{I_p}=W$. (ii) The fixed spaces of the ramification groups commute with
extension of scalars (they are kernels of linear maps defined over the smaller field), and they are unchanged by a twist which is trivial on
$I_p$ and its subgroups. (iii) If $P_p$ acts trivially, so do all higher ramification groups, which lie in $P_p$, and $\operatorname{sw}_p=0$.
\end{proof}

\begin{proposition}[Serre conductor of $\bar\rho_{E,\ell}$]\label{prop:galrep:serre-conductor-E}
Let $p\ne\ell$ be primes. Then
\[
a_p(\bar\rho_{E,\ell})=f_p(E)-\dim_{\Fl}\Phig_p[\ell].
\]
Hence $N(\bar\rho_{E,\ell})$ divides the prime-to-$\ell$ part of $N$, with $v_p(N(\bar\rho_{E,\ell}))=v_p(N)$ unless $\ell\mid\#\Phig_p$. Explicitly the exponent is: $0$ at a
good prime; $0$ or $1$ at a multiplicative prime according as $\ell\mid n_p$ or not; and, for $\ell$ odd, at an additive prime $f_p-1$ if $\ell=3$ and the type
is $\kod{IV}$ or $\kod{IV}^*$, and $f_p$ otherwise. In particular $\bar\rho_{E,\ell}$, $\ell$ odd, is ramified at every prime $p\ne\ell$ of additive reduction.
\end{proposition}

\begin{proof}
The image $\rho(P_p)$ of the wild inertia group is a pro-$p$ subgroup of $\GL_2(\Zl)$; it meets the pro-$\ell$ group $1+\ell M_2(\Zl)$ trivially, so it injects
into $\GL_2(\Fl)$ and is a finite $p$-group, of order prime to $\ell$. For a finite group $H$ of order prime to $\ell$ acting on $T$, taking invariants
commutes with reduction: the averaging operator $e_H=\frac1{\#H}\sum_{h\in H}h$ is defined over $\Zl$, and $T^H=e_HT$, $(T/\ell T)^H=e_H(T/\ell T)=T^H/\ell T^H$, so
$\dim_{\Fl}(T/\ell)^H=\rank_{\Zl}T^H=\dim_{\Ql}V^H$ (\Cref{ch:reps}). (In more detail: $e_H$ is an idempotent commuting with $H$, so $T=e_HT\oplus(1-e_H)T$; an element fixed
by $H$ satisfies $e_Hx=x$, and conversely $he_H=e_H$ for $h\in H$; the same holds on $T/\ell T$, and reducing the decomposition of $T$ modulo $\ell$ shows
$e_H(T/\ell T)=e_HT/\ell e_HT$, a free module over $\Fl$ of rank $\rank e_HT$.) The Swan conductor is $\sum_{i\ge1}\frac{\#G_i}{\#G_0}\dim(W/W^{G_i})$, where the $G_i$ are the ramification groups of a
finite Galois extension $L/\Q_p$ whose own wild inertia group acts trivially on $T$ (\Cref{ch:local}); one such $L$ serves for $V$ and for $E[\ell]$ at once. For
$i\ge1$ the group $G_i$ lies in $G_1$, the image of $P_p$, and acts on $T$ through $\rho(P_p)$, of order prime to $\ell$, so the codimension of its fixed space is the same on
$E[\ell]$ and on $V$. This gives
$\operatorname{sw}_p(\bar\rho)=\operatorname{sw}_p(V)$. Then, with \Cref{cor:galrep:delta},
\[
a_p(\bar\rho)=2-\dim E[\ell]^{I_p}+\operatorname{sw}_p(V)=f_p-\bigl(\dim E[\ell]^{I_p}-\delta_p\bigr).
\]
It remains to show $\dim E[\ell]^{I_p}-\delta_p=\dim\Phig_p[\ell]$. Good reduction: $E[\ell]^{I_p}=E[\ell]$, $\delta_p=2$, $\Phig_p=0$. Multiplicative: by
\Cref{lem:galrep:component-invariants}(ii) $\dim E[\ell]^{I_p}=1+\dim\Phig_p[\ell]$, and $\delta_p=1$. Additive: by \Cref{lem:galrep:component-invariants}(i)
$\dim E[\ell]^{I_p}=\dim\Phig_p[\ell]$, and $\delta_p=0$. The explicit values follow from $\Phig_p\simeq\Z/n_p$ at multiplicative primes and
\Cref{lem:galrep:component-invariants}(iii) at additive ones; $f_p\ge2$ at an additive prime, and $\dim\Phig_p[\ell]\le1$ there for $\ell$ odd, so the exponent is
at least $1$.
\end{proof}

\begin{paperlink}
The optimal-level paragraph of Section 3: ``Recall also that $N(\bar\rho_{E',\ell})$ divides the conductor of $E'$ \cite{Carayol89}, and that
$N(\bar\sigma)$ is prime to every prime at which $\bar\sigma$ is unramified.'' The second statement is
\Cref{prop:galrep:serre-conductor-basic}(i), the first is \Cref{prop:galrep:serre-conductor-E} (Carayol proved it for the residual representations
of all newforms; for elliptic curves the formula above is more precise, and is the form in which Livn\'e \cite{Livne89} and Serre
\cite[Section 4]{Serre87} discuss it). Lemma lem:IV's ``the Serre conductor of $\bar\rho_{E,3}$ has exponent $1$ at $q$'' is the case $\kod{IV}$, $\kod{IV}^*$; Lemma
lem:old's hypothesis ``$p\nmid M$ for every prime $p\,\|\,N$ at which $\bar\rho$ is unramified'' is the multiplicative case with $\ell\mid n_p$.
\end{paperlink}

\begin{example}\label{ex:galrep:serre-conductor-examples}
The following data were recomputed with PARI: the Kodaira types with the function \texttt{elllocalred}, and the exponents
from \Cref{prop:galrep:serre-conductor-E}. Every congruence was checked for all primes $p<2000$ not dividing the two conductors and different from
$\ell$, except at $\ell=5$, where $p<1224$, $p\ne2,5,7$.
\begin{center}\small
\begin{tabular}{llll}
\toprule
curve & local data & $N(\bar\rho_{E,\ell})$ & residual congruence\\
\midrule
\multicolumn{4}{l}{$\ell=3$}\\
\curve{11a1} & $\kod{I}_5$ at $11$ & $11$ & \\
\curve{121c1} & $\kod{IV}$ at $11$, $f_{11}=2$ & $11$ & $a_p(\curve{121c1})\equiv a_p(\curve{11a1})$\\
\curve{121a1} & $\kod{II}$ at $11$ & $11^2$ & $a_p(\curve{121a1})\equiv\chi_{-11}(p)a_p(\curve{11a1})$\\
\curve{121c2} & $\kod{IV}^*$ at $11$ (twist of \curve{121a1} by $\chi_{-11}$) & $11$ & \\
\curve{77c1} & $\kod{I}_3$ at $7$, $\kod{I}_2$ at $11$ & $11$ & $a_p(\curve{77c1})\equiv a_p(\curve{11a1})$\\
\curve{539c1} & $\kod{I}_3^*$ at $7$, $\kod{I}_2$ at $11$ & $7^2\cdot11$ & $a_p(\curve{539c1})\equiv\chi_{-7}(p)a_p(\curve{77c1})$\\
\curve{52a1} & $\kod{IV}^*$ at $2$, $\kod{I}_2$ at $13$ & $2\cdot13$ & $a_p(\curve{52a1})\equiv a_p(\curve{26b1})$\\
\curve{726a1} & $\kod{I}_6$ at $2$, $\kod{I}_3$ at $3=\ell$, $\kod{III}$ at $11$ & $11^2$ & \\
\curve{6912d1} & $\kod{III}$ at $2$ with $f_2=8$, $\kod{II}$ at $3=\ell$ & $2^8$ & \\
\midrule
\multicolumn{4}{l}{$\ell=5$}\\
\curve{11a1} & $\kod{I}_5$ at $11$; $E[5]\simeq\Z/5\oplus\mu_5$ & $1$ & \\
\curve{140b1} & $\kod{IV}^*$ at $2$, $\kod{I}_1$ at $5=\ell$, $\kod{I}_5$ at $7$ & $2^2$ & $a_p(\curve{140b1})\equiv a_p(\curve{20a1})$\\
\curve{980i1} & $\kod{IV}^*$ at $2$, $\kod{I}_1$ at $5=\ell$, $\kod{I}_5^*$ at $7$ & $2^2\cdot7^2$ & $a_p(\curve{980i1})\equiv\chi_{-7}(p)a_p(\curve{20a1})$\\
\bottomrule
\end{tabular}
\end{center}
For instance, for \curve{121c1} the exponent at $11$ is $f_{11}-\dim\Phig_{11}[3]=2-1=1$, so $\bar\rho_{\curve{121c1},3}$ has Serre conductor $11$ and is isomorphic
to $\bar\rho_{\curve{11a1},3}$, irreducible (\Cref{ex:galrep:eleven}), by \Cref{prop:galrep:irreducibility}(iii); for \curve{121a1} the exponent is $2$, and the
twist by $\chi_{-11}$ (to \curve{121c2}, type $\kod{IV}^*$) lowers it to $1$. For \curve{726a1}, the multiplicative prime $2$ has $3\mid n_2=6$ and drops out, and
$3=\ell$ is excluded by definition. For \curve{11a1} at $\ell=5$, $5\mid n_{11}=5$, so $\bar\rho_{\curve{11a1},5}$ is unramified everywhere outside $5$, as $E[5]\simeq\Z/5\oplus\mu_5$
requires. At $\ell=5$ the curves \curve{140b1} and \curve{980i1} have $5\mid N$, so they illustrate the mechanism of Theorem thm:qual (c), not the theorem.
\end{example}

\section{The images of \texorpdfstring{$\bar\rho_{E,\ell}$}{rhobar}}\label{sec:galrep:images}

Which subgroups of $\GL_2(\Fl)$ occur as images of $\bar\rho_{E,\ell}$ is a deep question in general and an elementary one at $\ell=3$. The paper
needs the answer at $\ell=3$ in Proposition prop:ns (the three irreducible images), in Section sec:anomaly (the curves with image \texttt{3Ns}),
in Section sec:eis (the reducible images \texttt{3B.1.1} and \texttt{3Cs.1.1}) and in Section sec:data (``the images of the mod $\ell$ Galois
representations''). This section states the general theorems of Serre, Mazur and Zywina, proves the list at $\ell=3$ for every $E/\Q$ from the
group theory of \Cref{ch:groups}, fixes the dictionary with Sutherland's labels, and records how images are computed.

\begin{theorem}[Serre; Mazur]\label{thm:galrep:open-image}
\begin{enumerate}[label=(\roman*)]
\item (Serre) If $E/\Q$ has no complex multiplication, $\rho_{E,\ell}(G_\Q)$ is open in $\GL_2(\Zl)$ for every $\ell$, and $\bar\rho_{E,\ell}$ is surjective for all but
finitely many $\ell$.
\item If $E$ has complex multiplication by an order in an imaginary quadratic field $K$ and $\ell$ is odd and unramified in $K$ (then $\ell$ is prime to
the conductor of the order, which for $E/\Q$ is $1$, $2$ or $3$, and $3$ only for $K=\Q(\sqrt{-3})$, \Cref{ch:ellcurves}), the image of
$\bar\rho_{E,\ell}$ lies in the normaliser of a Cartan subgroup, split or nonsplit according as $\ell$ splits or is inert in $K$, and is not contained
in the Cartan subgroup itself.
\item (Mazur) $E/\Q$ has a rational isogeny of prime degree $\ell$ (equivalently, $\bar\rho_{E,\ell}$ is reducible, for $\ell$ odd) only for
$\ell\in\{2,3,5,7,11,13,17,19,37,43,67,163\}$.
\item Whether $\bar\rho_{E,\ell}$ is surjective for every $\ell>37$ and every $E$ without complex multiplication (Serre's uniformity question) is open.
\end{enumerate}
\end{theorem}

We do not prove this; see Serre \cite{Serre72} for (i) and (ii) and Mazur \cite{Mazur78} for (iii). For (ii): the endomorphisms defined over $K$
commute with $G_K$, so $\bar\rho(G_K)$ lies in the centraliser $(\mathcal O_K/\ell)^\times$ of the CM action, a Cartan subgroup; $G_\Q$ normalises it; and
complex conjugation conjugates the CM action into its complex conjugate, so it is not in the Cartan subgroup. (iv) is stated only to mark the
boundary of what is known.

\begin{proposition}[constraints on the image]\label{prop:galrep:image-constraints}
Let $\ell$ be odd and $H=\bar\rho_{E,\ell}(G_\Q)$. Then $\det H=\Fl^\times$ and $H$ contains an element conjugate to $\operatorname{diag}(1,-1)$. Hence $H\not\subseteq\SL_2(\Fl)$ and
$H\not\subseteq C_{ns}(\Fl)$; and either $H=\GL_2(\Fl)$, or $H$ lies in a Borel subgroup ($\bar\rho$ reducible), or $H$ lies in the normaliser of a split or nonsplit
Cartan subgroup and, if $\bar\rho$ is irreducible, not in the Cartan subgroup, or the image of $H$ in $\PGL_2(\Fl)$ is $A_4$, $S_4$ or $A_5$.
\end{proposition}

\begin{proof}
$\det H=\Fl^\times$ and the element $\bar\rho(c)\sim\operatorname{diag}(1,-1)$ are \Cref{prop:galrep:rho-basic}(ii), (iii); $\operatorname{diag}(1,-1)$ has determinant $-1\ne1$. The
nonsplit Cartan subgroup $C_{ns}(\Fl)\simeq\F_{\ell^2}^\times$ is cyclic, so its only involution is $-1$, of determinant $1$; hence it does not contain a conjugate
of $\operatorname{diag}(1,-1)$. By Dickson's classification (\Cref{ch:groups}), if $\ell\mid\#H$ then $H$ lies in a Borel subgroup or contains $\SL_2(\Fl)$; in the
second case $\det H=\Fl^\times$ forces $H=\GL_2(\Fl)$ ($H/\SL_2(\Fl)$ injects into $\Fl^\times$ by the determinant, and is onto). If $\ell\nmid\#H$, $H$ lies in the
normaliser of a Cartan subgroup or has projective image $A_4$, $S_4$ or $A_5$. If $H$ lies in a split Cartan subgroup it is diagonal, so $\bar\rho$ is
reducible; and it is not in a nonsplit one.
\end{proof}

The normaliser $N_s$ is not always a maximal subgroup: $N_s(\F_3)\subset N_{ns}(\F_3)$ up to conjugacy ($N_s(\F_3)$, of order $8$, lies in a $2$-Sylow
subgroup of $\GL_2(\F_3)$, and the $2$-Sylow subgroups, of order $16$, are the conjugates of $N_{ns}(\F_3)$), and $N_s(\F_5)$ lies in a subgroup with projective image
$S_4$. So the large-$\ell$ picture of maximal subgroups must not be used at small $\ell$.

\begin{theorem}[the mod $3$ images]\label{thm:galrep:mod3-images}
For every $E/\Q$, with or without complex multiplication, the image $H=\bar\rho_{E,3}(G_\Q)$ is conjugate in $\GL_2(\F_3)$ to exactly one of the eight
subgroups in the following table; in particular never to $C_{ns}(\F_3)$ (label \texttt{3Cn}) and never to a subgroup of $\SL_2(\F_3)$. The orbits of $H$
on the four lines of $E[3]$ correspond to the irreducible factors of $\psi_3$ over $\Q$, the orbit lengths being their degrees, and the group
$H/(H\cap\{\pm1\})$ acting on the lines is the Galois group of $\psi_3$.
\begin{center}\footnotesize\setlength{\tabcolsep}{3pt}
\begin{tabular}{llllll}
\toprule
label & order & subgroup $H$ & arithmetic meaning & degrees of factors of $\psi_3$ & $\Gal(\psi_3)$\\
\midrule
(surjective) & $48$ & $\GL_2(\F_3)$ & generic & $4$ & $S_4$\\
\texttt{3Nn} & $16$ & $N_{ns}(\F_3)$ & irreducible, $\bar\rho\simeq\bar\rho\otimes\chi_d$, $d\ne-3$ & $4$ & $D_4$\\
\texttt{3Ns} & $8$ & $N_s(\F_3)$ & irreducible, $\bar\rho\simeq\bar\rho\otimes\chi_{-3}$ & $2+2$ & $V_4$\\
\texttt{3B} & $12$ & $\mathrm B(\F_3)$ & stable line, character $\ne1,\chi_{-3}$ & $1+3$ & $S_3$\\
\texttt{3B.1.1} & $6$ & $\sm{1}{*}{0}{*}$ & rational point of order $3$, non-split & $1+3$ & $S_3$\\
\texttt{3B.1.2} & $6$ & $\sm{*}{*}{0}{1}$ & $\mu_3\subset E[3]$, quotient $\Z/3$, non-split & $1+3$ & $S_3$\\
\texttt{3Cs} & $4$ & $C_s(\F_3)$ & $E[3]\simeq\eta\oplus\eta\chi_{-3}$, $\eta\ne1,\chi_{-3}$ & $1+1+2$ & $C_2$\\
\texttt{3Cs.1.1} & $2$ & $\{\operatorname{diag}(1,\pm1)\}$ & $E[3]\simeq\Z/3\oplus\mu_3$ & $1+1+2$ & $C_2$\\
\bottomrule
\end{tabular}
\end{center}
\end{theorem}

\begin{proof}
By \Cref{prop:galrep:image-constraints}, $H$ has surjective determinant and contains an element of determinant $-1$ and trace $0$. By \Cref{ch:groups}
the subgroups of $\GL_2(\F_3)$ with surjective determinant fall into nine conjugacy classes, namely the eight of the table and $C_{ns}(\F_3)$, cyclic of
order $8$, and $C_{ns}(\F_3)$ is the only one of them without an element of determinant $-1$ and trace $0$ (its elements of determinant $-1$ are its
generators, of order $8$, with nonzero trace). To see the last point, identify $C_{ns}(\F_3)$ with $\F_9^\times$ acting on $\F_9\simeq\F_3^2$ by multiplication; then the
determinant of multiplication by $x$ is the norm $x\cdot x^3=x^4$ and its trace is $x+x^3$. If $x$ has order $8$, then $x^4=-1$, and $x+x^3=x(1+x^2)\ne0$
because $x^2$ has order $4$, so $x^2\ne-1$. If $x$ has order dividing $4$, then $x^4=1$. A subgroup of $\SL_2(\F_3)$ does not have surjective determinant. So $H$ is conjugate to one of the
eight. The four lines of $E[3]$ correspond bijectively and Galois-equivariantly to the four roots of $\psi_3$ (\Cref{prop:galrep:irreducibility}(ii)), so
the orbits of $H$ on lines are the sets of roots of the irreducible factors of $\psi_3$, and the permutation group induced by $H$ on the lines, which is
$H/(H\cap\{\pm1\})$ ($\pm1$ are the only elements fixing every line), is the Galois group of $\psi_3$. The orbit structures and groups in the table are
read off from the subgroups (\Cref{ch:groups}): $\GL_2(\F_3)$ and $N_{ns}(\F_3)$ are transitive, with quotients $S_4$ and $D_4$; $N_s(\F_3)$ has the two orbits
$\{\langle e_1\rangle,\langle e_2\rangle\}$ and $\{\langle e_1+e_2\rangle,\langle e_1-e_2\rangle\}$, quotient $V_4$; the Borel groups fix $\langle e_1\rangle$ and permute the other three lines
as $S_3$; the subgroups of $C_s$ fix $\langle e_1\rangle$, $\langle e_2\rangle$ and swap the other two. The arithmetic meanings: a stable line with character $\eta_1$ and
quotient $\eta_2=\chi_{-3}\eta_1^{-1}$ (determinant $\chi_{-3}$) gives image in $\mathrm B$; it is $\sm1*0*$ iff $\eta_1=1$ (the points of the line are rational), $\sm**01$ iff
$\eta_2=1$, i.e.\ $\eta_1=\chi_{-3}$ (the line is $\mu_3$, by the Weil pairing); the image contains the unipotent group iff the extension is non-split; and in
the split case $E[3]\simeq\eta_1\oplus\eta_2$, with image $\{\operatorname{diag}(1,\pm1)\}$ iff $\{\eta_1,\eta_2\}=\{1,\chi_{-3}\}$. The irreducible cases are \Cref{prop:galrep:ns}
and \Cref{prop:galrep:dihedral}: the image is $N_s(\F_3)$ iff $\bar\rho\simeq\bar\rho\otimes\chi_{-3}$, and $N_{ns}(\F_3)$ is the normaliser of the Cartan subgroup $C_{ns}$, cut out
by a quadratic field $\Q(\sqrt d)$ with $\bar\rho\simeq\bar\rho\otimes\chi_d$ and $d\ne-3$ (as $C_{ns}(\F_3)\not\subseteq\SL_2(\F_3)$).
\end{proof}

Every case occurs. Among the $64{,}687$ curves of conductor below $10{,}000$ in Cremona's tables \cite{ecdata}, of which $462$ have complex
multiplication, we classified the image of every curve from $\psi_3$ by the recipe of \Cref{rem:galrep:computing-images} and compared it with
the label in Cremona's file \texttt{galrep}, which follows \cite{Sutherland16}:
\begin{center}\small
\begin{tabular}{lrrl}
\toprule
label & non-CM & CM & examples\\
\midrule
(surjective) & $53{,}760$ & $0$ & \curve{11a1}, \curve{15a1}\\
\texttt{3Nn} & $218$ & $184$ & \curve{245a1}; CM: \curve{32a1}\\
\texttt{3Ns} & $120$ & $26$ & \curve{338d1}, \curve{726a1}, \curve{6912d1}; CM: \curve{121b1}, \curve{256a1}\\
\texttt{3B} & $4{,}072$ & $154$ & \curve{50b1}, \curve{539a1}; CM: \curve{144a1}\\
\texttt{3B.1.1} & $2{,}772$ & $39$ & \curve{14a4}, \curve{20a1}; CM: \curve{36a1}\\
\texttt{3B.1.2} & $2{,}772$ & $39$ & \curve{14a3}, \curve{20a3}\\
\texttt{3Cs} & $249$ & $18$ & \curve{98a3}, \curve{539a2}\\
\texttt{3Cs.1.1} & $262$ & $2$ & \curve{14a1}, \curve{19a1}, \curve{26a1}; CM: \curve{27a1}, \curve{27a3}\\
\bottomrule
\end{tabular}
\end{center}
The classification from $\psi_3$ agrees with the \texttt{galrep} label for every curve that has one. The file gives no mod $3$ label to $364$ of the CM
curves, those whose image is the one forced by the complex multiplication: the $26$ curves with $j\in\{-32768,8000\}$ (image \texttt{3Ns}), the $184$ with
$j\in\{1728,287496,-3375,16581375,-884736,-884736000,-147197952000\}$ (image \texttt{3Nn}) and the $154$ with $j\in\{0,54000,-12288000\}$ (image
\texttt{3B}); the CM column counts all of them. Every count is for conductor below $10{,}000$.

The equal counts of \texttt{3B.1.1} and \texttt{3B.1.2} in the table are not an accident. They hold class by class, because the rational
$3$-isogenies arrange each isogeny class into paths whose two ends carry the two labels.

\begin{proposition}[paths of $3$-isogenies]\label{prop:galrep:3-isogeny-paths}
Let $E/\Q$ be an elliptic curve, with or without complex multiplication. Let $\Gamma$ be the graph whose vertices are the $\Q$-isomorphism classes of the
curves $\Q$-isogenous to $E$, two vertices being joined when there is an isogeny of degree $3$ over $\Q$ between them (the dual isogeny also has degree
$3$, so the relation is symmetric).
\begin{enumerate}[label=(\roman*)]
\item For each vertex $E'$, $E'[3]$ has at most two $G_\Q$-stable lines, and $C\mapsto E'/C$ is a bijection from the set of these lines to the set of
neighbours of $E'$.
\item Each connected component of $\Gamma$ is a path $E_0,E_1,\dots,E_k$ (pairwise distinct vertices, $E_{i-1}$ joined to $E_i$). If $k\ge1$, the two ends
have exactly one stable line in $E_0[3]$, respectively $E_k[3]$, so their images are \texttt{3B}, \texttt{3B.1.1} or \texttt{3B.1.2}; the interior vertices
have two, so their images are \texttt{3Cs} or \texttt{3Cs.1.1}.
\item If $k\ge1$ and $G_\Q$ acts on the stable line of $E_0[3]$ by the character $\eta$, then it acts on the stable line of $E_k[3]$ by $\chi_{-3}\eta^{-1}$.
Hence $E_0$ has image \texttt{3B.1.1} iff $E_k$ has image \texttt{3B.1.2}, $E_0$ has image \texttt{3B.1.2} iff $E_k$ has image \texttt{3B.1.1}, and
$E_0$ has image \texttt{3B} iff $E_k$ has.
\item In every isogeny class over $\Q$ the number of curves with image \texttt{3B.1.1} equals the number with image \texttt{3B.1.2}.
\end{enumerate}
\end{proposition}

\begin{proof}
All isogenies and isomorphisms are over $\Q$. We use two facts from \Cref{ch:ellcurves} and one from \Cref{ch:reduction} and the literature. (F1) A finite $G_\Q$-stable subgroup $C$ of $A(\Qbar)$ is the
kernel of an isogeny $A\to A/C$, and an isogeny is determined by its kernel up to an isomorphism of its target; for an isogeny $\varphi$ of degree $3$,
$\hat\varphi\varphi=[3]$, so $\ker\hat\varphi=\varphi(A[3])$, a line. (F2) Every $A/\Q$ has $\End_\Q(A)=\Z$ (the endomorphisms of a curve with complex
multiplication by $K$ that are not integers are defined only over $K$); so $\Aut_\Q(A)=\{\pm1\}$, and an endomorphism over $\Q$ of degree $3^m$ is
$\pm3^{m/2}$, whose kernel $A[3^{m/2}]$ is not cyclic when $m\ge1$. (F3) An isogeny class over $\Q$ is finite: isogenous curves have the same
conductor (\Cref{ch:reduction}), and there are finitely many curves over $\Q$ of given conductor (Shafarevich's theorem, \cite[Section IX.6]{SilvermanAEC}).

\emph{Step 1: cyclic composites.} Let $\varphi_i\colon A_{i-1}\to A_i$ ($1\le i\le m$) be isogenies of degree $3$ with $\ker\varphi_{i+1}\ne\varphi_i(A_{i-1}[3])$ for
$1\le i<m$. Then $\Phi_m=\varphi_m\circ\cdots\circ\varphi_1$ has cyclic kernel, of order $3^m$. For $m=1$ this is clear. Suppose $\Phi_{m-1}$ has cyclic
kernel and $\ker\Phi_m$ is not cyclic. A finite abelian $3$-group that is not cyclic contains a subgroup $(\Z/3)^2$, and in $A_0$ that subgroup is
$A_0[3]$; so $\Phi_m$ kills $A_0[3]$, that is, $\varphi_m$ kills $\Phi_{m-1}(A_0[3])$. As $\ker\Phi_{m-1}$ is cyclic, it meets $A_0[3]$ in a group of
order $3$, so $\Phi_{m-1}(A_0[3])$ has order $3$; it lies in $\varphi_{m-1}(A_{m-2}[3])$, which also has order $3$, so the two are equal, and
$\ker\varphi_m=\varphi_{m-1}(A_{m-2}[3])$: a contradiction.

(i) If three lines of $E'[3]$ were stable, every element of $H=\bar\rho_{E',3}(G_\Q)$ would have three eigenlines in a plane, hence be a scalar
$\pm1$, of determinant $1$; but $\det H=\F_3^\times$ (\Cref{prop:galrep:rho-basic}(ii)). The neighbours of $E'$ are the curves $E'/C$, by (F1), since the
kernel of an isogeny of degree $3$ is a stable line. Suppose $C\ne C'$ are stable lines and $\iota\colon E'/C\to E'/C'$ is an isomorphism; let
$\varphi\colon E'\to E'/C$ and $\psi\colon E'\to E'/C'$ be the quotient maps. If $\ker(\hat\psi\iota)=\varphi(E'[3])=\ker\hat\varphi$, then by (F1) and (F2)
$\hat\psi\iota=\pm\hat\varphi$; dualising (the dual of an isomorphism is its inverse), $\psi=\pm\iota\varphi$, so $C'=\ker\psi=\ker\varphi=C$, which is false. So
Step 1 applies to $\varphi$ and $\hat\psi\iota$: the endomorphism $\hat\psi\iota\varphi$ of $E'$ has degree $9$ and cyclic kernel, contradicting (F2). Hence
$C\mapsto E'/C$ is injective.

(ii) $\Gamma$ has no cycle. Let $v_0,v_1,\dots,v_m=v_0$ be a cycle, $m\ge3$, with $v_0,\dots,v_{m-1}$ distinct. Choose curves $A_i$ in the classes $v_i$, with
$A_m=A_0$, and isogenies $\varphi_i\colon A_{i-1}\to A_i$ of degree $3$. If $\ker\varphi_{i+1}=\varphi_i(A_{i-1}[3])=\ker\hat\varphi_i$ for some $1\le i<m$, then
$A_{i+1}\simeq A_i/\ker\hat\varphi_i\simeq A_{i-1}$ by (F1), so $v_{i+1}=v_{i-1}$, which is impossible in a cycle of length $m\ge3$. By Step 1,
$\varphi_m\circ\cdots\circ\varphi_1$ is an endomorphism of $A_0$ of degree $3^m$ with cyclic kernel, contradicting (F2). So each component is a tree,
finite by (F3), in which every vertex has degree at most $2$ by (i): a path. By (i) the degree of a vertex is its number of stable lines, $1$ at the
ends and $2$ in the interior when $k\ge1$. One stable line means $E[3]$ is a non-split extension (a split $E[3]$ has the two lines of its
decomposition), two mean that it is split; the labels follow from \Cref{thm:galrep:mod3-images}.

(iii) Let $\varphi_i\colon E_{i-1}\to E_i$ ($1\le i\le k$) be isogenies of degree $3$, $C_{i-1}=\ker\varphi_i\subset E_{i-1}[3]$ and $B_i=\varphi_i(E_{i-1}[3])=\ker\hat\varphi_i
\subset E_i[3]$, both stable lines. The map $\varphi_i$ induces a $G_\Q$-isomorphism $E_{i-1}[3]/C_{i-1}\simeq B_i$, and $\det\bar\rho_{E_{i-1},3}=\chi_{-3}$
(\Cref{prop:galrep:rho-basic}(ii)); so if $G_\Q$ acts on $C_{i-1}$ by $\alpha$, it acts on $B_i$ by $\chi_{-3}\alpha^{-1}$. For $0<i<k$, $E_i/B_i\simeq E_{i-1}$ and
$E_i/C_i\simeq E_{i+1}\ne E_{i-1}$, so $B_i\ne C_i$ by (i), $E_i[3]=B_i\oplus C_i$, and $G_\Q$ acts on $C_i$ by $\chi_{-3}$ divided by its character on
$B_i$. Starting from $\eta$ on $C_0$, induction gives $\chi_{-3}\eta^{-1}$ on every $B_i$ and $\eta$ on every $C_i$; at the end, $B_k$ is the stable line of
$E_k[3]$. By \Cref{thm:galrep:mod3-images} (the arithmetic meanings in its proof), a non-split $E'[3]$ whose stable line has character $\eta$ has image
\texttt{3B.1.1} if $\eta=1$, \texttt{3B.1.2} if $\eta=\chi_{-3}$, and \texttt{3B} otherwise; and $\eta\mapsto\chi_{-3}\eta^{-1}$ exchanges $1$ and $\chi_{-3}$ and maps
every other character to another character outside $\{1,\chi_{-3}\}$.

(iv) A curve with image \texttt{3B.1.1} or \texttt{3B.1.2} has exactly one stable line, so by (i) and (ii) it is an end of a component with $k\ge1$, whose
other end is a different curve. Sending each end to the other end of its component is an involution of the set of such ends in the class, and by
(iii) it exchanges the two labels.
\end{proof}

The data agree. In Cremona's tables below conductor $10{,}000$, the $3$-isogenies read from the isogeny matrices of the file \texttt{allisog} form
$4{,}924$ components with at least one edge in the $38{,}042$ classes; every one is a path whose interior vertices have image \texttt{3Cs} or
\texttt{3Cs.1.1} and whose ends have the images \texttt{3B} and \texttt{3B}, or \texttt{3B.1.1} and \texttt{3B.1.2}; and in every class the two labels
occur equally often.

\begin{remark}
The curves with a rational point of order $3$ are exactly those with image \texttt{3B.1.1} or \texttt{3Cs.1.1} (\Cref{prop:galrep:eis-labels}): $2{,}811$
and $264$ curves below conductor $10{,}000$, as the torsion orders in \cite{ecdata} confirm. The labels \texttt{3B.1.1} and \texttt{3B.1.2} occur equally
often in the table ($2{,}772$ non-CM and $39$ CM curves each), by \Cref{prop:galrep:3-isogeny-paths}(iv). The twist by $\chi_{-3}$ also exchanges them
(\Cref{prop:galrep:twist-isogeny}(ii)), but it changes the conductor in general, so it does not explain the equality below a conductor bound; the pairing
inside the class does. For example, the class \curve{14a} consists of the two paths \curve{14a3}, \curve{14a1}, \curve{14a4} and \curve{14a5}, \curve{14a2},
\curve{14a6}: the ends \curve{14a3}, \curve{14a5} have image \texttt{3B.1.2}, the ends \curve{14a4}, \curve{14a6} image \texttt{3B.1.1}, and the middle curves
\curve{14a1}, \curve{14a2} image \texttt{3Cs.1.1} ($E[3]\simeq\Z/3\oplus\mu_3$; the quotient by $\mu_3$ is the \texttt{3B.1.1} end, which has the rational
point of order $3$ coming from $E[3]/\mu_3$, and the quotient by $\Z/3$ is the \texttt{3B.1.2} end). The
image \texttt{3Cs.1.1} is not the split Cartan subgroup, which is \texttt{3Cs}, but its subgroup $\{1,\operatorname{diag}(1,-1)\}$ of order $2$. The image \texttt{3Nn} is
dihedral too: $\bar\rho$ is induced from the quadratic field cut out by $\bar\rho^{-1}(C_{ns})$ (\Cref{prop:galrep:dihedral}); what fails for \texttt{3Ns}, and
only for it, is irreducibility over $\Q(\sqrt{-3})$ (\Cref{prop:galrep:ns}).
\end{remark}

\begin{convention}[Sutherland's labels]\label{conv:galrep:labels}
A label such as \texttt{3B.1.1} names a conjugacy class of subgroups of $\GL_2(\Fl)$: the prime $\ell$; the type of the smallest standard subgroup
containing it, \texttt{B} (Borel), \texttt{Cs}, \texttt{Cn} (split, nonsplit Cartan), \texttt{Ns}, \texttt{Nn} (their normalisers), or \texttt{A4},
\texttt{S4}, \texttt{A5} (exceptional projective image); then optional numbers identifying the subgroup inside that one. Surjective images have no
label. We give the general rule only by reference \cite{Sutherland16} and always state the subgroup intrinsically; the dictionary at $\ell=3$ is the
table of \Cref{thm:galrep:mod3-images}, which we checked against Cremona's data. At $\ell=5$: \texttt{5Cs.1.1} for \curve{11a1}
($E[5]\simeq\Z/5\oplus\mu_5$), \texttt{5B.1.1} for \curve{11a3} (a rational point of order $5$, $E[5]$ non-split), \texttt{5B.1.2} for \curve{11a2}. The labels at
$\ell=5,7,11,13$ occurring for conductor below $10{,}000$ are:
\begin{center}\small
\begin{tabular}{lp{0.85\linewidth}}
\toprule
$\ell$ & labels\\
\midrule
$5$ & \texttt{5B}, \texttt{5B.1.1}, \texttt{5B.1.2}, \texttt{5B.1.3}, \texttt{5B.1.4}, \texttt{5B.4.1}, \texttt{5B.4.2}, \texttt{5Cs.1.1}, \texttt{5Cs.1.3}, \texttt{5Cs.4.1},
\texttt{5Nn}, \texttt{5Ns}, \texttt{5Ns.2.1}, \texttt{5S4}\\
$7$ & \texttt{7B}, \texttt{7B.1.1} to \texttt{7B.1.6}, \texttt{7B.2.1}, \texttt{7B.2.3}, \texttt{7B.6.1} to \texttt{7B.6.3}, \texttt{7Ns}, \texttt{7Ns.2.1}, \texttt{7Ns.3.1},
\texttt{7Ns.6.1.2}\\
$11$ & \texttt{11B.1.3} to \texttt{11B.1.8}, \texttt{11B.10.3} to \texttt{11B.10.5}, \texttt{11Nn.1.4}\\
$13$ & \texttt{13B}, \texttt{13B.3.1}, \texttt{13B.3.2}, \texttt{13B.4.1}, \texttt{13B.4.2}, \texttt{13B.5.1}, \texttt{13B.5.2}, \texttt{13Ns.5.1.4}\\
\bottomrule
\end{tabular}
\end{center}
\end{convention}

\begin{theorem}[Zywina]\label{thm:galrep:zywina}
For $E/\Q$ without complex multiplication and $\ell\in\{2,3,5,7,11,13\}$, the subgroups of $\GL_2(\Fl)$ that occur as $\bar\rho_{E,\ell}(G_\Q)$ are known up to conjugacy,
together with explicit families of $j$-invariants (rational functions on modular curves, and finitely many exceptional $j$) realising each of them.
\end{theorem}

We do not prove this; see Zywina \cite[Section 1]{Zywina15}. At $\ell=13$ the list in \cite{Zywina15} depended on the rational points of certain modular
curves, which later work determined; we give no reference for that. The book needs only $\ell=3$, which \Cref{thm:galrep:mod3-images} proves directly
and for CM curves as well; the paper cites \cite{Zywina15} for the three irreducible images at $\ell=3$ and derives them directly.

\begin{remark}[computing images]\label{rem:galrep:computing-images}
(a) At $\ell=3$ the image is determined by $\psi_3$: its factorisation pattern, the Galois group of an irreducible quartic factor, and, for each
rational root $x_0$, the character of the line, which is $\chi_D$ with $D=4x_0^3+b_2x_0^2+2b_4x_0+b_6$ (the line is $\{O,\pm P\}$ with
$(2y(P)+a_1x_0+a_3)^2=D$, and $g(P)=\pm P$ according as $g$ fixes $\sqrt D$). The following function returns the label of \Cref{thm:galrep:mod3-images}; on
all $64{,}687$ curves of conductor below $10{,}000$ it agrees with Cremona's \texttt{galrep} file wherever that has a label.
\begin{lstlisting}[language=GP]
mod3image(E) =
{
  my(F = factor(elldivpol(E, 3))[,1], d, f, r, ch);
  d = vecsort([poldegree(g) | g <- F]);
  if (d == [4], return (if (polgalois(F[1])[1] == 24, "GL2", "3Nn")));
  if (d == [2, 2], return ("3Ns"));
  f = 4*x^3 + E.b2*x^2 + 2*E.b4*x + E.b6;
  r = [subst(f, x, -polcoef(g, 0)/polcoef(g, 1)) | g <- F, poldegree(g) == 1];
  ch = [core(numerator(t)*denominator(t)) | t <- r];
  if (#r == 1, return (if (ch[1] == 1, "3B.1.1", if (ch[1] == -3, "3B.1.2", "3B"))));
  if (vecsort(ch) == [-3, 1], "3Cs.1.1", "3Cs");
}
\end{lstlisting}
The paper's appendix uses the shorter test \texttt{ns3}, which returns $1$ iff $\psi_3$ is a product of two quadratics; by the table this detects
$N_s(\F_3)$ exactly (a reducible $\bar\rho$ gives a rational root). (b) For general $\ell$, Sutherland's algorithm \cite{Sutherland16} computes the
image from the statistics of the characteristic polynomials of $\bar\rho(\Frob_p)$ for many $p$ (Chebotarev), together with division polynomials and
Zywina's families; we use its output through Cremona's files. (c) The inertia data of this chapter were computed from the fields $\Q(P)$, $P\in E[3]$:
for each irreducible factor $g$ of $\psi_3$, the field generated by a root $x_0$ and by $\sqrt{D(x_0)}$, and the ramification indices of $q$ in it
(\texttt{nfinit}, \texttt{idealprimedec}). When the image of inertia has order prime to $q$, its order is the least common multiple of these
indices, and $E[3]^{I_q}\ne0$ iff some prime above $q$ in some $\Q(P)$ is unramified.
\end{remark}

\begin{paperlink}
Section sec:data: ``The mod $3$ images of the tables were checked against the factorisation of the $3$-division polynomial (Proposition~prop:ns) on every curve for which it mattered.'' The appendix's
\texttt{ns3}: ``\texttt{ns3(E) = my(F = factor(elldivpol(E, 3))); (\#F[,1] == 2 \&\& poldegree(F[1,1]) == 2);}''. Both rest on
\Cref{thm:galrep:mod3-images}. Section sec:eis calls \texttt{3Cs.1.1} ``the subgroup of order $2$ of the split Cartan subgroup that fixes a point''.
\end{paperlink}

\section[Absolute irreducibility over \texorpdfstring{$\Q(\sqrt{\ell^*})$}{Q(sqrt l*)}]{Absolute irreducibility over \texorpdfstring{$\Q(\sqrt{\ell^*})$}{Q(sqrt l*)}, Proposition prop:ns and Lemma lem:h0}\label{sec:galrep:lstar}

The hypothesis of Theorem thm:main is that $\bar\rho$ stay absolutely irreducible on $\Gal(\Qbar/\Q(\sqrt{\ell^*}))$, which is the condition $\ell\notin S_f$ of
\cite[(31)]{DFG04}. This section shows what that condition means for elliptic curves: by Clifford theory its failure is $\bar\rho\simeq\bar\rho\otimes\chi_{\ell^*}$; by the
local analysis at $\ell$ this happens, for $\ell\nmid N$, only at $\ell=3$ and only for the image \texttt{3Ns}; and it is then measured by the group
$H^0(\Q,\ad T(1)\otimes\Qlmod)$ of Lemma lem:h0, whose length is at most one. It proves the paper's Proposition prop:ns, Lemma lem:h0, Proposition
prop:ns-structure and Proposition prop:Sf, with every compressed step expanded, and explains Section sec:anomaly.

\begin{proposition}[dihedral representations]\label{prop:galrep:dihedral}
Let $\ell$ be odd, $\bar\rho=\bar\rho_{E,\ell}$ irreducible, $K$ a quadratic field and $\chi_K$ its character. The following are equivalent:
\begin{enumerate}[label=(\alph*)]
\item $\bar\rho|_{G_K}$ is not absolutely irreducible;
\item $\bar\rho\simeq\bar\rho\otimes\chi_K$ (over $\Fl$, equivalently over $\Fbar_\ell$);
\item $\bar\rho\otimes\Fbar_\ell\simeq\Ind_{G_K}^{G_\Q}\psi$ for a character $\psi\colon G_K\to\Fbar_\ell^\times$;
\item $a_p(E)\equiv0\pmod\ell$ for all primes $p\nmid N\ell$ inert in $K$ (equivalently, for those outside a set of density zero).
\end{enumerate}
When they hold, every isomorphism $\phi\colon\bar\rho\to\bar\rho\otimes\chi_K$ has trace $0$, and $H^0(\Q,\ad\bar\rho\otimes\chi_K)=\Hom_{G_\Q}(\bar\rho,\bar\rho\otimes\chi_K)$ is a line.
\end{proposition}

\begin{proof}
By \Cref{prop:galrep:irreducibility}(i), $\bar\rho$ is absolutely irreducible; so by Schur's lemma $\End_{G_\Q}(\bar\rho\otimes\Fbar_\ell)=\Fbar_\ell$.

(b)$\Rightarrow$(a). An isomorphism $\phi\colon\bar\rho\to\bar\rho\otimes\chi_K$ satisfies $\phi\bar\rho(g)=\chi_K(g)\bar\rho(g)\phi$, so it commutes with $\bar\rho(G_K)$, and it is not a scalar
(a scalar would force $\chi_K=1$). So $\End_{G_K}(\bar\rho\otimes\Fbar_\ell)$ has dimension at least $2$, and $\bar\rho|_{G_K}$ is not absolutely irreducible.

(a)$\Rightarrow$(c). Let $L\subset\bar\rho\otimes\Fbar_\ell$ be a $G_K$-stable line, with character $\psi$. By Frobenius reciprocity,
$\Hom_{G_\Q}(\Ind_{G_K}^{G_\Q}\psi,\bar\rho\otimes\Fbar_\ell)=\Hom_{G_K}(\psi,\bar\rho|_{G_K}\otimes\Fbar_\ell)\ne0$ (\Cref{ch:reps}). A nonzero map from the two-dimensional $\Ind\psi$
to the absolutely irreducible two-dimensional $\bar\rho\otimes\Fbar_\ell$ is surjective, hence an isomorphism.

(c)$\Rightarrow$(b). $\Ind\psi\otimes\chi_K\simeq\Ind(\psi\otimes\chi_K|_{G_K})=\Ind\psi$, as $\chi_K$ is trivial on $G_K$. So $\Hom_{G_\Q}(\bar\rho,\bar\rho\otimes\chi_K)\otimes\Fbar_\ell\ne0$; the space of
$G_\Q$-maps is the kernel of a system of linear equations over $\Fl$, and its dimension does not change under extension of scalars, so there is a
nonzero $\Fl$-linear $G_\Q$-map $\bar\rho\to\bar\rho\otimes\chi_K$, which is an isomorphism as both are irreducible of the same dimension.

(b)$\Rightarrow$(d). For $p\nmid N\ell$ inert in $K$, $\chi_K(\Frob_p)=-1$, and taking traces at $\Frob_p$ gives $a_p\equiv-a_p$, so $2a_p\equiv0$ and $a_p\equiv0\pmod\ell$.

(d)$\Rightarrow$(b). The representations $\bar\rho$ and $\bar\rho\otimes\chi_K$ have the same determinant ($\chi_K^2=1$) and the same trace $\pm a_p$ at $\Frob_p$ for every
$p\nmid N\ell d_K$: at split primes $\chi_K(p)=1$, and at inert primes both traces are $0$ outside a set of density zero. By
\Cref{prop:galrep:irreducibility}(iii) (Chebotarev and Brauer-Nesbitt), $\bar\rho\simeq\bar\rho\otimes\chi_K$, both being irreducible.

Trace zero. Choose $g\in G_\Q\setminus G_K$; then $\phi\bar\rho(g)=-\bar\rho(g)\phi$, so $\bar\rho(g)\phi\bar\rho(g)^{-1}=-\phi$ and $\tr\phi=-\tr\phi$, whence $\tr\phi=0$ ($\ell$ odd). So
$\Hom_{G_\Q}(\bar\rho,\bar\rho\otimes\chi_K)$ lies in $(\ad\bar\rho\otimes\chi_K)^{G_\Q}$, and conversely; it is a line because any two isomorphisms differ by an automorphism of
$\bar\rho$, a scalar.
\end{proof}

For $K=\Q(\sqrt{\ell^*})$ condition (a) is exactly the failure of the hypothesis of Theorem thm:main. Not every dihedral $\bar\rho$ fails it: the
image \texttt{3Nn} is induced from a quadratic field other than $\Q(\sqrt{-3})$, and those curves satisfy the hypothesis.

\begin{theorem}[the hypothesis on $\Q(\sqrt{\ell^*})$ is automatic for $\ell\ge5$]\label{thm:galrep:lstar-automatic}
Let $\ell$ be odd, $\ell\nmid N$, and $\bar\rho_{E,\ell}$ irreducible. If $\ell\ge5$, then $\bar\rho_{E,\ell}$ is absolutely irreducible on $\Gal(\Qbar/\Q(\sqrt{\ell^*}))$. If $\ell=3$ and
$\bar\rho_{E,3}$ is not absolutely irreducible on $\Gal(\Qbar/\Q(\sqrt{-3}))$, then $E$ is ordinary at $3$ (and the image is $N_s(\F_3)$, \Cref{prop:galrep:ns}).
\end{theorem}

\begin{proof}
If $\bar\rho$ is not absolutely irreducible on $G_{\Q(\sqrt{\ell^*})}$, then $\bar\rho\simeq\bar\rho\otimes\chi_{\ell^*}$ by \Cref{prop:galrep:dihedral}, and \Cref{cor:galrep:ordinary}(ii)
gives $\ell=3$ and ordinary reduction at $3$.
\end{proof}

The data agree. Below conductor $10{,}000$, the curves with a label \texttt{$\ell$Ns} or \texttt{$\ell$Nn}, $\ell\ge5$, for which $a_p\equiv0\pmod\ell$ at each of the first
$150$ primes $p\nmid N\ell$ inert in $\Q(\sqrt{\ell^*})$ (all below $2{,}000$) are $24$ curves (labels \texttt{5Ns.2.1}, \texttt{7Ns.2.1}, \texttt{7Ns.3.1},
\texttt{7Ns.6.1.2}), and all of them have $\ell\mid N$. The $28$ curves with label \texttt{5Ns}, \texttt{5Nn} or \texttt{7Ns} and $\ell\nmid N$ are induced from
other quadratic fields: for instance \curve{608b1} (\texttt{5Ns}) from $\Q(i)$, \curve{1044a1} (\texttt{5Nn}) from $\Q(\sqrt{-3})$, and \curve{9225a1}
(\texttt{7Ns}) from $\Q(\sqrt{-3})$.

\begin{paperlink}
The paper states this as Proposition prop:Sf: ``Let $\ell$ be an odd prime with $\ell\nmid N$ and $E[\ell]$ irreducible. If
$\ell\ge5$, then $\ell\notin S_f$. If $\ell=3$, then $3\in S_f$ if and only if the image of $\bar\rho_{E,3}$ is conjugate to $N_s(\F_3)$, and then $E$ is ordinary at
$3$.'' Its proof is the argument of \Cref{cor:galrep:ordinary}(ii). Accordingly the paper's abstract reads ``(for $\ell\ge5$ this only asks that
$E[\ell]$ be irreducible)'', and Section sec:forms records that there are no other failures of the hypothesis among the non-CM curves with
$E[\ell]$ irreducible, ``the curves whose mod $\ell$ representation is induced from $\Q(\sqrt{\ell^*})$ with $\ell\ge5$ all having $\ell\mid N$''.
\end{paperlink}

\begin{proposition}[the paper's Proposition prop:ns]\label{prop:galrep:ns}
Let $E/\Q$, with or without complex multiplication, have $E[3]$ irreducible, and let $H$ be the image of $\bar\rho=\bar\rho_{E,3}$. Then $H$ is conjugate to
$\GL_2(\F_3)$, $N_{ns}(\F_3)$ or $N_s(\F_3)$, and $\bar\rho$ is absolutely irreducible on $\Gal(\Qbar/\Q(\sqrt{-3}))$ unless $H\sim N_s(\F_3)$. In that case:
\begin{enumerate}[label=(\roman*)]
\item $\bar\rho(G_{\Q(\sqrt{-3})})=C_4=\langle J\rangle$, and $\bar\rho|_{G_{\Q(\sqrt{-3})}}$ is irreducible over $\F_3$ but not absolutely irreducible;
\item $\bar\rho\otimes\F_9\simeq\Ind_{G_{\Q(\sqrt{-3})}}^{G_\Q}\psi$ for a character $\psi$ of $G_{\Q(\sqrt{-3})}$ of order $4$ with values in $\F_9^\times$;
\item $\bar\rho\simeq\bar\rho\otimes\chi_{-3}$ over $\F_3$, $H^0(\Q,\ad\bar\rho(1))\ne0$, and $h_3(E)=1$ (\Cref{lem:galrep:h0}).
\end{enumerate}
The $3$-division polynomial is a product of two irreducible quadratics in $\Q[x]$ if $H\sim N_s(\F_3)$, and irreducible otherwise.
\end{proposition}

\begin{proof}
(1) By \Cref{prop:galrep:rho-basic}(ii), $\bar\rho(G_{\Q(\sqrt{-3})})=H\cap\SL_2(\F_3)$.

(2) By \Cref{thm:galrep:mod3-images} the irreducible images are $\GL_2(\F_3)$, $N_{ns}(\F_3)$ and $N_s(\F_3)$ (the others lie in Borel subgroups). The proof used
only the determinant and complex conjugation, so it covers curves with complex multiplication.

(3) The determinant-one subgroups are $\SL_2(\F_3)$, $Q_8$ and $C_4=\langle J\rangle$, respectively (\Cref{ch:groups}).

(4) $\SL_2(\F_3)$ and $Q_8$ act absolutely irreducibly on $\F_3^2$. Being non-abelian and acting faithfully is not enough, since a Borel subgroup is
both. The argument: a subgroup of $\GL_2(\Fbar_3)$ fixing a line lies in a Borel subgroup $\sm**0*$, and its commutator subgroup lies in the
unipotent radical $\sm1*01$, abelian of exponent $3$. The commutator subgroup of $Q_8$ is $\{\pm1\}$, which has an element of order $2$, and that of
$\SL_2(\F_3)$ is $Q_8$, non-abelian (\Cref{ch:groups}). So neither fixes a line over $\Fbar_3$.

(5) $C_4$ is abelian. Its generator $J$ has characteristic polynomial $X^2+1$, irreducible over $\F_3$, so $C_4$ fixes no line of $\F_3^2$ but over $\F_9=\F_3[i]$ it
has the two eigenlines with eigenvalues $i$ and $-i$ ($J=m_{-i}$ in the model $\F_3^2=\F_9$; \Cref{ch:groups}). So $\bar\rho|_{G_{\Q(\sqrt{-3})}}$ is irreducible over $\F_3$
and $\bar\rho|_{G_{\Q(\sqrt{-3})}}\otimes\F_9\simeq\psi\oplus\psi^{-1}$, where $\psi$ is the character on one eigenline; $\psi(G_{\Q(\sqrt{-3})})=\langle i\rangle$ has order $4$. This proves (i).

(6) By \Cref{prop:galrep:dihedral} with $K=\Q(\sqrt{-3})$ ((a)$\Rightarrow$(c) and (a)$\Rightarrow$(b)), $\bar\rho\otimes\F_9\simeq\Ind\psi$ (the proof of (a)$\Rightarrow$(c) works over $\F_9$,
where the line is defined), and $\bar\rho\simeq\bar\rho\otimes\chi_{-3}$ over $\F_3$, the isomorphism descending because spaces of homomorphisms commute with
extension of scalars. This proves (ii) and the first part of (iii). Conversely, if $H$ is $\GL_2(\F_3)$ or $N_{ns}(\F_3)$, (4) shows that $\bar\rho|_{G_{\Q(\sqrt{-3})}}$ is
absolutely irreducible.

(7) $\Hom_{G_\Q}(\bar\rho,\bar\rho\otimes\bar\chi_3)\ne0$ since $\bar\chi_3=\chi_{-3}$, and it lies in $\ad\bar\rho(1)$ by the trace-zero statement of \Cref{prop:galrep:dihedral}; so
$H^0(\Q,\ad\bar\rho(1))\ne0$, and $h_3(E)=1$ by the equivalence $h_3(E)=1\iff\bar\rho\simeq\bar\rho\otimes\chi_{-3}$ of \Cref{lem:galrep:h0}(iii), whose proof does
not use the present proposition (only the last equivalence there, with the image $N_s(\F_3)$, does).

(8) The roots of $\psi_3$ correspond to the four lines of $E[3]$ (\Cref{prop:galrep:irreducibility}(ii)). $N_s(\F_3)$ has two orbits of size two on them, and
$N_{ns}(\F_3)$ and $\GL_2(\F_3)$ are transitive (\Cref{thm:galrep:mod3-images}); $\psi_3$ has no rational root as $E[3]$ is irreducible. So $\psi_3$ is a product of two
irreducible quadratics in the first case and irreducible in the other two.
\end{proof}

\begin{example}\label{ex:galrep:ns-examples}
For \curve{726a1}, $\psi_3=(x^2-x-74)(3x^2+8x+20)$; the discriminants $1+4\cdot74=297=3^3\cdot11$ and $8^2-4\cdot3\cdot20=-176=-2^4\cdot11$ (a quadratic $ax^2+bx+c$
has its roots in $\Q(\sqrt{b^2-4ac})$, and square factors can be dropped) give the quadratic fields $\Q(\sqrt{33})$ and
$\Q(\sqrt{-11})$, which with $\Q(\sqrt{-3})$ are the three quadratic subfields of the biquadratic splitting field ($33\cdot(-11)=-3\cdot11^2$). For \curve{6912d1},
$y^2=x^3-6x-8$ (the example of \Cref{ch:galois}), $\psi_3=3(x^2-4x-2)(x^2+4x+6)$, with fields $\Q(\sqrt6)$, $\Q(\sqrt{-2})$. For \curve{338d1},
$\psi_3=(x^2+29x+1858)(3x^2-82x-172)$, with fields $\Q(\sqrt{-39})$, $\Q(\sqrt{13})$. For \curve{11a1}, $\psi_3$ is irreducible with Galois group $S_4$
(\Cref{ex:galrep:eleven}). In the three \texttt{3Ns} cases, $a_p\equiv0\pmod3$ at every prime $p\equiv2\pmod3$ of good reduction: for \curve{726a1}, $a_5=0$,
$a_{17}=6$, $a_{23}=6$, $a_{29}=-6$, $a_{41}=-6$, $a_{47}=6$, $a_{53}=-12$, $a_{59}=-12$.
\end{example}

\begin{paperlink}
The paper's proof of Proposition prop:ns contains steps (1) to (8) in compressed form: ``The first two act absolutely
irreducibly on $\F_3^2$: a subgroup of $\GL_2(\overline\F_3)$ fixing a line lies in a Borel subgroup, so its commutator subgroup is abelian of exponent $3$,
whereas the commutator subgroups $\{\pm1\}$ of $Q_8$ and $Q_8$ of $\SL_2(\F_3)$ are not.'' Being non-abelian and acting faithfully would not suffice, as step (4) notes. The restriction to $\Gal(\Qbar/\Q(\sqrt{-3}))$ in the
\texttt{3Ns} case is irreducible over $\F_3$ and becomes reducible only over $\F_9$, where the induction of (ii) takes place.
\end{paperlink}

\begin{lemma}[the paper's Lemma lem:h0]\label{lem:galrep:h0}
Let $\ell$ be odd and $\bar\rho=\bar\rho_{E,\ell}$ irreducible, $E$ with or without complex multiplication. Let $W^D=\ad T(1)\otimes\Qlmod$, so $W^D[\ell^k]=\ad(T/\ell^k)(1)$.
\begin{enumerate}[label=(\roman*)]
\item $H^0(\Q,W^D)[\ell^k]=H^0(\Q,\ad(T/\ell^k)(1))$, and this group equals $\Hom_{G_\Q}\bigl(T/\ell^k,(T/\ell^k)(1)\bigr)$: every $G_\Q$-map $T/\ell^k\to(T/\ell^k)(1)$ has trace $0$.
\item $H^0(\Q,W^D)$ is finite and cyclic, $H^0(\Q,W^D)\simeq\Z/\ell^{h}$, where $h=h_\ell(E)$ is the largest $k\ge0$ such that
$\rho\bmod\ell^k\simeq(\rho\otimes\chi_\ell)\bmod\ell^k$.
\item $h_\ell(E)=0$ for every $\ell\ge5$, and $h_3(E)\le1$, with $h_3(E)=1$ iff $\bar\rho_{E,3}\simeq\bar\rho_{E,3}\otimes\chi_{-3}$, iff $\bar\rho_{E,3}$ is not absolutely irreducible on
$\Gal(\Qbar/\Q(\sqrt{-3}))$, iff the image of $\bar\rho_{E,3}$ is $N_s(\F_3)$.
\item $h_\ell$ is invariant under quadratic twists and, when $E[\ell]$ is irreducible, under isogenies.
\end{enumerate}
In particular $H^0(\Q,W^D)=0$ whenever $\ell\notin S_f$.
\end{lemma}

\begin{proof}
(i) For $\ell$ odd, $\End(T/\ell^k)=\ad(T/\ell^k)\oplus(\Z/\ell^k)\cdot1$ (a scalar $a$ has trace $2a$, and $2$ is a unit), so $\ad T\otimes\Z/\ell^k=\ad(T/\ell^k)$ and
$W^D[\ell^k]=\ad T(1)\otimes\ell^{-k}\Zl/\Zl\simeq\ad(T/\ell^k)(1)$. Taking $G_\Q$-invariants commutes with taking $\ell^k$-torsion. The $G_\Q$-invariants of $\End(T/\ell^k)(1)$
are the $G_\Q$-maps $\phi\colon T/\ell^k\to(T/\ell^k)(1)$, and the trace is a $G_\Q$-map $\End(T/\ell^k)(1)\to(\Z/\ell^k)(1)=\mu_{\ell^k}$; so $\tr\phi\in H^0(\Q,(\Z/\ell^k)(1))=0$, as $\Q$
contains no nontrivial $\ell$-th root of unity for $\ell$ odd. (Equivalently, as in the paper: $\tr\phi=\chi_\ell(g)\tr\phi$ for every $g$, and $1-\chi_\ell(c)=2$ is a unit.)

(ii) By (i) with $k=1$, $H^0(\Q,W^D)[\ell]=\Hom_{G_\Q}(\bar\rho,\bar\rho\otimes\bar\chi_\ell)$. Both are absolutely irreducible (\Cref{prop:galrep:irreducibility}(i)), so this space is $0$
or, if they are isomorphic, isomorphic to $\End_{G_\Q}(\bar\rho)=\Fl$: its dimension is at most one. A submodule $X$ of $(\Qlmod)^3$ with $\dim_{\Fl}X[\ell]\le1$ is $0$, $\Z/\ell^h$
or $\Qlmod$ (\Cref{ch:modules}: its Pontryagin dual is a quotient of $\Zl^3$ with at most one generator). An element of $X$ of exact order $\ell^k$ is a
$G_\Q$-map $\phi\colon T/\ell^k\to(T/\ell^k)(1)$ not divisible by $\ell$, that is, whose reduction modulo $\ell$ is a nonzero $G_\Q$-map $\bar\rho\to\bar\rho\otimes\bar\chi_\ell$, hence an
isomorphism (both irreducible); then $\phi$ is an isomorphism by Nakayama's lemma, and $\rho\bmod\ell^k\simeq(\rho\otimes\chi_\ell)\bmod\ell^k$. Conversely an isomorphism
$T/\ell^k\to(T/\ell^k)(1)$ is a $G_\Q$-map whose reduction is nonzero, an element of exact order $\ell^k$ by (i). So $X$ has an element of order $\ell^k$ iff
$\rho\bmod\ell^k\simeq(\rho\otimes\chi_\ell)\bmod\ell^k$; that $X$ is finite, hence $\simeq\Z/\ell^h$ with $h$ as stated, follows from (iii).

(iii) If $\phi\colon T/\ell^k\to(T/\ell^k)(1)$ is a $G_\Q$-isomorphism, $\phi\rho(g)\phi^{-1}=\chi_\ell(g)\rho(g)$ modulo $\ell^k$, and taking determinants,
$\det\rho(g)=\chi_\ell(g)^2\det\rho(g)$, so $\chi_\ell(g)^2\equiv1\pmod{\ell^k}$ for every $g\in G_\Q$. The cyclotomic character is onto $\Zl^\times$ (\Cref{ch:galois}), and
$2\in\Zl^\times$; so $4\equiv1\pmod{\ell^k}$, that is, $\ell^k\mid3$. For $\ell\ge5$ this forces $k=0$, and for $\ell=3$ it forces $k\le1$ (\Cref{ch:rings}: $a^2\equiv1\pmod{\ell^k}$
for all units $a$ iff $\ell^k=3$). So $h_\ell(E)=0$ for $\ell\ge5$ and $h_3(E)\le1$, with $h_3(E)=1$ iff $\bar\rho_{E,3}\simeq\bar\rho_{E,3}\otimes\bar\chi_3$, and $\bar\chi_3=\chi_{-3}$
(\Cref{prop:galrep:rho-basic}(ii)). By \Cref{prop:galrep:dihedral} with $K=\Q(\sqrt{-3})$ this is equivalent to $\bar\rho_{E,3}$ not being absolutely irreducible on
$G_{\Q(\sqrt{-3})}$, and by \Cref{prop:galrep:ns} to the image being $N_s(\F_3)$.

(iv) $W^D$ depends only on $\ad T$, which is unchanged by twists (\Cref{prop:galrep:twist-isogeny}(i): conjugation by $\chi_d(g)\rho(g)$ equals conjugation by
$\rho(g)$), and $T$ is unchanged by isogenies when $E[\ell]$ is irreducible (\Cref{prop:galrep:twist-isogeny}(iii)). Finally $\ell\notin S_f$ means that $\bar\rho$ is
absolutely irreducible on $G_{\Q(\sqrt{\ell^*})}$ (and $\ell\nmid2N$); for $\ell=3$ this excludes $h_3(E)=1$ by (iii), and for $\ell\ge5$, $h_\ell(E)=0$ anyway.
\end{proof}

\begin{paperlink}
The paper's Lemma lem:h0 states (i) to (iii): $H^0(\Q,W^D)$ is ``finite and cyclic, of length $h_\ell(E)$, the largest $k$ such that
$\rho\bmod\ell^k\simeq(\rho\otimes\chi_\ell)\bmod\ell^k$'', with the bound $h_3(E)\le1$, the same chain of equivalences for $h_3(E)=1$, and the clause ``This holds for every
$E$, with or without complex multiplication''. Part (iv) is added here. The paper's proof is the one above in compressed form; the determinant argument
of (iii) is what gives finiteness, and it needs no hypothesis on complex multiplication. (For $\ell\nmid N$ the CM case is anyway harmless: complex
multiplication by $\Q(\sqrt{\ell^*})$ forces $\ell\mid N$, \Cref{ch:ellcurves}.) The criterion for $h_\ell(E)>0$ uses $\bar\chi_\ell$, not $\chi_{\ell^*}$; they agree
only for $\ell=3$, and for $\ell\ge5$ the isomorphism $\bar\rho\simeq\bar\rho\otimes\chi_{\ell^*}$ can hold only if $\ell\mid N$ (\Cref{thm:galrep:lstar-automatic}).
\end{paperlink}

\begin{proposition}[the paper's Proposition prop:ns-structure]\label{prop:galrep:ns-structure}
Let $E/\Q$ (complex multiplication allowed) have $\bar\rho=\bar\rho_{E,3}$ with image conjugate to $N_s(\F_3)=\langle J,h\rangle$, and let $K_1$, $K_2$ be the fields
cut out by the two quadratic factors of $\psi_3$, the fixed fields of $\bar\rho^{-1}(C_s)$ and of $\bar\rho^{-1}(\{\pm1,\pm w\})$.
\begin{enumerate}[label=(\roman*)]
\item $K_1$, $K_2$ and $\Q(\sqrt{-3})$ are the three quadratic subfields of the splitting field of $\psi_3$, which is biquadratic.
\item For every prime $p\ne3$, $\bar\rho(I_p)\subseteq C_4=N_s(\F_3)\cap\SL_2(\F_3)=\langle J\rangle$. Consequently every prime $p\ne3$ with $v_p(j)<0$ has $3\mid v_p(j)$ (so
$t_3(E,p)\ge1$); every prime $p\ge5$ of type $\kod{I}_n^*$ has $3\mid n$; no prime $p\ge5$ has type $\kod{II}$ or $\kod{II}^*$; no prime $p\ne3$ has type $\kod{IV}$ or $\kod{IV}^*$.
\item Let $R$ be the set of primes $p\ne3$ ramified in $K_1$. Then $R$ is also the set of primes $p\ne3$ ramified in $K_2$, it equals $\{p\ne3:\bar\rho(I_p)=C_4\}$, and
it is nonempty. For $p\in R$ the reduction is potentially good and $\rho_{E,\ell}(I_p)$ is cyclic of order $4$ for every $\ell\ne p$; if $p\ge5$, $p$ has type
$\kod{III}$ or $\kod{III}^*$. Conversely every prime $p\ge5$ of type $\kod{III}$ or $\kod{III}^*$ lies in $R$.
\item For $p\in R$ and $\ell\ne p$, $\Frob_p$ (arithmetic or geometric) acts on $(\ad V_\ell E)^{I_p}=\Ql\tau$ by $\chi_{-3}(p)$; so $L_p(A_f,s)^{-1}=1-\chi_{-3}(p)p^{-s}$,
$P_p(X)=1-\chi_{-3}(p)pX$ and $e_3(E,p)=v_3(p-\chi_{-3}(p))\ge1$. If $p\ge5$, then $\chi_{-4}(p)=\chi_{-3}(p)$, that is, $p\equiv1$ or $11\pmod{12}$. If $p=2$, then $P_2(X)=1+2X$
and $e_3(E,2)=1$.
\item $a_p(E)\equiv0\pmod3$ for every prime $p\nmid3N$ with $p\equiv2\pmod3$.
\item If $3\nmid N$, $E$ is ordinary at $3$.
\item $h_3(E)=1$.
\end{enumerate}
\end{proposition}

\begin{proof}
Recall $h=\operatorname{diag}(1,-1)$, $w=hJ=\sm0110$, and $N_s(\F_3)=\{\pm1,\pm J,\pm h,\pm w\}$, dihedral of order $8$ with centre $\{\pm1\}$.

(i) The splitting field of $\psi_3$ is the fixed field of the kernel of the action on lines, $\bar\rho^{-1}(\{\pm1\})$, with Galois group $N_s(\F_3)/\{\pm1\}\simeq V_4$; so it
is biquadratic, and its three quadratic subfields are the fixed fields of the three subgroups of index $2$ of $N_s(\F_3)$ containing $-1$: $C_4=\langle J\rangle$,
$C_s=\{\pm1,\pm h\}$ and $\{\pm1,\pm w\}$. The first is the determinant-one subgroup, with fixed field $\Q(\sqrt{-3})$ (\Cref{prop:galrep:rho-basic}(ii)). The
stabiliser in $N_s(\F_3)$ of the line $\langle e_1\rangle$ is $C_s$, and that of $\langle e_1+e_2\rangle$ is $\{\pm1,\pm w\}$ ($J(e_1+e_2)=e_1-e_2$, $h(e_1+e_2)=e_1-e_2$); so the
roots of the factor of $\psi_3$ belonging to the orbit $\{\langle e_1\rangle,\langle e_2\rangle\}$ generate $K_1$, and those of the other factor generate $K_2$.

(ii) For $p\ne3$, $\det\bar\rho=\chi_{-3}$ is unramified at $p$, so $\bar\rho(I_p)\subseteq H\cap\SL_2(\F_3)=C_4$, which has no element of order $3$ or $6$. If $v_p(j)<0$, then on $I_p$,
$\rho_{E,3}(\sigma)=\chi(\sigma)\sm{1}{\nu t_3(\sigma)}{0}{1}$ with $\nu=-v_p(j)$ and $\chi$ trivial or quadratic (\Cref{thm:galrep:tate-shape}, \Cref{prop:galrep:potmult}); on the kernel of
$\chi$ (of index at most $2$) $t_3$ is still onto $\Z_3$, so if $3\nmid\nu$ the image modulo $3$ contains a unipotent of order $3$, which is impossible. So $3\mid\nu$, and
$t_3(E,p)=v_3(\nu)\ge1$; at $p\ge5$ the type $\kod{I}_n^*$ has $n=\nu$ (\Cref{prop:galrep:potmult}). At $p\ge5$ of type $\kod{II}$, $\kod{II}^*$ the image of inertia is
cyclic of order $6$ and injects into $\GL_2(\F_3)$ (\Cref{thm:galrep:serre-tate}), so it contains an element of order $3$ in $C_4$: impossible. At $p\ne3$ of type
$\kod{IV}$, $\kod{IV}^*$, $\bar\rho(I_p)$ has order $3$ (\Cref{lem:galrep:IV}): impossible.

(iii) The subgroups of $C_4$ are $1$, $\{\pm1\}$ and $C_4$; the first two lie in $C_s$ and in $\{\pm1,\pm w\}$, and $C_4$ lies in neither. The prime $p$ ramifies in $K_1$
iff $I_p$ is not contained in the index-two normal subgroup $\bar\rho^{-1}(C_s)$, iff $\bar\rho(I_p)\not\subseteq C_s$, iff $\bar\rho(I_p)=C_4$; likewise for $K_2$. As $C_s\ne C_4$,
$K_1\ne\Q(\sqrt{-3})$; the only quadratic field unramified outside $3$ is $\Q(\sqrt{-3})$ (the discriminant of $\Q(\sqrt d)$, $d$ squarefree, is $d$ if $d\equiv1\pmod4$ and $4d$
otherwise, and its prime factors are the ramified primes; so only $3$ ramifies iff $d=\pm3$ and $d\equiv1\pmod4$, that is, $d=-3$); so some $p\ne3$ ramifies in $K_1$, and $R\ne\emptyset$. Let $p\in R$. If the
reduction at $p$ were potentially multiplicative, $\bar\rho(I_p)$ would consist of elements $\pm u$ with $u$ unipotent, of orders $1$, $2$, $3$ or $6$; but $J$ has
order $4$. So the reduction is potentially good, $\rho_{E,3}(I_p)$ is finite and injects into $\GL_2(\F_3)$ (\Cref{thm:galrep:serre-tate}(ii)): it is cyclic of order $4$,
generated by an element $\tau$ with $\tau\equiv J^{\pm1}$, $\tau^2=-1$. By \Cref{thm:galrep:serre-tate}(i), $\rho_{E,\ell}(I_p)$ is cyclic of order $4$ for every $\ell\ne p$; for
$p\ge5$ this is the semistability defect $4$, that is, $v_p(\Delta_{\min})\in\{3,9\}$, types $\kod{III}$, $\kod{III}^*$ (\Cref{thm:galrep:serre-tate}(iv)). Conversely, at $p\ge5$
of type $\kod{III}$ or $\kod{III}^*$, $\rho_{E,3}(I_p)$ is cyclic of order $4$ and injects modulo $3$, so $\bar\rho(I_p)$ has order $4$ in $C_4$: $\bar\rho(I_p)=C_4$ and $p\in R$.

(iv) Let $p\in R$ and $\ell\ne p$. The commutant of $\tau$ in $\ad V$ is $\Ql\tau$ (\Cref{thm:galrep:additive-shapes}, whose argument uses only that $\tau$ has the
distinct eigenvalues $\pm i$). Frobenius normalises $\rho(I_p)$, so $\rho(\Frob_p)\tau\rho(\Frob_p)^{-1}=\tau^{\pm1}$, and acts on $\Ql\tau$ by that sign. The sign does not depend on $\ell$: it
is the action of $\Frob_p$ by conjugation on the cyclic group $\Gal(L/\Q_p^{\ur})\simeq\rho_{E,\ell}(I_p)$ (\Cref{thm:galrep:serre-tate}(i)); so it can be read at $\ell=3$, and
modulo $3$, where reduction is injective. In $N_s(\F_3)$ the elements of $C_4$ commute with $J$, and those of the other coset $hC_4$ conjugate $J$ to $J^{-1}$
($hJh^{-1}=\sm0{-1}10=J^{-1}$). Since $\bar\rho(I_p)=C_4$, the coset of $\bar\rho(\Frob_p)$ modulo $C_4$ is well defined, and it is $C_4$ iff $\det\bar\rho(\Frob_p)=\chi_{-3}(p)=1$. So
$\Frob_p$, and also $\Frob_p^{-1}$, acts on $(\ad V)^{I_p}$ by $\chi_{-3}(p)$; then $L_p(A_f,s)^{-1}=1-\chi_{-3}(p)p^{-s}$, $P_p(X)=1-\chi_{-3}(p)pX$, and, the naive factor being $1$,
$e_3(E,p)=v_3(1-\chi_{-3}(p)p^{-1})=v_3(p-\chi_{-3}(p))\ge1$ because $p\equiv\chi_{-3}(p)\pmod3$. If $p\ge5$, the image of inertia is tame and the tame relation gives
$\rho(\Frob_p)\tau\rho(\Frob_p)^{-1}=\tau^p=\tau^{\chi_{-4}(p)}$ (as in \Cref{thm:galrep:additive-shapes}); comparing, $\chi_{-4}(p)=\chi_{-3}(p)$: $p\equiv1\pmod{12}$ or
$p\equiv11\pmod{12}$. If $p=2$: $\chi_{-3}(2)=-1$, $P_2(X)=1+2X$ and $e_3(E,2)=v_3(3)=1$.

(v) $\bar\rho\simeq\bar\rho\otimes\chi_{-3}$ (\Cref{prop:galrep:ns}) and \Cref{prop:galrep:dihedral}(d); the primes $p\equiv2\pmod3$ are those inert in $\Q(\sqrt{-3})$.
(vi) \Cref{thm:galrep:lstar-automatic}. (vii) \Cref{lem:galrep:h0}(iii).
\end{proof}

These conclusions were checked on all $146$ curves with image $N_s(\F_3)$ and conductor below $10{,}000$ ($120$ without complex multiplication and $26$ with):
in every case $R$ computed from the discriminants of the two quadratic factors of $\psi_3$ is the same for both factors and agrees with the
primes $p\ne3$ at which the image of inertia on $E[3]$ is cyclic of order $4$ with no fixed point; every multiplicative prime $p\ne3$ has $3\mid n_p$; no
type $\kod{II}$, $\kod{II}^*$ occurs at $p\ge5$ and no type $\kod{IV}$, $\kod{IV}^*$ at $p\ne3$; $a_p\equiv0\pmod3$ for $p\equiv2\pmod3$, $p<410$, $p\nmid3N$; and $3\nmid a_3$ whenever
$3\nmid N$. Of the $120$ non-CM curves, $100$ have primes $p\ge5$ in $R$, all of type $\kod{III}$ or $\kod{III}^*$ and $\equiv1$ or $11\pmod{12}$; for the other $20$ (the
classes \curve{3840h}, \curve{3840k}, \curve{3840o}, \curve{3840r}, \curve{4864b}, \curve{4864c}, \curve{4864o}, \curve{4864p}, \curve{6912d}, \curve{6912f}, \curve{6912j},
\curve{6912l}, \curve{6912p}, \curve{6912r}, \curve{6912v}, \curve{6912x}), $R=\{2\}$, and the type at $2$ is $\kod{III}$ or $\kod{III}^*$ with $f_2=8$. PARI's Euler factors are
$1+11X$ for \curve{726a1} at $11$, $1-13X$ for \curve{338d1} at $13$, and $1+2X$ for \curve{3840h1} and \curve{6912d1} at $2$, as (iv) predicts. Of the $26$
curves with complex multiplication, $14$ have $R=\{11\}$ and $12$ have $R=\{2\}$.

For the optimal curves (the curve of Cremona's file \texttt{degphi} in each class, which is the first curve of the class except for \curve{990h3}) we
recomputed the counts: below conductor $10{,}000$ there are $90$ optimal curves without complex multiplication with image $N_s(\F_3)$, of which $44$ have
$3\mid N$, $74$ have a prime $p\ge5$ in $R$ and $16$ have $R=\{2\}$; and $13$ optimal curves with complex multiplication with this image, namely \curve{121b1},
\curve{256a1}, \curve{256d1}, \curve{1089e1}, \curve{1936f1}, \curve{2304c1}, \curve{2304k1}, \curve{3025a1}, \curve{5929a1}, \curve{6400f1}, \curve{6400r1},
\curve{7744a1} and \curve{7744s1}, of which $3$ have $3\mid N$, $7$ (those with $j=-32768$) have $R=\{11\}$ and $6$ (those with $j=8000$) have $R=\{2\}$.

\begin{paperlink}
Section sec:anomaly: ``Conversely every $\Ns$ curve of the tables has a prime of type $\mathrm{III}$ or $\mathrm{III}^*$ with a positive term $e_3$''; ``The
residue classes of the type $\mathrm{III}$ primes of the violating curves are $q\equiv1$ and $q\equiv11\pmod{12}$ in comparable numbers, with some
$q=2$''. \Cref{prop:galrep:ns-structure} proves these observations, except that the Kodaira type at $p=2$ of a prime of $R$ is not determined by the
argument: the paper says so (``That the type at $2$ is always $\mathrm{III}$ or $\mathrm{III}^*$ when $R=\{2\}$ is not proved''), and the data show $\kod{III}$ or
$\kod{III}^*$ in every case. What is proved at $2$ is the cyclic image of order $4$ and the Euler factor $1+2X$. The paper's ``the violation involves a prime
$p\equiv-1\pmod3$ with $t_3(E,p)\ge1$'', recorded there as an observation, is consistent with (ii): every potentially multiplicative prime $p\ne3$ of a \texttt{3Ns} curve has $t_3(E,p)\ge1$.

The paragraph after the proof of Proposition prop:ns-structure says: ``Among the $90$ optimal curves without complex multiplication with image
$\Ns$ of conductor below $10{,}000$ ($44$ of them with $3\mid N$), $74$ have a prime $p\ge5$ in $R$, of type $\mathrm{III}$ or $\mathrm{III}^*$ and $p\equiv1$ or
$11\pmod{12}$, and $16$ have $R=\{2\}$''; it then lists the $13$ optimal curves with complex multiplication and image $N_s(\F_3)$, $7$ with $R=\{11\}$ and $6$
with $R=\{2\}$. Our recomputation above agrees with all these counts.
\end{paperlink}

\begin{paperlink}
The paper's Proposition prop:ns-structure has the content of \Cref{prop:galrep:ns-structure}, with the parts in another order: its (iv) is (i)
here; its (i) is (ii) here, where the clause on type $\kod{I}_n^*$ at $p\ge5$ is the case $v_p(j)=-n$ of the clause on $v_p(j)<0$; its (ii) and (iii) are (iii) and
(iv) here; and its (v) to (vii) are (v) to (vii) here. The paper's proof takes its (iv) second, after its (i): ``We prove (iv) next, because (ii) uses the fields $K_1$ and $K_2$ it
describes.'' Putting that part first here reflects the same dependence.
\end{paperlink}

\begin{corollary}[$S_f$ for elliptic curves]\label{cor:galrep:Sf}
Let $\ell$ be odd and $\ell\nmid N$. The following are equivalent: (a) $\ell\notin S_f$, that is, $\bar\rho_{E,\ell}$ is absolutely irreducible on $\Gal(\Qbar/\Q(\sqrt{\ell^*}))$;
(b) $E[\ell]$ is irreducible and $h_\ell(E)=0$; (c) $E[\ell]$ is irreducible and, if $\ell=3$, the image of $\bar\rho_{E,3}$ is not $N_s(\F_3)$. So for $\ell\nmid2N$ the hypothesis of
Theorem thm:main reads: $E[\ell]$ is irreducible, and not ($\ell=3$ and image \texttt{3Ns}).
\end{corollary}

\begin{proof}
By definition \cite[(31)]{DFG04}, $S_f$ consists of the primes dividing $2N$ and those $\ell$ for which $\bar\rho|_{G_{\Q(\sqrt{\ell^*})}}$ is not absolutely irreducible. If $E[\ell]$
is reducible, so is its restriction, and $\ell\in S_f$; so all three conditions require irreducibility. Given it, for $\ell\ge5$ (a) holds by
\Cref{thm:galrep:lstar-automatic} and (b), (c) hold by \Cref{lem:galrep:h0}(iii); for $\ell=3$ all three are equivalent to the image not being $N_s(\F_3)$, by
\Cref{prop:galrep:ns} and \Cref{lem:galrep:h0}(iii).
\end{proof}

Below conductor $10{,}000$ the curves with $3\nmid N$, $E[3]$ irreducible and $3\in S_f$ are the $62$ curves with label \texttt{3Ns} and $3\nmid N$ and the
$20$ curves with complex multiplication, image $N_s(\F_3)$ and $3\nmid N$ (with $j=-32768$ or $8000$).

\section{The Eisenstein case}\label{sec:galrep:eisenstein}

When $E[\ell]$ is reducible, the Hecke maximal ideal attached to $(f,\ell)$ is Eisenstein, multiplicity one and the theorems of Diamond, Flach and
Guo are unavailable, and the modular degree can lose the local divisibility of Theorem thm:main. The paper studies this case in Section
sec:eis and Conjecture conj:eis, where the controlling invariant is rational torsion. This section proves the structure of a reducible $\bar\rho$
at a prime $\ell\nmid N$ and the dictionary between the reducible mod $3$ images and torsion.

\begin{proposition}[reducible $\bar\rho$]\label{prop:galrep:reducible}
Let $\ell$ be odd, $\ell\nmid N$, and $E[\ell]$ reducible, with a stable line on which $G_\Q$ acts by $\eta_1$ and quotient character $\eta_2$.
\begin{enumerate}[label=(\roman*)]
\item $\bar\rho^{\mathrm{ss}}=\eta_1\oplus\eta_2$ and $\eta_1\eta_2=\bar\chi_\ell$.
\item $E$ is ordinary at $\ell$, and $\{\eta_1|_{I_\ell},\eta_2|_{I_\ell}\}=\{\omega,1\}$.
\item $\eta_1$ and $\eta_2$ are unramified at every prime of multiplicative reduction.
\item If $E$ is semistable, then $\{\eta_1,\eta_2\}=\{1,\bar\chi_\ell\}$, so $a_p\equiv1+p\pmod\ell$ for every $p\nmid N\ell$, and $E$ or a curve $\ell$-isogenous to it over $\Q$ has a
rational point of order $\ell$.
\end{enumerate}
\end{proposition}

\begin{proof}
(i) The composition factors of $\bar\rho$ are $\eta_1$ and $\eta_2$, and $\eta_1\eta_2=\det\bar\rho=\bar\chi_\ell$ (\Cref{prop:galrep:rho-basic}(ii)). (ii) is
\Cref{cor:galrep:ordinary}(i). (iii) At a multiplicative prime $q$, $\bar\rho(I_q)$ consists of unipotent elements (\Cref{thm:galrep:tate-shape}(i)); a unipotent element
has the single eigenvalue $1$, so it acts trivially on every one-dimensional subquotient, in particular on the line and on the quotient. (iv) If $E$
is semistable, $\eta_1$ and $\eta_2$ are unramified at every $p\ne\ell$: at good primes by \Cref{thm:galrep:nos}, at multiplicative primes by (iii). A character
of $G_\Q$ unramified outside $\ell$ (and $\infty$) factors through $\Gal(\Q(\zeta_{\ell^\infty})/\Q)\simeq\Zl^\times$ by the theorem of Kronecker and Weber (\Cref{ch:ant}); if it
takes values in $\Fbar_\ell^\times$, its order is prime to $\ell$, so it factors through the quotient $\Gal(\Q(\zeta_\ell)/\Q)\simeq\Fl^\times$ and is a power $\bar\chi_\ell^a$. As
$\Q(\zeta_\ell)/\Q$ is totally ramified at $\ell$, $\bar\chi_\ell^a$ is determined by its restriction $\omega^a$ to $I_\ell$; by (ii) the two characters are $\bar\chi_\ell$ and $1$.
Then $a_p\equiv\tr\bar\rho(\Frob_p)=\eta_1(\Frob_p)+\eta_2(\Frob_p)=p+1\pmod\ell$. If $\eta_1=1$ the line consists of rational points: a rational point of order $\ell$. If
$\eta_1=\bar\chi_\ell$, the line $C$ is isomorphic to $\mu_\ell$, and on $E'=E/C$ the image of $E[\ell]$ is $E[\ell]/C$, a subgroup of order $\ell$
on which $G_\Q$ acts by $\eta_2=1$: a rational point of order $\ell$ on the $\ell$-isogenous curve $E'$.
\end{proof}

We checked (iv) on every curve of conductor below $10{,}000$ with square-free $N$, $\ell\nmid N$, and a Borel or split Cartan label at $\ell=3$, $5$, $7$ ($1{,}315$,
$158$ and $26$ curves): $a_p\equiv1+p\pmod\ell$ for every prime $p<282$ with $p\nmid N\ell$.

\begin{proposition}[the reducible mod $3$ images and torsion]\label{prop:galrep:eis-labels}
\begin{enumerate}[label=(\roman*)]
\item The mod $3$ image of $E$ is \texttt{3B.1.1} if and only if $E$ has a rational point of order $3$ and $E[3]$ is not split; it is \texttt{3Cs.1.1} iff
$E[3]\simeq\Z/3\oplus\mu_3$; and it is \texttt{3B.1.2} iff $\mu_3\subset E[3]$ with quotient $\Z/3$ and $E[3]$ not split. The twist by $\chi_{-3}$ exchanges \texttt{3B.1.1} and \texttt{3B.1.2}.
\item $3\mid\#E(\Q)_{\mathrm{tors}}$ iff the image is \texttt{3B.1.1} or \texttt{3Cs.1.1}.
\item If $E$ has no complex multiplication and $E[3]\simeq\Z/3\oplus\mu_3$, then $E(\Q)$ has no point of order $9$. Hence, for $E$ without complex multiplication,
$v_3(\#E(\Q)_{\mathrm{tors}})+[E[3]\simeq\Z/3\oplus\mu_3]\le2$.
\item At $\ell=5$: the image \texttt{5Cs.1.1} means $E[5]\simeq\Z/5\oplus\mu_5$ (\curve{11a1}), and \texttt{5B.1.1} a rational point of order $5$ with $E[5]$ non-split (\curve{11a3}).
\end{enumerate}
\end{proposition}

\begin{proof}
(i) is the table of \Cref{thm:galrep:mod3-images}, read through the arithmetic meanings given in its proof, and \Cref{prop:galrep:twist-isogeny}(ii). (ii) A
rational point of order $3$ spans a line on which $G_\Q$ acts trivially, so $H$ fixes a nonzero vector. Of the eight images, \texttt{3B.1.1} and \texttt{3Cs.1.1} fix
$e_1$; \texttt{3B}, \texttt{3Cs}, \texttt{3Ns}, \texttt{3Nn} and $\GL_2(\F_3)$ contain $-1$, which fixes no nonzero vector; and \texttt{3B.1.2} $=\{\sm{a}{b}{0}{1}\}$ contains
$\operatorname{diag}(-1,1)$, whose fixed vectors are the multiples of $e_2$, and $u$, which moves $e_2$. (iii) This is proved in \Cref{ch:ellcurves}, and is the paper's
Proposition prop:no9: if $P$ had order $9$, then with $\phi\colon E\to E'=E/\mu_3$ the group $C=\hat\phi^{-1}(\langle P\rangle)\subset E'$ would be a $G_\Q$-stable cyclic subgroup
of order $27$ (it does not contain $E'[3]$, since $\hat\phi(E'[3])=\mu_3\not\subseteq\langle P\rangle$), and a rational cyclic $27$-isogeny forces complex multiplication. The
inequality follows from Mazur's theorem on torsion \cite{Mazur77}: $v_3(\#E(\Q)_{\mathrm{tors}})\le2$, with equality only for $\Z/9$. (iv) is read from Cremona's
data as in \Cref{conv:galrep:labels} (and \Cref{ex:galrep:eleven}).
\end{proof}

Below conductor $10{,}000$ the curves with $3\mid\#E(\Q)_{\mathrm{tors}}$ are exactly the $2{,}811$ curves with image \texttt{3B.1.1} (torsion orders $3$, $6$, $9$, $12$)
and the $264$ with image \texttt{3Cs.1.1} (torsion orders $3$ and $6$ only, as (iii) requires).

\begin{paperlink}
Section sec:eis describes \texttt{3Cs.1.1} as ``the subgroup of order $2$ of the split Cartan subgroup that fixes a
point'', lists \curve{14a1} (torsion $\Z/6$, image \texttt{3Cs.1.1}) among the curves with $D=2$, and proves in Proposition prop:no9 that the two sources of a
deficit $2$ never combine, so that the bound $D\le2$ of Conjecture conj:eis follows from its first bound: this is \Cref{prop:galrep:eis-labels}(iii).
\end{paperlink}

\begin{remark}[what reducibility costs]\label{rem:galrep:eisenstein}
If $E[\ell]$ is reducible, the lattices $T_\ell E'$ of the curves $E'$ in the isogeny class are not all homothetic (\Cref{ch:reps}: the stable lattices of a
reducible residual representation), so $\ad T$ and $h_\ell$ may change along the class, and the choice of the optimal curve matters at $\ell$. The
maximal ideal of the Hecke algebra attached to $(f,\ell)$ is Eisenstein (\Cref{ch:moddeg}), multiplicity one and the theorems of Diamond, Flach and
Guo are unavailable, and the modular degree can fall short of the local bound: the deficit $D$ of Conjecture conj:eis. This is why the
hypothesis ``$E[\ell]$ irreducible'' appears in Theorem thm:quant, Conjecture conj:euler and Proposition prop:forms.
\end{remark}

\section[Serre's conjecture and the optimal level]{Serre's conjecture for \texorpdfstring{$\bar\rho_{E,\ell}$}{rhobar} and the optimal level}\label{sec:galrep:serre}

Theorem thm:qual lowers the level of $f_E$, or of a twist, to the Serre conductor, and then twists back. This section states the optimal-level
theorem in the form the paper uses and lists the elliptic-curve inputs, all proved above; the general statement and the passage to trivial
nebentypus are in \Cref{ch:modgalrep}.

\begin{theorem}[Serre's conjecture; the optimal level]\label{thm:galrep:optimal-level}
\begin{enumerate}[label=(\roman*)]
\item (Khare and Wintenberger) Every odd irreducible continuous $\bar\sigma\colon G_\Q\to\GL_2(\Fbar_\ell)$ arises from a newform of weight $k(\bar\sigma)$, level $N(\bar\sigma)$ and
character determined by $\det\bar\sigma$.
\item (Ribet, Carayol, Diamond; \cite[Theorem 3.15]{DDT} as quoted by the paper) Let $\ell$ be odd and let $E'/\Q$ have $\ell\nmid N(E')$ and $\bar\rho=\bar\rho_{E',\ell}$
irreducible, and, if $\ell=3$, image not $N_s(\F_3)$. Then there are a newform $g$ of weight $2$ on $\Gamma_0(N(\bar\rho))$ and a prime $\lambda\mid\ell$ of its coefficient
field with $\bar\rho_{g,\lambda}\simeq\bar\rho\otimes\Fbar_\ell$; in particular $a_p(g)\equiv a_p(E')\pmod\lambda$ for all $p\nmid\ell N(E')$.
\end{enumerate}
\end{theorem}

We do not prove this. Part (i) is \cite{KhareWintenberger}, with Serre's formulation in \cite{Serre87}. Part (ii) is the level-lowering theorem of
Ribet \cite{Ribet90}, Carayol \cite{Carayol89} and Diamond \cite{Diamond95}, which needs only the modularity of $\bar\rho$; \Cref{ch:modgalrep} states it in general and
proves the specialisation to $\bar\rho_{E',\ell}$, including the triviality of the nebentypus. The inputs from elliptic curves are: $\bar\rho$ is modular of level $N(E')$
(modularity, \cite{Wiles95,BCDT}; \Cref{ch:moddeg}); it is odd, $\det\bar\rho(c)=-1$ (\Cref{prop:galrep:rho-basic}(iii)); it is absolutely irreducible, and absolutely
irreducible on $\Gal(\Qbar/\Q(\sqrt{-3}))$ when $\ell=3$ (\Cref{prop:galrep:irreducibility}(i), \Cref{prop:galrep:ns}); its restriction to $D_\ell$ is finite flat, so $\delta(\bar\rho)=0$
(\Cref{thm:galrep:good-at-l}(i)); its weight is $2$ (\Cref{thm:galrep:good-at-l}(iii)); $\det\bar\rho=\bar\chi_\ell$ (\Cref{prop:galrep:rho-basic}(ii)); and $N(\bar\rho)$ is given by
\Cref{prop:galrep:serre-conductor-E}.

\begin{paperlink}
The optimal-level paragraph of Section 3: ``For $\bar\sigma=\bar\rho_{E',\ell}$ with $E'$ an elliptic curve of conductor prime to $\ell$, the restriction to $D_\ell$ comes
from the finite flat group scheme $E'[\ell]$ and is good, and $\det\bar\sigma=\bar\chi_\ell$ forces the nebentypus, which is trivial modulo $\lambda$ and of order
prime to $\ell$, to be trivial.'' The first clause is \Cref{thm:galrep:good-at-l}(i), the determinant is \Cref{prop:galrep:rho-basic}(ii), and the nebentypus
argument is \Cref{ch:modgalrep}'s. Theorem thm:qual states the extra hypothesis at $\ell=3$ as the paper's ``absolutely irreducible when restricted to
$\Gal(\overline\Q/\Q(\sqrt{-3}))$''; by \Cref{prop:galrep:ns} this is equivalent to the image not being
$N_s(\F_3)$, the form in which \Cref{thm:galrep:optimal-level}(ii) states it.
\end{paperlink}

\begin{example}[the paper's Example, read through this chapter]\label{ex:galrep:level-lowering}
The congruences of the Example of Section 3 of the paper (case (a) at $q=2$, case (b), case (c) at $\ell=3$ and at $\ell=5$), together with one case (a)
at $q=11$ that the paper does not list and the book adds as an illustration (the curve \curve{121c1} of \Cref{ex:galrep:serre-conductor-examples}),
with the Serre conductors of \Cref{ex:galrep:serre-conductor-examples} and the modular degrees from Cremona's file \texttt{degphi}:
\begin{center}\footnotesize\setlength{\tabcolsep}{3pt}
\begin{tabular}{lllll}
\toprule
case & curve, $\degphi$ & type at $q$ & level lowering & congruent form\\
\midrule
(a), $q=11$ (added) & \curve{121c1}, $6$ & $\kod{IV}$ at $11$ & $N(\bar\rho_3)=11$ & $f_{\curve{11a}}$\\
(a), $q=2$ & \curve{52a1}, $3$ & $\kod{IV}^*$ at $2$ & $N(\bar\rho_3)=26$ & $f_{\curve{26b}}$\\
(b) & \curve{121a1}, $6$ & $\kod{II}$ at $11$ & twist by $\chi_{-11}$: \curve{121c2}, $N(\bar\rho_3)=11$ & $f_{\curve{11a}}$, twisted\\
(c), $\ell=3$ & \curve{539c1}, $2^5\cdot3^2$ & $\kod{I}_3^*$ at $7$ & twist by $\chi_{-7}$: \curve{77c1}, $N(\bar\rho_3)=11$ & $f_{\curve{11a}}$, twisted\\
(c), $\ell=5$ & \curve{980i1}, $2^6\cdot3^2\cdot5$ & $\kod{I}_5^*$ at $7$ & twist by $\chi_{-7}$: \curve{140b1}, $N(\bar\rho_5)=4$ & $f_{\curve{20a}}$, twisted\\
\bottomrule
\end{tabular}
\end{center}
The paper identifies the twisted forms: $f_{\curve{11a}}\otimes\chi_{-11}=f_{\curve{121d}}$, $f_{\curve{11a}}\otimes\chi_{-7}=f_{\curve{539d}}$ and
$f_{\curve{20a}}\otimes\chi_{-7}=f_{\curve{980g}}$ (recomputed: the twists of \curve{11a1} and \curve{20a1} have the traces of Frobenius of \curve{121d1},
\curve{539d1}, \curve{980g1}). In case (a) at $q=11$, $3\mid\deg\varphi_{\curve{121c1}}=6$; at $q=2$, $3=\deg\varphi_{\curve{52a1}}$, where the bound $e_3(E,2)=1$ of \Cref{cor:galrep:IV-at-2} is attained. In
case (b) the Serre exponent $2$ of $\bar\rho_{\curve{121a1},3}$ at $11$ drops to $1$ after the twist (\Cref{prop:galrep:mod3-inertia}). In case (c) the type $\kod{I}_3^*$ has
$3\mid n=3$, so the twist $\bar\rho\otimes\chi_{-7}$ is unramified at $7$ (\Cref{lem:galrep:tamagawa-exp} for \curve{77c1}, type $\kod{I}_3$) and the level drops by $7^2$. The
congruences of \Cref{ex:galrep:serre-conductor-examples} were checked for all $p<2000$ ($p<1224$ at $\ell=5$) away from the bad primes; Sturm's bound
turns the paper's versions into identities for all coefficients (\Cref{ch:modforms}). At $\ell=5$, $5\mid980$, so this case illustrates the congruence, not
Theorem thm:qual.
\end{example}

\section{Exercises}\label{sec:galrep:exercises}

\begin{exercise}\label{ex:galrep:complex-conjugation}
Let $\ell$ be odd. (a) Show directly that $\rho_{E,\ell}(c)$ is conjugate in $\GL_2(\Zl)$ to $\operatorname{diag}(1,-1)$ for a complex conjugation $c$, and deduce that $\bar\rho_{E,\ell}$
is irreducible iff it is absolutely irreducible. (b) Where does the argument fail for $\ell=2$? Show that $\bar\rho_{E,2}(c)$ is either trivial or a transvection,
according as $E(\R)$ has two components or one (consider the real $2$-torsion points).

\emph{Hint.} (b) The idempotents $\frac12(1\pm\rho(c))$ need $2$ to be invertible, and modulo $2$ the eigenvalues $1$ and $-1$ coincide. In a model $y^2=g(x)$,
$E(\R)$ has two components iff $g$ has three real roots; then $c$ fixes $E[2]$ pointwise, and otherwise it fixes the one real point of order $2$ and swaps
the other two, which on $\F_2^2$ is a transvection.
\end{exercise}

\begin{exercise}\label{ex:galrep:eleven-five}
(The paper's Lemma lem:tamagawa-exp on \curve{11a1} and \curve{11a3}.) Both curves have multiplicative reduction at $11$, with $n_{11}=5$, resp.\ $1$. Show that
$\bar\rho_{\curve{11a1},5}$ is unramified at $11$ and $\bar\rho_{\curve{11a3},5}$ is ramified there, compute $\#E[5^t]^{I_{11}}$ for both curves and every $t\ge1$, and explain why
$E[5]\simeq\Z/5\oplus\mu_5$ for \curve{11a1} is consistent with the first statement.

\emph{Answer.} By \Cref{lem:galrep:tamagawa-exp}(ii), $\#E[5^t]^{I_{11}}=5^{t+1}$ for \curve{11a1} and $5^t$ for \curve{11a3}. The module $\Z/5\oplus\mu_5$ is unramified at every
prime other than $5$.
\end{exercise}

\begin{exercise}\label{ex:galrep:torsion-ur}
Let $K/\Q_p$ be finite, $\ell\ne p$, and $E/K$. Using the filtration $E_1\subseteq E_0\subseteq E$ over the finite unramified extensions of $K$, show: (a) if the
reduction is additive, $E(K^{\ur})[\ell^\infty]$ is finite, of order dividing $\#\Phi_E(\bar k)$; (b) if it is multiplicative, $E(K^{\ur})[\ell^\infty]$ is an extension of a
finite group by $\mu_{\ell^\infty}$, so it is infinite; (c) if it is good, $E(K^{\ur})[\ell^\infty]\simeq(\Qlmod)^2$. Deduce once more the criterion of N\'eron, Ogg and
Shafarevich, and the dichotomy of \Cref{thm:galrep:nos}(e), (f) for $E$ with bad reduction over $K$ itself.
\end{exercise}

\begin{exercise}\label{ex:galrep:tamagawa-l-part}
Let $q\ne\ell$ and let $E$ have additive reduction at $q$. (a) Deduce from \Cref{lem:galrep:component-invariants} that the $\ell$-part of $c_q$ is
$\#E(\Q_q)[\ell^\infty]$. (b) Show by an example that (a) fails for split multiplicative reduction (take $q\equiv1\pmod\ell$ and $\ell\nmid n_q$). (c) Compare with the final
check in the proof of the paper's Lemma lem:tamfactor, where the $\ell$-part of $c_q$ at a multiplicative prime is read from the quotient
$\Phig_q[\ell^\infty]$ of $E[\ell^\infty]^{I_q}$, not from the sub $\mu_{\ell^\infty}\otimes\varepsilon$.
\end{exercise}

\begin{exercise}\label{ex:galrep:121c1}
(The paper's Lemma lem:IV for \curve{121c1}.) The curve \curve{121c1} has type $\kod{IV}$ at $11$. Show that $N(\bar\rho_{\curve{121c1},3})=11$ and check
$a_p(\curve{121c1})\equiv a_p(\curve{11a1})\pmod3$ for all primes $p<100$, $p\ne11$, including $p=3$, where both curves are ordinary. Which statement of
\Cref{thm:galrep:deligne-fontaine} does the case $p=3$ illustrate?

\emph{Answer.} The exponent at $11$ is $f_{11}-\dim\Phig_{11}[3]=2-1=1$. At $p=3$ one has $a_3(\curve{121c1})=2$ and $a_3(\curve{11a1})=-1$; the congruence is part (iii) of
\Cref{thm:galrep:deligne-fontaine}, since \Cref{prop:galrep:irreducibility}(iii) only covers primes not dividing $3N$.
\end{exercise}

\begin{exercise}\label{ex:galrep:IV-at-2}
(The paper's Remark rem:wild on the type $\kod{IV}$ at $2$.) Prove that a curve of type $\kod{IV}$ or $\kod{IV}^*$ at $2$ has $P_2(X)=1+2X$ and $e_3(E,2)=1$ using only
\Cref{lem:galrep:component-invariants} at $q=2$, \Cref{thm:galrep:serre-tate}(i), (ii) and the tame relation; do not use Kraus's classification. Check the
result on \curve{20a1} and \curve{980i1} (type $\kod{IV}^*$ at $2$) with \texttt{lfuneuler(lfunsympow(E,2),2)}.
\end{exercise}

\begin{exercise}\label{ex:galrep:kraus}
Carry out \Cref{rem:galrep:kraus-from-components}: for $q\ge5$ derive the orders $\kod{I}_0^*\mapsto2$, $\kod{IV},\kod{IV}^*\mapsto3$, $\kod{III},\kod{III}^*\mapsto4$, $\kod{II},\kod{II}^*\mapsto6$
of the image of inertia from the component groups, \Cref{lem:galrep:component-invariants} at $\ell=2$ and $\ell=3$, and the independence of $\ell$. In particular
prove that an element $\tau\in\SL_2(\Z_2)$ with $\tau^2=-1$ is not congruent to $1$ modulo $2$, and that its reduction fixes a line of $\F_2^2$.

\emph{Hint.} The reduction $\bar\tau$ satisfies $\bar\tau^2=1$, so its characteristic polynomial is $X^2+1=(X+1)^2$ over $\F_2$; an element of $\GL_2(\F_2)$ with this
characteristic polynomial is $1$ or a transvection, and it is not $1$ by the argument of \Cref{rem:galrep:kraus-from-components}.
\end{exercise}

\begin{exercise}\label{ex:galrep:types-rows}
(The paper's Remark rem:types.) Verify the table of \Cref{prop:galrep:mod3-inertia} row by row on the curves \curve{121a1} ($\kod{II}$), \curve{121c2} ($\kod{IV}^*$),
\curve{726a1} ($\kod{III}$ at $11$), \curve{539d1} ($\kod{I}_0^*$ at $7$) and \curve{539c1} ($\kod{I}_3^*$ at $7$), by computing the ramification indices of the prime in the fields of
definition of the $3$-torsion points. Which statements of the remark concern $\bar\rho_{E,3}$ and which concern $E$, and in which row of the table
do the twist's effects on the two differ?
\end{exercise}

\begin{exercise}\label{ex:galrep:cns}
Let $\ell$ be odd. Show that $\bar\rho_{E,\ell}(G_\Q)$ is never contained in $\SL_2(\Fl)$ and never conjugate to a subgroup of $C_{ns}(\Fl)$. Show that it can be contained in
$C_s(\Fl)$, and characterise this case by the splitting of $E[\ell]$.

\emph{Hint.} Use the determinant and complex conjugation (\Cref{prop:galrep:image-constraints}). The image lies in a conjugate of $C_s(\Fl)$ iff $E[\ell]$ is a direct
sum of two $G_\Q$-stable lines; \curve{14a1} at $\ell=3$ is an example.
\end{exercise}

\begin{exercise}\label{ex:galrep:726a1}
(The paper's Proposition prop:ns for \curve{726a1}.) Factor $\psi_3$ for \curve{726a1}, identify the three quadratic subfields $\Q(\sqrt{33})$, $\Q(\sqrt{-11})$, $\Q(\sqrt{-3})$ of
its splitting field, deduce that the image is $N_s(\F_3)$, and check $a_p\equiv0\pmod3$ for the primes $p\equiv2\pmod3$, $p<60$, $p\ne2,11$.
\end{exercise}

\begin{exercise}\label{ex:galrep:h3-traces}
(The paper's Section sec:anomaly.) Let $E$ have image \texttt{3Ns}. (a) Show $h_3(E)=1$ using \Cref{lem:galrep:h0} only. (b) Show without determinants that
$\rho\bmod9\simeq(\rho\otimes\chi_3^{\mathrm{cyc}})\bmod9$ would give $a_p\equiv pa_p\pmod9$ for $p\nmid3N$, hence $9\mid a_p$ for every such $p\equiv2\pmod3$. (c) For \curve{726a1}, $a_{17}=6$: so the trace
test of (b) also excludes $h_3\ge2$ for this curve. Why is such a test unnecessary?

\emph{Hint.} (b) Take traces at $\Frob_p$: $\chi_3^{\mathrm{cyc}}(\Frob_p)=p$, and $1-p$ is a unit modulo $3$ when $p\equiv2\pmod3$. (c) $9\nmid6$; the test is unnecessary
because the determinant argument in the proof of \Cref{lem:galrep:h0}(iii) gives $h_3(E)\le1$ for every curve.
\end{exercise}

\begin{exercise}\label{ex:galrep:R-examples}
Prove parts (iii) and (iv) of \Cref{prop:galrep:ns-structure} in the case of \curve{726a1} ($R=\{11\}$), \curve{338d1} ($R=\{13\}$) and \curve{6912d1} ($R=\{2\}$): find $R$
from the discriminants of the quadratic factors of $\psi_3$, check the Kodaira types at the primes of $R$, the classes modulo $12$, and the Euler
factors $1+11X$, $1-13X$, $1+2X$ returned by \texttt{lfuneuler(lfunsympow(E,2),p)}.
\end{exercise}

\begin{exercise}\label{ex:galrep:lstar}
Let $\ell\ge5$ and $\ell\nmid N$. (a) Using the fundamental characters, show that $\bar\rho_{E,\ell}\not\simeq\bar\rho_{E,\ell}\otimes\chi_{\ell^*}$, without appealing to
\Cref{cor:galrep:ordinary} (redo its proof for the ordinary and the supersingular case separately). (b) The curve \curve{608b1} ($N=2^5\cdot19$) has mod $5$
image \texttt{5Ns}. Find the quadratic field $K$ with $\bar\rho_{\curve{608b1},5}\simeq\bar\rho_{\curve{608b1},5}\otimes\chi_K$ by testing $a_p\equiv0\pmod5$ at the primes inert in the
quadratic fields of discriminant dividing $4N$, and check that $K\ne\Q(\sqrt5)$.
\end{exercise}

\begin{exercise}\label{ex:galrep:eisenstein}
Let $E$ be semistable, $\ell$ odd, $\ell\nmid N$, and $E[\ell]$ reducible. Prove that $a_p\equiv1+p\pmod\ell$ for every $p\nmid N\ell$ and that $E$ or an $\ell$-isogenous curve has a
rational point of order $\ell$. Check both statements for \curve{14a1} and \curve{26a1} at $\ell=3$, and find the isogenous curve with a point of order $3$ for \curve{14a3}.

\emph{Answer for the last part.} \curve{14a3} has image \texttt{3B.1.2}, its stable line is $\mu_3$, and its torsion group has order $2$. Its only
$3$-isogenous neighbour is \curve{14a1} (the isogeny matrix in \texttt{allisog}), whose torsion group is $\Z/6$.
\end{exercise}

\begin{exercise}\label{ex:galrep:3cs11}
Determine the mod $3$ images of \curve{14a1} and of \curve{539a2} from $\psi_3$ and the characters of the rational lines. Explain why \texttt{3Cs.1.1}, the image of
\curve{14a1}, is not the full split Cartan subgroup, and which curve of the pair has the full split Cartan image.
\end{exercise}

\begin{exercise}\label{ex:galrep:conductors}
Compute $N(\bar\rho_{E,3})$ from \Cref{prop:galrep:serre-conductor-E} for \curve{539c1}, \curve{77c1}, \curve{726a1} and \curve{6912d1}, and $N(\bar\rho_{E,5})$ for \curve{980i1} and
\curve{140b1}. Which of these curves occur in the paper's Example of Section 3, directly or as the twist of a curve there, and which are illustrations
added here? Which congruences of that Example do the former explain, and why does the one at $\ell=5$ not fall under Theorem thm:qual?

\emph{Answer for the conductors.} $7^2\cdot11$, $11$, $11^2$, $2^8$ at $\ell=3$, and $2^2\cdot7^2$, $2^2$ at $\ell=5$; compare \Cref{ex:galrep:serre-conductor-examples}. At
$\ell=5$ both curves have $5\mid N$.
\end{exercise}

\begin{exercise}\label{ex:galrep:ns-ordinary}
Let $E$ have image \texttt{3Ns} and $3\nmid N$. (a) Show that $\omega_2$ has order $8$ for $\ell=3$, so that in the supersingular case $\bar\rho_{E,3}(I_3)$ would contain an
element of order $8$. (b) Show that $N_s(\F_3)$ has no element of order $8$, and deduce that $E$ is ordinary at $3$: a second proof of
\Cref{prop:galrep:ns-structure}(vi). (c) Show that in the ordinary case $\bar\rho(I_3)$ has order $2$ or $6$, and use $N_s(\F_3)$ to show it has order $2$.

\emph{Hint.} (a) $\omega_2$ maps onto $\F_9^\times$. (b) $N_s(\F_3)$ is dihedral of order $8$; its elements have orders $1$, $2$, $4$. (c) In the ordinary case
$\bar\rho(I_3)$ lies in $\sm{\omega}{*}{0}{1}$ with $\omega$ of order $2$, a group of order $6$, and $\bar\rho(I_3)$ surjects onto the values of $\omega$; $N_s(\F_3)$ has no element
of order $3$.
\end{exercise}

\section{Notes and references}\label{sec:galrep:notes}

\paragraph{Sources followed.} The Tate module, the Weil pairing and the criterion of N\'eron, Ogg and Shafarevich follow Silverman
\cite[Chapters III and VII]{SilvermanAEC} (III.7, III.8 and VII.7 are the relevant sections in the editions we know; these are unverified pointers) and
Serre and Tate \cite[Sections 1 and 2]{SerreTate68}; the proof of (d)$\Rightarrow$(a) in \Cref{thm:galrep:nos} is the one in \cite[Section 1]{SerreTate68}. The Tate
curve and its Galois module are in \cite[Chapter V]{SilvermanATAEC} (V.3 and V.5, unverified pointers), and the matrix form at a multiplicative prime
is the elliptic-curve case of \cite[Theorem 3.1(e)]{DDT}, the number the paper cites. The image of inertia at a potentially good prime is due to
Serre and Tate \cite[Section 2]{SerreTate68} and Kraus \cite{Kraus90}; we cite Kraus's paper as a whole (\Cref{ch:reduction} explains why). The fundamental
characters and the description of $\bar\rho|_{I_\ell}$ in the ordinary and supersingular cases are in Serre \cite[Section 1]{Serre72}; the description for
newforms, due to Deligne and Fontaine, is quoted from Edixhoven \cite[Section 2]{Edixhoven92}, as in the paper. Serre's conjecture, the Serre conductor
and the weight are from Serre \cite{Serre87} (Sections 1 and 2 for the definitions and Section 4 for examples from elliptic curves: unverified
pointers); the divisibility $N(\bar\rho)\mid N$ is Carayol's \cite{Carayol89} for newforms, and the conductor of a residual representation of a modular
form is studied by Livn\'e \cite{Livne89}. The open image theorem is Serre's \cite{Serre72}; the list of degrees of rational isogenies is Mazur's
\cite{Mazur78}; the images for non-CM curves at $\ell\le13$ are Zywina's \cite{Zywina15}, and the labels Sutherland's \cite{Sutherland16}. Serre's conjecture
is the theorem of Khare and Wintenberger \cite{KhareWintenberger}; the level-lowering theorem is due to Ribet \cite{Ribet90}, Carayol \cite{Carayol89} and Diamond
\cite{Diamond95}, in the form of \cite[Theorem 3.15]{DDT} quoted by the paper. Faltings's theorem is \cite{Faltings83}; crystalline representations are
introduced in Berger's survey \cite{Berger}.

\paragraph{Theorem numbers.} We cite theorem numbers only where the paper cites them: \cite[Theorem 3.1(e)]{DDT}, \cite[Theorem 3.15]{DDT}, \cite[Section
2]{Edixhoven92}, and the definition \cite[(31)]{DFG04} of $S_f$. All other references are to sections, and the section numbers of books given above are
pointers we have not checked against printed editions.

\paragraph{What is new here.} Three statements of this chapter, also proved in the paper, are not standard in the literature as far as we
know: that the hypothesis on $\Q(\sqrt{\ell^*})$ is automatic for $\ell\ge5$, $\ell\nmid N$ (\Cref{thm:galrep:lstar-automatic}, the paper's Proposition prop:Sf);
that the types $\kod{IV}$, $\kod{IV}^*$ at $2$ have cyclic inertia of order $3$ (\Cref{lem:galrep:IV}, \Cref{cor:galrep:IV-at-2}); and the structure of the curves
with image \texttt{3Ns} (\Cref{prop:galrep:ns-structure}). The first may well be known; we found no reference. The formula $a_p(\bar\rho_{E,\ell})=f_p-\dim\Phig_p[\ell]$
(\Cref{prop:galrep:serre-conductor-E}) is standard in substance (it is how the Serre conductor of an elliptic curve is computed in practice) but we
know no reference that states it in this form. The pairing of the images \texttt{3B.1.1} and \texttt{3B.1.2} by the ends of the $3$-isogeny paths
(\Cref{prop:galrep:3-isogeny-paths}) is elementary and surely known; we found no reference that states it.

\paragraph{Computations.} All numbers were computed with PARI/GP 2.17.2 \cite{PARI} through \texttt{cypari2}, on Cremona's files \texttt{allcurves},
\texttt{allisog}, \texttt{galrep} and \texttt{degphi} \cite{ecdata} for conductors below $10{,}000$ ($64{,}687$ curves, $462$ with complex multiplication, recognised
by their $j$-invariant among the $13$ rational CM $j$-invariants). The checks reported in
the text are: the mod $3$ image of every curve from $\psi_3$ against the \texttt{galrep} label (\Cref{thm:galrep:mod3-images}); the $3$-isogeny paths of all
$38{,}042$ classes (\Cref{prop:galrep:3-isogeny-paths}); the image of inertia modulo $3$
by Kodaira type at $q\ge5$ and at $q=2$ (\Cref{prop:galrep:mod3-inertia}, \Cref{rem:galrep:wild-inertia}); all $4{,}750$ curves with type $\kod{IV}$ or $\kod{IV}^*$ at $2$
(\Cref{cor:galrep:IV-at-2}); all $146$ curves with image \texttt{3Ns}, and the $103$ optimal ones among them, the optimal curve being the one listed in \texttt{degphi}
(\Cref{prop:galrep:ns-structure}); the ordinarity at $\ell$ of the curves with reducible
$\bar\rho_{E,\ell}$, $\ell\nmid N$ (\Cref{cor:galrep:ordinary}); the congruences $a_p\equiv1+p$ (\Cref{prop:galrep:reducible}); the dihedral images at $\ell=5,7$
(\Cref{thm:galrep:lstar-automatic}); and the Serre conductors and congruences of \Cref{ex:galrep:serre-conductor-examples}. Every count is for conductor below
$10{,}000$; nothing beyond that range was checked.
